% Lutte contre la mauvaise alimentation et l'inactivité physique — SVT 1ère D — Cours signé OnBuch+
% Programme officiel MINESEC. Compiler avec : tectonic ND10-mauvaise-alimentation-inactivite-d.tex
\newcommand{\DOCMATIERE}{SVT}
\newcommand{\DOCNIVEAU}{1ère D}
\newcommand{\DOCMODULE}{Module 2 : L'Éducation à la Santé}
\newcommand{\DOCLECON}{ND10}
\newcommand{\DOCTITRE}{Lutte contre la mauvaise alimentation et l'inactivité physique}
% OnBuch+ — préambule « cours signés » ENRICHI (compilé avec tectonic / XeLaTeX)
% Système visuel « Pop » d'OnBuch+ : fond crème (jamais blanc), bordures encre
% épaisses, ombres « dures » décalées, titres Archivo Black en capitales.
% Version MATHÉMATIQUES : graphes (pgfplots/tikz), boîtes pédagogiques
% (définition, propriété, méthode, attention, le savais-tu, exercice résolu,
% exercice, TP, bilan), page de garde et signature OnBuch+ ancrée en bas.
%
% Macros attendues dans main.tex AVANT \input{preamble} :
%   \DOCMATIERE  \DOCNIVEAU  \DOCMODULE  \DOCLECON (numéro)  \DOCTITRE
\documentclass[11pt]{article}

\usepackage{fontspec}
\usepackage[french]{babel}
\usepackage[a4paper,margin=2cm,top=3.4cm,bottom=2.8cm,headheight=1.4cm,footskip=1.3cm]{geometry}
\usepackage{amsmath,amssymb,amsfonts}
\usepackage{mathtools} % \xmapsto, \coloneqq... souvent utilisés par les modèles
\usepackage{cancel} % simplifications barrées (bilans de forces)
% \abs{z} (module, valeur absolue) et \norm{u} (norme), utilisés par les modèles
\providecommand{\abs}{}\let\abs\relax\DeclarePairedDelimiter{\abs}{\lvert}{\rvert}
\providecommand{\norm}{}\let\norm\relax\DeclarePairedDelimiter{\norm}{\lVert}{\rVert}
\usepackage[table]{xcolor} % [table] : fournit aussi \rowcolors (alternance de lignes)
\usepackage{pagecolor}
\usepackage{tikz}
\usetikzlibrary{arrows.meta,positioning,calc,shapes.geometric,shapes.misc,shapes.arrows,decorations.pathmorphing,decorations.pathreplacing,decorations.markings,patterns,shadows,mindmap,trees,fit,backgrounds,matrix,intersections,angles,quotes}
\usepackage{pgfplots}
\usepackage[version=4]{mhchem} % équations nucléaires \ce{^{235}_{92}U}
\usepackage[european,siunitx]{circuitikz} % schémas électriques
\pgfplotsset{compat=1.18}
\usepgfplotslibrary{fillbetween,groupplots} % aires entre courbes (calcul intégral)
\usepackage{enumitem}
\usepackage{fancyhdr}
% [most] charge toutes les bibliothèques tcolorbox (listings, minted, raster,
% documentation...) et gonfle inutilement la mémoire TeX sur un document long
% (~300 boîtes) : on ne charge que ce qui est réellement utilisé.
\usepackage[skins,breakable]{tcolorbox}
\usepackage{graphicx}
\usepackage{colortbl}
\usepackage{booktabs}
\usepackage{tabularx}
\usepackage{diagbox} % case d'en-tête diagonale des tableaux à double entrée
% Colonnes extensibles centrée / gauche / droite (C, L, R), souvent utilisées par les modèles
\newcolumntype{C}{>{\centering\arraybackslash}X}
\newcolumntype{L}{>{\raggedright\arraybackslash}X}
\newcolumntype{R}{>{\raggedleft\arraybackslash}X}
\usepackage{array}
\usepackage{multicol}
\usepackage{siunitx}
% parse-numbers=false : accepte aussi \SI{2,5 \times 10^{-3}}{...}
\sisetup{locale=FR,output-decimal-marker={,},group-minimum-digits=4,parse-numbers=false}
\DeclareSIUnit\ohm{\ensuremath{\symup{\Omega}}} % U+2126 absent de la police
\DeclareSIUnit\division{div} % réglages d'oscilloscope (ms/div, V/div)
\DeclareSIUnit\tr{tr} % tours (vitesse de rotation, ex. \SI{1500}{\tr\per\minute})
\DeclareSIUnit\ch{ch} % cheval-vapeur (puissance moteur, unité usuelle hors SI)
\DeclareSIUnit\franccfa{FCFA} % franc CFA (coûts, contexte camerounais)
\DeclareSIUnit\dioptre{\ensuremath{\delta}} % vergence d'une lentille (dioptrie)
\usepackage{qrcode}
% \degree : commande fréquemment utilisée par les modèles pour un angle ou une
% température en dehors de \SI{}{} (non fournie par les paquets chargés).
\providecommand{\degree}{^{\circ}}
% \qed (fin de démonstration), utilisé par les modèles sans amsthm
\providecommand{\qed}{\hfill$\square$}
% \barre : double barre des valeurs interdites dans un tableau de variations (tkz-tab)
\providecommand{\barre}{\ensuremath{\|}}
% environnement proof (amsthm non chargé) utilisé par les modèles
\ifdefined\proof\else\newenvironment{proof}[1][Démonstration]{\par\noindent\textit{#1.}\ }{\qed\par}\fi
\usepackage{lastpage}
\usepackage{hyperref}
\hypersetup{hidelinks}

% ============================================================
% Palette « Pop » OnBuch+ — mêmes hex que la classe P (plus_ui.dart)
% ============================================================
\definecolor{poporange}{HTML}{F59321}
\definecolor{popdark}{HTML}{E07A0C}
\definecolor{popcream}{HTML}{FFF6E8}
\definecolor{popink}{HTML}{1C1714}
\definecolor{poppurple}{HTML}{7A5AE0}
\definecolor{popgreen}{HTML}{2E9E6A}
\definecolor{poppink}{HTML}{E0457B}
\definecolor{popblue}{HTML}{2F7DE1}
\definecolor{popgold}{HTML}{C98A12}
\definecolor{popmuted}{HTML}{8A7F76}
\definecolor{popline}{HTML}{EADFD2}
\definecolor{popgray}{HTML}{8A7F76}
\colorlet{popgrey}{popgray}
% Alias tolérants (le modèle nomme parfois les couleurs différemment)
\colorlet{popwhite}{white}
\colorlet{popblack}{popink}
\colorlet{popred}{poppink}
\colorlet{popyellow}{popgold}
% Teintes claires pour fonds de tableaux / zones
\colorlet{popblueL}{popblue!10!white}
\colorlet{poporangeL}{poporange!14!white}
\colorlet{popgreenL}{popgreen!12!white}
\colorlet{poppurpleL}{poppurple!12!white}
\colorlet{poppinkL}{poppink!12!white}
\colorlet{popgoldL}{popgold!14!white}

\pagecolor{popcream}

% ============================================================
% Typographie
% ============================================================
\defaultfontfeatures{Ligatures=TeX}
\setmainfont{PlusJakartaSans-Medium.ttf}[Path=../../../fonts/,BoldFont=PlusJakartaSans-Bold.ttf]
% Maths en sans-serif (Fira Math) avec chiffres et lettres droites de Plus
% Jakarta, pour que les formules se fondent dans le texte.
\usepackage[math-style=ISO,bold-style=ISO]{unicode-math}
\setmathfont{FiraMath-Regular.otf}[Path=../../../fonts/]
\setmathfont{PlusJakartaSans-Medium.ttf}[Path=../../../fonts/,range={up/num,up/{Latin,latin},\mathup}]
\usepackage{icomma} % virgule décimale sans espace en mode math
% Caractères mathématiques Unicode absents des polices de texte (Plus Jakarta,
% Archivo Black) : rendus par leur équivalent mathématique, en texte comme en
% formule — sinon ils s'affichent en carré vide (ex. « Arithmétique dans ℤ »).
\usepackage{newunicodechar}
\newunicodechar{ℝ}{\ensuremath{\mathbb{R}}}
\newunicodechar{ℤ}{\ensuremath{\mathbb{Z}}}
\newunicodechar{ℂ}{\ensuremath{\mathbb{C}}}
\newunicodechar{ℕ}{\ensuremath{\mathbb{N}}}
\newunicodechar{ℚ}{\ensuremath{\mathbb{Q}}}
\newunicodechar{𝔻}{\ensuremath{\mathbb{D}}}
\newunicodechar{α}{\ensuremath{\alpha}}
\newunicodechar{β}{\ensuremath{\beta}}
\newunicodechar{γ}{\ensuremath{\gamma}}
\newunicodechar{δ}{\ensuremath{\delta}}
\newunicodechar{ε}{\ensuremath{\varepsilon}}
\newunicodechar{θ}{\ensuremath{\theta}}
\newunicodechar{λ}{\ensuremath{\lambda}}
\newunicodechar{μ}{\ensuremath{\mu}}
\newunicodechar{σ}{\ensuremath{\sigma}}
\newunicodechar{τ}{\ensuremath{\tau}}
\newunicodechar{φ}{\ensuremath{\varphi}}
\newunicodechar{ψ}{\ensuremath{\psi}}
\newunicodechar{ω}{\ensuremath{\omega}}
\newunicodechar{Σ}{\ensuremath{\Sigma}}
% majuscules produites par \MakeUppercase dans les titres (α-aminés -> Α)
\newunicodechar{Α}{\ensuremath{\alpha}}
\newunicodechar{Β}{\ensuremath{\beta}}
\newunicodechar{Γ}{\ensuremath{\Gamma}}
\newunicodechar{Δ}{\ensuremath{\Delta}}
\newunicodechar{Ε}{\ensuremath{\varepsilon}}
\newunicodechar{Θ}{\ensuremath{\Theta}}
\newunicodechar{Λ}{\ensuremath{\Lambda}}
\newunicodechar{Μ}{\ensuremath{\mu}}
\newunicodechar{Τ}{\ensuremath{\tau}}
\newunicodechar{Φ}{\ensuremath{\Phi}}
\newunicodechar{Ψ}{\ensuremath{\Psi}}
\newunicodechar{Ω}{\ensuremath{\Omega}}
\newunicodechar{↦}{\ensuremath{\mapsto}}
\newunicodechar{∈}{\ensuremath{\in}}
\newunicodechar{∉}{\ensuremath{\notin}}
\newunicodechar{∧}{\ensuremath{\wedge}}
\newunicodechar{⇔}{\ensuremath{\Leftrightarrow}}
\newunicodechar{⟺}{\ensuremath{\Longleftrightarrow}}
\newunicodechar{⇒}{\ensuremath{\Rightarrow}}
\newunicodechar{⟹}{\ensuremath{\Longrightarrow}}
\newunicodechar{∘}{\ensuremath{\circ}}
\newunicodechar{∪}{\ensuremath{\cup}}
\newunicodechar{∩}{\ensuremath{\cap}}
\newunicodechar{≡}{\ensuremath{\equiv}}
\newunicodechar{⊂}{\ensuremath{\subset}}
\newunicodechar{⊥}{\ensuremath{\perp}}
\newunicodechar{∃}{\ensuremath{\exists}}
\newunicodechar{∀}{\ensuremath{\forall}}
\newunicodechar{∛}{\ensuremath{\sqrt[3]{\,}}}
\newunicodechar{∜}{\ensuremath{\sqrt[4]{\,}}}
\newunicodechar{⃗}{\ensuremath{{}^{\to}}}
\newunicodechar{ⁿ}{\ifmmode^{n}\else\textsuperscript{\ensuremath{n}}\fi}
\newunicodechar{ˣ}{\ifmmode^{x}\else\textsuperscript{\ensuremath{x}}\fi}
\newunicodechar{ᵇ}{\ifmmode^{b}\else\textsuperscript{\ensuremath{b}}\fi}
\newunicodechar{ʳ}{\ifmmode^{r}\else\textsuperscript{\ensuremath{r}}\fi}
\newunicodechar{ᵘ}{\ifmmode^{u}\else\textsuperscript{\ensuremath{u}}\fi}
\newunicodechar{ʸ}{\ifmmode^{y}\else\textsuperscript{\ensuremath{y}}\fi}
\newunicodechar{ᵃ}{\ifmmode^{a}\else\textsuperscript{\ensuremath{a}}\fi}
\newunicodechar{ᵉ}{\ifmmode^{e}\else\textsuperscript{\ensuremath{e}}\fi}
\newunicodechar{ᵏ}{\ifmmode^{k}\else\textsuperscript{\ensuremath{k}}\fi}
\newunicodechar{ᵖ}{\ifmmode^{p}\else\textsuperscript{\ensuremath{p}}\fi}
\newunicodechar{ᴺ}{\ifmmode^{N}\else\textsuperscript{\ensuremath{N}}\fi}
\newunicodechar{ᶜ}{\ifmmode^{c}\else\textsuperscript{\ensuremath{c}}\fi}
\newunicodechar{ʰ}{\ifmmode^{h}\else\textsuperscript{\ensuremath{h}}\fi}
\newunicodechar{ᵠ}{\ifmmode^{q}\else\textsuperscript{\ensuremath{q}}\fi}
\newunicodechar{ₙ}{\ifmmode_{n}\else\textsubscript{\ensuremath{n}}\fi}
\newunicodechar{ₐ}{\ifmmode_{a}\else\textsubscript{\ensuremath{a}}\fi}
\newunicodechar{ᵢ}{\ifmmode_{i}\else\textsubscript{\ensuremath{i}}\fi}
\newunicodechar{ᵣ}{\ifmmode_{r}\else\textsubscript{\ensuremath{r}}\fi}
\newunicodechar{ₖ}{\ifmmode_{k}\else\textsubscript{\ensuremath{k}}\fi}
\newunicodechar{ᵦ}{\ifmmode_{\beta}\else\textsubscript{\ensuremath{\beta}}\fi}
\newunicodechar{ₓ}{\ifmmode_{x}\else\textsubscript{\ensuremath{x}}\fi}
\newfontfamily\popdisplay{ArchivoBlack-Regular.ttf}[Path=../../../fonts/]
\newfontfamily\poplogo{SpaceGrotesk-Bold.ttf}[Path=../../../fonts/]
\newfontfamily\popsemi{PlusJakartaSans-SemiBold.ttf}[Path=../../../fonts/]
\newfontfamily\popxbold{PlusJakartaSans-ExtraBold.ttf}[Path=../../../fonts/]
\newfontfamily\popxboldls{PlusJakartaSans-ExtraBold.ttf}[Path=../../../fonts/,LetterSpace=3.0]

\setlist[enumerate]{leftmargin=1.7em,itemsep=5pt,topsep=4pt}
\setlist[itemize]{leftmargin=1.7em,itemsep=4pt,topsep=4pt,label={\color{poporange}\textbullet}}
\setlength{\parindent}{0pt}
\setlength{\parskip}{5pt}
\renewcommand{\arraystretch}{1.25}

% pgfplots : style Pop par défaut
\pgfplotsset{
  every axis/.append style={axis line style={popink,line width=1pt},tick style={popink},
    grid style={popline,line width=0.5pt},label style={font=\small},tick label style={font=\footnotesize}},
  popcurve/.style={poporange,line width=1.8pt,no markers,smooth},
}
% Même style disponible dans un \draw TikZ simple (hors pgfplots)
\tikzset{popcurve/.style={poporange,line width=1.8pt}}
\tikzset{popgrid/.style={popline,very thin}}
\colorlet{popcurve}{poporange} % le nom du style est parfois utilisé comme couleur

% ============================================================
% Pill / squiggle
% ============================================================
\newcommand{\poppill}[3]{%
  \tikz[baseline=-0.55ex]\node[draw=popink,line width=1.2pt,rounded corners=2.6pt,%
    fill=#2,inner xsep=6.5pt,inner ysep=3pt]{%
    \popxboldls\fontsize{7}{7}\selectfont\color{#3}\MakeUppercase{#1}};%
}
\newcommand{\popsquiggle}[1]{%
  \tikz[baseline=-0.4ex]{
    \draw[#1,line width=2.6pt,line cap=round]
      plot [smooth] coordinates {(0,0.09) (0.34,0.01) (0.68,0.21) (1.02,0.01) (1.36,0.10)};
  }%
}
% Mot-clé surligné (marqueur orange sous le texte)
\newcommand{\cle}[1]{{\popxbold\color{popdark}#1}}

% ============================================================
% En-tête / pied
% ============================================================
\pagestyle{fancy}
\fancyhf{}
\renewcommand{\headrulewidth}{0pt}
\renewcommand{\footrulewidth}{0pt}
\lhead{\raisebox{-0.15ex}{\poplogo\fontsize{13}{13}\selectfont\color{popink}OnBuch\color{poporange}+}}
\rhead{\poppill{\DOCMATIERE\ \textbullet\ \DOCNIVEAU\ \textbullet\ Leçon \DOCLECON}{white}{popink}}
\lfoot{\popxbold\fontsize{7.6}{7.6}\selectfont\color{popmuted}\MakeUppercase{Cours signé OnBuch+}\ \textbullet\ {\popsemi\DOCTITRE}}
\rfoot{\tikz[baseline=-0.6ex]\node[rounded corners=8.5pt,fill=popink,minimum height=17pt,minimum width=17pt,inner xsep=5pt,inner sep=0]{\popxbold\fontsize{8}{8}\selectfont\color{popcream}\thepage\,/\,\pageref*{LastPage}};}

% ============================================================
% Titres
% ============================================================
\newcommand{\chaptitre}[2]{%
  \begin{tcolorbox}[enhanced,colback=poporange,colframe=popink,boxrule=2.6pt,arc=11pt,
      drop shadow=popink,left=15pt,right=15pt,top=12pt,bottom=11pt,before skip=3pt,after skip=18pt]
    \poppill{#2}{popink}{popcream}\par\vspace{9pt}
    {\raggedright\hyphenpenalty=10000\exhyphenpenalty=10000\popdisplay\fontsize{22}{24}\selectfont\color{popcream}\MakeUppercase{#1}\par}
  \end{tcolorbox}
}
% Section numérotée : pastille orange avec le numéro + titre Archivo
\newcounter{popsec}
\newcommand{\coursec}[1]{%
  \stepcounter{popsec}\setcounter{popsub}{0}%
  \par\vspace{16pt}\noindent
  \begin{minipage}{\linewidth}
  \tikz[baseline=-0.7ex]\node[draw=popink,line width=1.6pt,fill=poporange,rounded corners=4pt,minimum size=22pt,inner sep=0]{\popdisplay\fontsize{12}{12}\selectfont\color{popcream}\thepopsec};%
  \hspace{8pt}{\popdisplay\fontsize{15}{17}\selectfont\color{popink}\MakeUppercase{#1}}\par
  \vspace{4pt}\noindent{\color{poporange}\rule{36pt}{3.6pt}}
  \end{minipage}\par\nopagebreak\vspace{8pt}
}
\newcounter{popsub}
\newcommand{\courssub}[1]{%
  \stepcounter{popsub}%
  \par\vspace{9pt}\noindent{\popxbold\fontsize{11.5}{13}\selectfont\color{popink}{\color{poporange}\thepopsec.\thepopsub}\hspace{6pt}#1}\par\nopagebreak\vspace{4pt}
}

% ============================================================
% Boîtes pédagogiques — même recette : fond blanc, bordure épaisse colorée,
% ombre dure encre, badge plein. Titre optionnel [#1].
% ============================================================
\tcbset{popbase/.style={breakable,enhanced,colback=white,boxrule=2.3pt,arc=9pt,
  shadow={3pt}{-3pt}{0pt}{popink},left=14pt,right=14pt,top=11pt,bottom=11pt,before skip=11pt,after skip=15pt,
  fontupper=\popsemi\fontsize{10.6}{14.5}\selectfont\color{popink},
  fontlower=\popsemi\fontsize{10.6}{14.5}\selectfont\color{popink}}}
% Coche dessinée (le glyphe ✓ n'existe pas dans les polices du document)
\newcommand{\popcheck}[1]{\tikz[baseline=-0.5ex]\draw[#1,line width=1.6pt,line cap=round,line join=round](-0.13,0.0)--(-0.04,-0.09)--(0.13,0.1);}
\providecommand{\coche}{\popcheck{popgreen}}
\providecommand{\resultat}[1]{\par\smallskip\noindent\fcolorbox{popgreen}{popgreenL}{\parbox{\dimexpr\linewidth-2\fboxsep-2\fboxrule\relax}{\textbf{Résultat.} #1}}\par\smallskip}
\newcommand{\popboxhead}[3]{\poppill{#1}{#2}{white}\ifx\relax#3\relax\else\hspace{6pt}{\popxbold\fontsize{10.6}{12}\selectfont\color{popink}#3}\fi\par\vspace{7pt}}

\newtcolorbox{aretenir}[1][]{popbase,colframe=popgreen,before upper={\popboxhead{À retenir}{popgreen}{#1}}}
\newtcolorbox{exemplebox}[1][]{popbase,colframe=poporange,before upper={\popboxhead{Exemple}{poporange}{#1}}}
\newtcolorbox{definition}[1][]{popbase,colframe=popblue,before upper={\popboxhead{Définition}{popblue}{#1}}}
\newtcolorbox{propriete}[1][]{popbase,colframe=poppurple,before upper={\popboxhead{Propriété}{poppurple}{#1}}}
\newtcolorbox{methode}[1][]{popbase,colframe=popgold,before upper={\popboxhead{Méthode}{popgold}{#1}}}
\newtcolorbox{attention}[1][]{popbase,colframe=poppink,before upper={\popboxhead{Attention}{poppink}{#1}}}
\newtcolorbox{experience}[1][]{popbase,colframe=popblue,colback=popblueL,before upper={\popboxhead{Expérience}{popblue}{#1}}}
% « Le savais-tu ? » : carte violette pleine (contexte, culture, Cameroun)
\newtcolorbox{savaistu}[1][]{popbase,colback=poppurple,colframe=popink,
  fontupper=\popsemi\fontsize{10.4}{14.2}\selectfont\color{white},
  before upper={\poppill{Le savais-tu\,?}{white}{poppurple}\ifx\relax#1\relax\else\hspace{6pt}{\popxbold\color{white}#1}\fi\par\vspace{7pt}}}
% Exercice résolu : énoncé en haut, \tcblower puis la solution sur fond crème
\newcounter{popexo}\newcounter{popexr}
\newtcolorbox{exoresolu}[1][]{popbase,colframe=poporange,colbacklower=popcream,
  segmentation style={popink,line width=1.2pt,dashed},
  before upper={\stepcounter{popexr}\popboxhead{Exercice résolu \thepopexr}{poporange}{#1}},
  before lower={\poppill{Solution}{popink}{popcream}\par\vspace{6pt}}}
% Exercice (non résolu en place) — la correction est dans la partie Corrigés
\newtcolorbox{exercice}[1][]{popbase,colframe=popink,
  before upper={\stepcounter{popexo}\popboxhead{Exercice \thepopexo}{popink}{#1}}}
% Corrigé
\newtcolorbox{corrige}[1][]{popbase,colframe=popgreen,colback=popgreenL,
  before upper={\popboxhead{Corrigé}{popgreen}{#1}}}
% Objectifs / prérequis / bilan
\newtcolorbox{objectifs}{popbase,colframe=popink,colback=poporangeL,
  before upper={\popboxhead{Objectifs de la leçon}{popink}{}}}
\newtcolorbox{prerequis}{popbase,colframe=popink,colback=popblueL,
  before upper={\popboxhead{Prérequis}{popblue}{}}}
\newtcolorbox{bilan}{popbase,colframe=popink,colback=white,boxrule=2.8pt,
  before upper={\popboxhead{Fiche bilan}{poporange}{L'essentiel à connaître pour l'examen}}}
% Figure encadrée avec légende
\newtcolorbox{popfigure}[1][]{enhanced,breakable=false,colback=white,colframe=popline,boxrule=1.4pt,arc=8pt,
  left=10pt,right=10pt,top=10pt,bottom=8pt,before skip=10pt,after skip=12pt,halign=center,
  lower separated=false,halign lower=center,
  fontlower=\popsemi\fontsize{9}{11.5}\selectfont\color{popmuted},
  #1}
\newcommand{\legende}[1]{\par\vspace{6pt}{\popsemi\fontsize{9}{11.5}\selectfont\color{popmuted}#1}}

% ============================================================
% Signature OnBuch+
% ============================================================
\newcommand{\poplogoinline}[1]{{\poplogo\fontsize{#1}{#1}\selectfont\color{popink}OnBuch\color{poporange}+}}
\newcommand{\signatureonbuch}{%
  \begin{tcolorbox}[enhanced,colback=popink,colframe=popink,boxrule=2.2pt,arc=10pt,
      drop shadow=poporange,left=14pt,right=14pt,top=10pt,bottom=10pt,before skip=0pt,after skip=0pt]
    \tikz[baseline=-0.7ex]\node[circle,fill=popgreen,draw=popcream,line width=1.3pt,minimum size=22pt,inner sep=0]{\popcheck{white}};%
    \hspace{9pt}\begin{minipage}[c]{0.62\linewidth}
      {\popxbold\fontsize{12}{13}\selectfont\color{popcream}Signé\ \textbullet\ {\poplogo OnBuch\color{poporange}+}}\par\vspace{2pt}
      {\raggedright\popsemi\fontsize{8}{10.5}\selectfont\color{popline}\MakeUppercase{Équipe pédagogique OnBuch+}\\\MakeUppercase{Conforme au programme officiel MINESEC}\par}
    \end{minipage}\hfill
    {\poplogo\fontsize{22}{22}\selectfont\color{popcream}OnBuch\color{poporange}+}
  \end{tcolorbox}%
}

% ============================================================
% Page de garde
% #1 titre · #2 matière · #3 niveau · #4 module · #5 n° leçon · #6 durée ·
% #7 description / objectifs (texte ou itemize)
% ============================================================
\newcommand{\pagegarde}[7]{%
  \thispagestyle{empty}
  \newgeometry{a4paper,margin=1.7cm,top=1.6cm,bottom=1.4cm}%
  \begin{tikzpicture}[remember picture,overlay]
    \fill[poporange!18] ([xshift=-3.2cm,yshift=-3.6cm]current page.north east) circle (4.6cm);
    \fill[poppurple!14] ([xshift=2.4cm,yshift=7.6cm]current page.south west) circle (3.8cm);
  \end{tikzpicture}%
  {\poplogo\fontsize{30}{30}\selectfont\color{popink}OnBuch\color{poporange}+}\hfill
  \raisebox{0.6ex}{\poppill{Cours signé \textbullet\ Premium}{popink}{popcream}}\par
  \vspace{1pt}\popsquiggle{poppurple}\par
  \vspace{8pt}
  \begin{tcolorbox}[enhanced,colback=poporange,colframe=popink,boxrule=2.8pt,arc=13pt,
      drop shadow=popink,left=18pt,right=18pt,top=12pt,bottom=13pt,after skip=10pt]
    \poppill{#2 \textbullet\ #3}{popink}{popcream}\hspace{5pt}\poppill{#4}{white}{popink}\par\vspace{8pt}
    {\popdisplay\fontsize{14}{14}\selectfont\color{popink}LEÇON #5}\par\vspace{4pt}
    {\raggedright\hyphenpenalty=10000\exhyphenpenalty=10000\popdisplay\fontsize{25}{27}\selectfont\color{popcream}\MakeUppercase{#1}\par}
    \vspace{8pt}
    {\popxbold\fontsize{9.5}{9.5}\selectfont\color{popink}\MakeUppercase{Durée conseillée : #6}}
  \end{tcolorbox}
  \begin{tcolorbox}[enhanced,colback=white,colframe=popink,boxrule=2.2pt,arc=10pt,
      drop shadow=popink,left=16pt,right=16pt,top=10pt,bottom=10pt,after skip=8pt]
    \poppill{Dans cette leçon}{poppurple}{white}\par\vspace{5pt}
    \popsemi\fontsize{9.3}{12.6}\selectfont\color{popink}#7
  \end{tcolorbox}
  \noindent\begin{minipage}[t]{0.487\linewidth}
    \begin{tcolorbox}[enhanced,colback=popgreen,colframe=popink,boxrule=2.2pt,arc=10pt,
        drop shadow=popink,left=11pt,right=11pt,top=10pt,bottom=10pt,height=3.1cm,valign=center]
      \tcbox[colback=white,colframe=popink,boxrule=1.3pt,arc=5pt,boxsep=0pt,nobeforeafter,
        left=3pt,right=3pt,top=3pt,bottom=3pt,valign=center]{\qrcode[height=1.9cm]{https://play.google.com/store/apps/details?id=com.luvvix.onbuch&hl=fr}}%
      \hspace{8pt}\begin{minipage}[c]{0.5\linewidth}
        {\popdisplay\fontsize{11}{12.5}\selectfont\color{white}\MakeUppercase{Télécharge OnBuch}}\par\vspace{4pt}
        {\raggedright\popsemi\fontsize{8.4}{11}\selectfont\color{white}Gratuit \textbullet\ Google Play\\Tuteur IA, annales, concours, résultats\par}
      \end{minipage}
    \end{tcolorbox}
  \end{minipage}\hfill
  \begin{minipage}[t]{0.487\linewidth}
    \begin{tcolorbox}[enhanced,colback=poppurple,colframe=popink,boxrule=2.2pt,arc=10pt,
        drop shadow=popink,left=11pt,right=11pt,top=10pt,bottom=10pt,height=3.1cm,valign=center]
      \tikz[baseline=-0.7ex]\node[circle,fill=white,draw=popink,line width=1.3pt,minimum size=1.6cm,inner sep=0]{\popdisplay\fontsize{24}{24}\selectfont\color{poppurple}+};%
      \hspace{8pt}\begin{minipage}[c]{0.6\linewidth}
        {\popdisplay\fontsize{11}{12.5}\selectfont\color{white}\MakeUppercase{Passe à OnBuch+}}\par\vspace{4pt}
        {\raggedright\popsemi\fontsize{8.4}{11}\selectfont\color{white}Tous les cours signés, Tuteur IA illimité, fiches PDF — onglet OnBuch+ de l'app\par}
      \end{minipage}
    \end{tcolorbox}
  \end{minipage}
  \vspace{8pt}
  \popcontenu
  \vfill
  \signatureonbuch
  \restoregeometry
  \newpage
}

% Rangée « Ce cours contient » (page de garde)
\newcommand{\poptile}[3]{% #1 couleur #2 gros texte #3 légende
  \begin{tcolorbox}[enhanced,colback=#1,colframe=popink,boxrule=2pt,arc=9pt,shadow={2.5pt}{-2.5pt}{0pt}{popink},
      left=6pt,right=6pt,top=9pt,bottom=9pt,halign=center,valign=center,height=1.75cm,width=\linewidth,nobeforeafter]
    {\popdisplay\fontsize{12.5}{13}\selectfont\color{white}\MakeUppercase{#2}}\par\vspace{3pt}
    {\centering\popsemi\fontsize{7.8}{9.5}\selectfont\color{white}#3\par}
  \end{tcolorbox}}
\newcommand{\popcontenu}{%
  \par\noindent{\popxboldls\fontsize{7.5}{7.5}\selectfont\color{popmuted}CE COURS SIGNÉ CONTIENT}\par\vspace{7pt}
  \noindent\begin{minipage}[t]{0.235\linewidth}\poptile{popblue}{Cours}{complet, illustré, pas à pas}\end{minipage}\hfill
  \begin{minipage}[t]{0.235\linewidth}\poptile{poporange}{Méthodes}{et exercices résolus}\end{minipage}\hfill
  \begin{minipage}[t]{0.235\linewidth}\poptile{poppink}{Exercices}{type examen + corrigés}\end{minipage}\hfill
  \begin{minipage}[t]{0.235\linewidth}\poptile{popgreen}{Fiche bilan}{l'essentiel à retenir}\end{minipage}\par}

% Sommaire
\newtcolorbox{sommaire}{enhanced,colback=white,colframe=popink,boxrule=2.2pt,arc=10pt,
  drop shadow=popink,left=18pt,right=18pt,top=13pt,bottom=13pt,before skip=6pt,after skip=20pt,
  before upper={\poppill{Sommaire}{poppurple}{white}\par\vspace{9pt}\popsemi\fontsize{10.6}{14.5}\selectfont\color{popink}}}

% Signature de fin de cours — poussée tout en bas de la dernière page
\newcommand{\findecours}{%
  \par\vfill\nopagebreak
  \begin{center}\popsquiggle{poporange}\end{center}\vspace{4pt}
  \signatureonbuch
}
\AtBeginDocument{\def\llbracket{\mathopen{[\mkern-3.5mu[}}\def\rrbracket{\mathclose{]\mkern-3.5mu]}}} % intervalles d'entiers (glyphe ⟦ absent de la police)
% Glyphes absents de Fira Math : redéfinis à partir de glyphes disponibles
\AtBeginDocument{%
  \def\setminus{\mathbin{\backslash}}\let\smallsetminus\setminus
  \def\vdots{\mathord{\vbox{\baselineskip3.2pt\lineskiplimit0pt\kern4pt\hbox{.}\hbox{.}\hbox{.}}}}%
  \def\ddots{\mathinner{\mkern1mu\raise6.5pt\hbox{.}\mkern2mu\raise3.6pt\hbox{.}\mkern2mu\raise0.7pt\hbox{.}\mkern1mu}}%
  \def\triangle{\mathord{\scalebox{0.9}{$\symup{\Delta}$}}}%
  \def\nabla{\mathord{\raisebox{\depth}{\scalebox{1}[-1]{$\symup{\Delta}$}}}}%
  \def\checkmark{\popcheck{popdark}}%
  \def\star{\mathbin{\ast}}\def\bigstar{\mathord{\ast}}%
  \def\diamond{\mathbin{\scalebox{0.6}{$\Diamond$}}}%
  \def\bigcup{\mathop{\mathchoice{\raisebox{-0.3em}{\scalebox{1.7}{$\cup$}}}{\scalebox{1.25}{$\cup$}}{\cup}{\cup}}\displaylimits}%
  \def\bigcap{\mathop{\mathchoice{\raisebox{-0.3em}{\scalebox{1.7}{$\cap$}}}{\scalebox{1.25}{$\cap$}}{\cap}{\cap}}\displaylimits}%
  \def\lesseqqgtr{\mathrel{\lessgtr}}\def\bigoplus{\mathop{\oplus}\displaylimits}%
}
\providecommand{\ssi}{si et seulement si} % abréviation fréquente
\AtBeginDocument{\providecommand{\wideparen}{\overparen}} % arc de cercle

%% FIGFIX-BEGIN v5 — correctifs globaux des figures (docs/onbuch-generation/tools/figfix)
\usepackage{adjustbox}\usepackage{etoolbox}\tcbuselibrary{raster}
% toute figure TikZ/pgfplots plus large que la ligne est réduite à la largeur disponible
\BeforeBeginEnvironment{tikzpicture}{\begin{adjustbox}{max width=\linewidth}}
\AfterEndEnvironment{tikzpicture}{\end{adjustbox}}
% légendes pgfplots lisibles par-dessus les courbes
\pgfplotsset{every axis/.append style={legend style={fill=white,fill opacity=0.92,text opacity=1,draw=black!15,inner sep=3pt}}}
% étiquettes d'arêtes (syntaxe edge["..."]) avec fond blanc
\tikzset{every edge quotes/.append style={fill=white,inner sep=1.2pt}}
%% FIGFIX-END
\begin{document}

\pagegarde{Lutte contre la mauvaise alimentation et l'inactivité physique}{SVT}{1ère D}{Module 2 : L'Éducation à la Santé}{ND10}{8 h}{Cette leçon étudie les besoins nutritionnels de l'adolescent et de l'adulte, les conséquences sanitaires de la sédentarité et de la mauvaise alimentation (obésité, diabète de type 2, maladies cardiovasculaires), ainsi que les recommandations pratiques pour une alimentation équilibrée et une activité physique régulière. Elle fournit les outils pour évaluer une ration alimentaire, interpréter l'IMC et élaborer un programme de prévention des maladies liées au mode de vie.
\vspace{4pt}
\begin{multicols}{2}\small
\begin{itemize}
\item À la fin de cette leçon, je sais définir l'équilibre alimentaire et citer les grandes familles de nutriments et leurs rôles
\item À la fin de cette leçon, je sais décrire les besoins énergétiques et nutritionnels spécifiques de l'adolescent et de l'adulte
\item À la fin de cette leçon, je sais évaluer une ration alimentaire réelle par rapport aux apports recommandés
\item À la fin de cette leçon, je sais calculer et interpréter l'IMC et le relier aux risques sanitaires
\item À la fin de cette leçon, je sais expliquer les mécanismes conduisant de la sédentarité et de la malbouffe à l'obésité, au diabète de type 2 et aux maladies cardiovasculaires
\item À la fin de cette leçon, je sais énoncer les recommandations officielles en matière d'activité physique et d'alimentation équilibrée
\end{itemize}\end{multicols}}

\begin{sommaire}
\begin{multicols}{2}
\begin{enumerate}
\item Équilibre alimentaire et besoins nutritionnels
\item L'IMC : calcul et interprétation
\item Conséquences de la sédentarité et de la mauvaise alimentation
\item Recommandations : alimentation équilibrée
\item Recommandations : activité physique
\item Élaborer un programme de prévention
\item Activité d'intégration
\item Méthodes et exercices résolus
\item Exercices
\item Corrigés des exercices
\item Fiche bilan
\end{enumerate}
\end{multicols}
\end{sommaire}

\begin{objectifs}
\begin{itemize}
\item À la fin de cette leçon, je sais définir l'équilibre alimentaire et citer les grandes familles de nutriments et leurs rôles
\item À la fin de cette leçon, je sais décrire les besoins énergétiques et nutritionnels spécifiques de l'adolescent et de l'adulte
\item À la fin de cette leçon, je sais évaluer une ration alimentaire réelle par rapport aux apports recommandés
\item À la fin de cette leçon, je sais calculer et interpréter l'IMC et le relier aux risques sanitaires
\item À la fin de cette leçon, je sais expliquer les mécanismes conduisant de la sédentarité et de la malbouffe à l'obésité, au diabète de type 2 et aux maladies cardiovasculaires
\item À la fin de cette leçon, je sais énoncer les recommandations officielles en matière d'activité physique et d'alimentation équilibrée
\item À la fin de cette leçon, je sais élaborer un programme hebdomadaire combinant alimentation équilibrée et activité physique adapté à un profil donné
\item À la fin de cette leçon, je sais argumenter un message de prévention des maladies liées au mode de vie auprès de ma communauté
\end{itemize}
\end{objectifs}

\begin{prerequis}
\begin{itemize}
\item Les nutriments (glucides, lipides, protéines, vitamines, sels minéraux, eau) vus en classe de Seconde
\item Notions de digestion et d'absorption des aliments (Seconde)
\item Le métabolisme de base et la dépense énergétique (notions de début d'année de Première)
\item Calculs élémentaires : proportionnalité, pourcentages, quotient (mathématiques)
\end{itemize}
\end{prerequis}

\newpage

\coursec{Équilibre alimentaire et besoins nutritionnels}

\textbf{Situation d'accroche.} Il est 6 h 30 à Yaoundé, quartier Mvog-Ada. Amina, 16 ans, élève de Première D, avale en vitesse un beignet-haricot acheté à la vendeuse du carrefour, arrosé d'une boisson sucrée bien fraîche. À midi, au réfectoire, elle déjeune d'un généreux plat de riz à l'huile rouge. Le soir, affalée sur le canapé familial, elle fait défiler les vidéos sur son téléphone jusqu'à 22 h, en grignotant des chips. Lors de la visite médicale scolaire, l'infirmier note un IMC de 29 et fronce les sourcils : la mère d'Amina est diabétique. Pourtant, Amina « mange bien », comme le dit sa grand-mère... Manger beaucoup, est-ce manger bien ? Comment expliquer scientifiquement la situation d'Amina ? Quels risques encourt-elle ? Et surtout, quel programme concret peut-on lui proposer ? Toute la leçon répondra à ces questions ; cette première section pose les fondations : \emph{que signifie « bien s'alimenter », et de quoi notre organisme a-t-il exactement besoin ?}

\courssub{Qu'est-ce que l'équilibre alimentaire ?}

\textbf{Intuition.} Imagine une moto-taxi : pour rouler, elle a besoin d'essence (l'énergie), mais aussi d'huile moteur, de freins en bon état, de pneus gonflés (l'entretien et la construction). Donne-lui uniquement de l'essence en excès et jamais d'entretien : elle finira en panne, le réservoir débordant inutilement. Notre corps fonctionne pareillement : il lui faut du « carburant » (les nutriments énergétiques) \emph{et} des « pièces de rechange et produits d'entretien » (protéines, vitamines, minéraux, eau), dans des proportions précises.

\begin{definition}[Équilibre alimentaire]
L'\cle{équilibre alimentaire} est l'adéquation entre les \textbf{apports alimentaires} (ce que l'on mange et boit) et les \textbf{besoins de l'organisme}, à la fois :
\begin{itemize}
\item en \textbf{quantité} : l'énergie apportée doit couvrir les dépenses (fonctionnement des organes, croissance, activité physique) sans excès durable ;
\item en \textbf{qualité} : chaque catégorie de nutriments (glucides, lipides, protéines, vitamines, minéraux, fibres, eau) doit être présente en proportion adaptée.
\end{itemize}
L'ensemble des aliments consommés en une journée constitue la \cle{ration alimentaire} ; elle est dite \emph{équilibrée} lorsqu'elle satisfait ces deux conditions.
\end{definition}

\begin{attention}[Manger beaucoup n'est pas manger équilibré]
Erreur fréquente, y compris au Probatoire : croire qu'un repas copieux est forcément un bon repas. Un grand plat de riz à l'huile rouge suivi d'une boisson sucrée peut dépasser \SI{1000}{kcal} tout en étant \textbf{déséquilibré} : excès de lipides et de sucres rapides, carence en fibres, en vitamines et en protéines de qualité. L'équilibre se juge sur la \emph{composition} de la ration, pas sur son volume. Retiens : \textbf{quantité suffisante + diversité + bonnes proportions = équilibre}.
\end{attention}

\courssub{Les nutriments énergétiques : le carburant du corps}

Les aliments contiennent trois grandes familles de nutriments capables de fournir de l'énergie une fois digérés et oxydés dans nos cellules (rappelle-toi la respiration cellulaire : \ce{C6H12O6 + 6O2 -> 6CO2 + 6H2O} + énergie).

\begin{propriete}[Valeurs énergétiques des nutriments (admises)]
Les valeurs énergétiques des nutriments sont des constantes physiologiques mesurées expérimentalement par \textbf{calorimétrie} (on mesure la chaleur dégagée par l'oxydation complète d'un gramme de nutriment). On les admet et on les utilise dans tous les calculs de ration :
\begin{itemize}
\item \textbf{Glucides} : \SI{4}{kcal/g} (sucres, amidon — riz, maïs, manioc, igname, pain) ;
\item \textbf{Lipides} : \SI{9}{kcal/g} (graisses et huiles — huile rouge, huile d'arachide, beurre) ;
\item \textbf{Protéines} : \SI{4}{kcal/g} (viande, poisson, œufs, niébé, arachides).
\end{itemize}
Conséquence immédiate : \textbf{1 g de lipides fournit plus du double d'énergie qu'1 g de glucides ou de protéines}. C'est pourquoi les plats très gras (fritures, huile rouge en abondance) font exploser le compteur calorique sans rassasier davantage.
\end{propriete}

\begin{exemplebox}[Rôles respectifs des trois nutriments énergétiques]
\begin{itemize}
\item Les \cle{glucides} sont le carburant « à usage rapide » : cerveau et muscles en sont les grands consommateurs. L'amidon du riz ou du foutou est hydrolysé en glucose.
\item Les \cle{lipides} sont la réserve énergétique concentrée (stockée dans le tissu adipeux) et constituent les membranes cellulaires ; ils transportent les vitamines A, D, E, K.
\item Les \cle{protéines} ont d'abord un rôle \textbf{bâtisseur} (muscles, enzymes, anticorps, hémoglobine) ; elles ne servent de carburant qu'en dernier recours. C'est le nutriment irremplaçable de la croissance de l'adolescent.
\end{itemize}
\end{exemplebox}

\courssub{Les nutriments non énergétiques : petits mais indispensables}

Ils ne fournissent \textbf{aucune calorie}, pourtant leur absence provoque des maladies graves, même chez quelqu'un qui mange « à sa faim ».

\begin{itemize}
\item Les \cle{vitamines} : vitamine C (agrumes, goyave, folong) pour les défenses et l'absorption du fer ; vitamine A (huile rouge, carotte, mangue) pour la vision ; vitamines du groupe B pour le métabolisme énergétique.
\item Les \cle{sels minéraux} : le \textbf{fer} (folong, niébé, viande, poisson) entre dans l'hémoglobine — sa carence provoque l'\textbf{anémie}, fréquente chez la jeune fille réglée ; le \textbf{calcium} (lait, poissons fumés consommés avec les arêtes, légumes verts) construit le squelette ; l'\textbf{iode} (sel iodé, poisson de mer) prévient le goitre.
\item Les \cle{fibres} (légumes, fruits, céréales complètes) : non digérées, elles régulent le transit, prolongent la satiété et freinent l'absorption des sucres.
\item L'\cle{eau} : 1,5 à 2 litres par jour ; aucune réaction biochimique n'est possible sans elle.
\end{itemize}

\begin{savaistu}[Une stratégie locale anti-anémie]
Le fer des végétaux (folong, niébé, feuilles de manioc) est un fer « non héminique », mal absorbé seul (moins de 5 \%). Mais en présence de \textbf{vitamine C}, son absorption est multipliée par 2 à 4 ! Associer un plat de légumes verts à un fruit riche en vitamine C (goyave, orange, citron pressé sur le ndolé) est donc une stratégie simple, gratuite et redoutablement efficace contre l'anémie — particulièrement précieuse pour les adolescentes. À l'inverse, le thé pris pendant le repas freine l'absorption du fer.
\end{savaistu}

\courssub{Les besoins énergétiques : combien de calories par jour ?}

\textbf{Démarche d'investigation.} \emph{Observation} : un joueur de football de 17 ans mange nettement plus qu'une secrétaire de 40 ans, sans grossir. \emph{Hypothèse} : les besoins énergétiques dépendent de l'âge, du sexe et de l'activité physique. \emph{Données} : les mesures de dépense énergétique (calorimétrie indirecte) confirment que la dépense totale = métabolisme de base + croissance (chez l'adolescent) + activité physique. \emph{Conclusion} : les besoins varient fortement d'un individu à l'autre.

\begin{aretenir}[Ordres de grandeur des besoins énergétiques quotidiens]
\begin{itemize}
\item \textbf{Adolescent (10--18 ans)} : environ \num{2500} à \SI{3000}{kcal/j} selon le sexe et l'activité (les garçons et les sportifs en haut de la fourchette). La croissance rapide de la puberté justifie des besoins \emph{supérieurs} à ceux de l'adulte.
\item \textbf{Adulte} : environ \num{2000} à \SI{2500}{kcal/j} selon le sexe et le métier (un manœuvre dépense plus qu'un employé de bureau).
\end{itemize}
L'adolescent a aussi des besoins \textbf{spécifiques accrus} en protéines (construction musculaire), en calcium (ossification : près de la moitié de la masse osseuse adulte se construit pendant l'adolescence) et en fer, surtout chez la \textbf{jeune fille} qui perd du fer chaque mois par les règles.
\end{aretenir}

\begin{center}
\begin{popfigure}
\begin{tikzpicture}
\begin{axis}[
ybar, bar width=14pt,
width=0.95\linewidth, height=6.2cm,
ylabel={Besoins (kcal/j)},
symbolic x coords={Ado fille,Ado garçon,Femme active,Homme actif,Adulte sédentaire},
xtick=data, x tick label style={font=\small, rotate=15, anchor=east},
ymin=0, ymax=3400,
nodes near coords, nodes near coords style={font=\small},
ymajorgrids, grid style={popline!40},
]
\addplot[fill=popblue!70, draw=popblue] coordinates {(Ado fille,2500) (Ado garçon,3000) (Femme active,2200) (Homme actif,2600) (Adulte sédentaire,2000)};
\end{axis}
\end{tikzpicture}
\legende{Besoins énergétiques quotidiens approximatifs selon l'âge, le sexe et le niveau d'activité physique. L'adolescent en croissance a des besoins comparables, voire supérieurs, à ceux d'un adulte actif.}
\end{popfigure}
\end{center}

\courssub{Comment répartir l'énergie dans la ration ?}

Il ne suffit pas d'atteindre le bon total de calories : encore faut-il que cette énergie provienne des bonnes sources, dans les bonnes proportions.

\begin{aretenir}[Répartition énergétique idéale de la ration]
\begin{itemize}
\item \textbf{Glucides : 50--55 \%} de l'énergie totale (en privilégiant les sucres lents : céréales, tubercules) ;
\item \textbf{Lipides : 30--35 \%} (sans excès de graisses animales et de fritures) ;
\item \textbf{Protéines : 12--15 \%} (animales et végétales : poisson, œufs, niébé, arachides).
\end{itemize}
\end{aretenir}

\begin{center}
\begin{popfigure}
\begin{tikzpicture}[scale=1.05]
% Camembert : glucides 52%, lipides 33%, protéines 15%
\fill[popblue!70] (0,0) -- (90:2.2) arc (90:90-187.2:2.2) -- cycle;
\fill[poporange!75] (0,0) -- (90-187.2:2.2) arc (90-187.2:90-306:2.2) -- cycle;
\fill[popgreen!70] (0,0) -- (90-306:2.2) arc (90-306:90-360:2.2) -- cycle;
\draw[popdark, thick] (0,0) circle (2.2);
\node[font=\small\bfseries, popdark] at (20:1.35) {Glucides};
\node[font=\small\bfseries, popdark] at (-40:1.5) {50--55 \%};
\node[font=\small\bfseries, popdark] at (150:1.35) {Lipides};
\node[font=\small\bfseries, popdark] at (185:1.5) {30--35 \%};
\node[font=\small\bfseries, popdark] at (250:1.35) {Protéines};
\node[font=\small\bfseries, popdark] at (285:1.6) {12--15 \%};
\end{tikzpicture}
\legende{Diagramme circulaire de la répartition énergétique idéale de la ration quotidienne : les glucides fournissent environ la moitié de l'énergie, les lipides environ un tiers, les protéines le reste.}
\end{popfigure}
\end{center}

Cette répartition dans la journée se traduit concrètement par une distribution entre les repas :

\begin{center}
\begin{tabularx}{\linewidth}{|l|X|X|}
\hline
\rowcolor{popblueL} \textbf{Repas} & \textbf{Part de l'énergie quotidienne} & \textbf{Exemple adapté} \\
\hline
Petit-déjeuner & environ 25 \% & Bouillie de mil enrichie, pain, œuf, fruit \\
\hline
Déjeuner & 35--40 \% & Riz + ndolé + poisson + fruit \\
\hline
Dîner & 25--30 \% & Repas plus léger : igname + légumes + haricots \\
\hline
Collation(s) & 5--10 \% & Fruit, arachides grillées (et non chips ni boisson sucrée) \\
\hline
\end{tabularx}
\end{center}

\begin{attention}[Le piège du petit-déjeuner sauté ou sucré]
Sauter le petit-déjeuner (ou le réduire à un beignet et une boisson sucrée comme Amina) provoque hypoglycémie et fatigue en classe, puis compulsion sucrée dans la journée. Un petit-déjeuner doit apporter environ le \textbf{quart} de l'énergie quotidienne, avec des sucres lents et des protéines.
\end{attention}

\courssub{Diversité alimentaire et pyramide alimentaire camerounaise}

Une \cle{ration équilibrée} repose sur la \textbf{diversité} : aucun aliment seul ne contient tous les nutriments. La pyramide alimentaire hiérarchise les priorités : à la base, les aliments à consommer en plus grande quantité ; au sommet, ceux à limiter.

\begin{center}
\begin{popfigure}
\begin{tikzpicture}[scale=1]
% Pyramide en 4 étages
\fill[popgreen!60] (-3,0) -- (3,0) -- (2.1,1.4) -- (-2.1,1.4) -- cycle;
\fill[popgold!65] (-2.1,1.4) -- (2.1,1.4) -- (1.3,2.7) -- (-1.3,2.7) -- cycle;
\fill[popblue!55] (-1.3,2.7) -- (1.3,2.7) -- (0.6,3.9) -- (-0.6,3.9) -- cycle;
\fill[poporange!70] (-0.6,3.9) -- (0.6,3.9) -- (0,5.1) -- cycle;
\draw[popdark, thick] (-3,0) -- (0,5.1) -- (3,0) -- cycle;
\node[font=\footnotesize, align=center, popdark] at (0,0.65) {\textbf{Base} : céréales et tubercules\\ (riz, maïs, manioc, igname, mil) + eau};
\node[font=\footnotesize, align=center, popdark] at (0,2.0) {Légumes et fruits\\ (folong, ndolé, goyave, mangue)};
\node[font=\footnotesize, align=center, popdark] at (0,3.25) {Protéines : poisson, œufs,\\ niébé, arachides, lait};
\node[font=\footnotesize, align=center, popdark] at (0,4.45) {À limiter :};
\node[font=\footnotesize, align=center, popdark] at (0,4.15) {huiles, sucre, sel};
% flèches latérales corrigées
\draw[-{Stealth}, popgreen!80!black, thick] (-3.3,0.7) -- (-3.3,2.6) node[midway, left, font=\footnotesize, align=right, popdark]{en plus grande\\ quantité};
\draw[-{Stealth}, poporange!90!black, thick] (3.3,2.6) -- (3.3,0.7) node[midway, right, font=\footnotesize, align=left, popdark]{avec\\ modération};
\end{tikzpicture}
\legende{Pyramide alimentaire adaptée aux aliments du Cameroun : la base (céréales, tubercules, eau) est consommée en plus grande quantité ; le sommet (graisses, sucre, sel) doit rester exceptionnel.}
\end{popfigure}
\end{center}

\begin{exemplebox}[Le niébé, champion local]
Le niébé (haricot blanc) illustre parfaitement la diversité à bas coût : riche en protéines végétales, en fer, en fibres et en glucides lents, il complète idéalement un plat de riz ou de maïs (les acides aminés des céréales et des légumineuses se complètent). Un repas « riz + niébé + folong + une goyave » couvre quasiment tous les besoins pour un coût modeste : l'équilibre alimentaire n'est pas une question de richesse mais de choix.
\end{exemplebox}

\courssub{Évaluer une ration alimentaire : la méthode}

\begin{methode}[Évaluer une ration alimentaire en 4 étapes]
\begin{enumerate}
\item \textbf{Inventorier} tous les aliments et boissons consommés dans la journée, avec leurs quantités (en grammes) ;
\item \textbf{Convertir} chaque aliment en nutriments puis en énergie à l'aide d'une table de composition des aliments et des valeurs 4--9--4 kcal/g ;
\item \textbf{Comparer} le total obtenu aux apports recommandés (quantité totale en kcal, répartition 50--55 / 30--35 / 12--15, présence de vitamines, minéraux, fibres) ;
\item \textbf{Conclure} (ration équilibrée, excessive, carencée ?) et \textbf{proposer des corrections} concrètes et réalistes.
\end{enumerate}
\end{methode}

\begin{exoresolu}[Le déjeuner d'Amina est-il équilibré ?]
Énoncé : Amina déjeune d'un plat composé de 150 g de riz cuit ($\approx$ \SI{195}{kcal}), de 100 g de ndolé préparé avec arachides et huile ($\approx$ \SI{250}{kcal}) et de 50 g de poisson fumé ($\approx$ \SI{100}{kcal}). (a) Calculer l'énergie totale du repas. (b) La comparer aux besoins d'un adolescent (environ \SI{2500}{kcal/j}) et à la part attendue d'un déjeuner. (c) Ce repas est-il équilibré qualitativement ?
\tcblower
Solution détaillée :
\begin{itemize}
\item (a) Énergie totale : $195 + 250 + 100 = \SI{545}{kcal}$.
\item (b) Part des besoins quotidiens : $\dfrac{545}{2500} \times 100 \approx 22\,\%$. Or un déjeuner devrait apporter 35--40 \% de l'énergie quotidienne, soit $0{,}35 \times 2500 \approx \SI{875}{kcal}$. Ce déjeuner est donc \textbf{quantitativement insuffisant} pour le repas principal — ce qui pousse Amina à grignoter le soir des aliments très caloriques et pauvres en nutriments.
\item (c) Qualitativement, le repas est plutôt bon : glucides (riz), protéines (poisson fumé), lipides et protéines végétales (arachides du ndolé), fibres et fer (feuilles de ndolé). \textbf{Correction proposée} : augmenter la portion de riz ou ajouter un fruit (goyave ou orange, dont la vitamine C améliorera l'absorption du fer du ndolé) et un verre d'eau plutôt qu'une boisson sucrée.
\end{itemize}
Conclusion : ce plat traditionnel est de bonne qualité ; c'est l'ensemble de la journée d'Amina (petit-déjeuner sucré et gras, grignotage du soir, absence d'activité physique) qui est déséquilibré — ce que nous analyserons dans les sections suivantes.
\end{exoresolu}

\begin{exercice}[À toi de jouer]
Kamdem, 17 ans, consomme en une journée 400 g de glucides, 120 g de lipides et 90 g de protéines. (1) Calcule son apport énergétique total. (2) Vérifie si la répartition glucides/lipides/protéines respecte les proportions recommandées. (3) Formule un conseil.
\end{exercice}

\begin{corrige}[Exercice : ration de Kamdem]
(1) Glucides : $400 \times 4 = \SI{1600}{kcal}$ ; lipides : $120 \times 9 = \SI{1080}{kcal}$ ; protéines : $90 \times 4 = \SI{360}{kcal}$. Total : $1600 + 1080 + 360 = \SI{3040}{kcal}$.
(2) Proportions : glucides $= 1600/3040 \approx 53\,\%$ (correct, 50--55 \%) ; lipides $= 1080/3040 \approx 35,5\,\%$ (légèrement au-dessus des 30--35 \%) ; protéines $= 360/3040 \approx 11,8\,\%$ (légèrement en dessous du minimum recommandé de 12 \%).
(3) Conseil : l'apport total convient à un adolescent très actif, mais Kamdem devrait réduire légèrement les lipides (moins de fritures et d'huile) et augmenter légèrement la part protéique (poisson, œufs, niébé) pour atteindre au moins 12 \%, tout en veillant à la diversité (fruits, légumes, eau).
\end{corrige}

\begin{aretenir}[L'essentiel de la section]
L'équilibre alimentaire, c'est la bonne \textbf{quantité} d'énergie (adolescent : \num{2500}--\SI{3000}{kcal/j} ; adulte : \num{2000}--\SI{2500}{kcal/j}) \emph{et} la bonne \textbf{qualité} : 50--55 \% de glucides, 30--35 \% de lipides, 12--15 \% de protéines, plus vitamines, minéraux (fer, calcium, iode), fibres et eau. L'adolescent en croissance a des besoins accrus en protéines, calcium et fer (surtout la jeune fille). La diversité alimentaire, accessible avec les aliments locaux (niébé, folong, poisson, fruits), est la clé — et non la quantité servie dans l'assiette.
\end{aretenir}

\coursec{L'IMC : calcul et interprétation}

Dans la section précédente, nous avons appris à construire une alimentation équilibrée et à estimer les besoins nutritionnels d'un adolescent ou d'un adulte. Mais une question pratique se pose immédiatement : comment savoir si, concrètement, une personne mange trop, pas assez, ou correctement par rapport à ses besoins ? Comment un médecin, à l'hôpital de district de Bafoussam ou dans un centre de santé de Yaoundé, peut-il évaluer rapidement l'état nutritionnel d'un patient ? Il lui faut un \cle{indicateur} simple, mesurable et comparable d'une personne à l'autre. Cet indicateur existe : c'est l'\cle{Indice de Masse Corporelle}, ou \cle{IMC}.

\courssub{Pourquoi un indice ? De l'observation à la définition}

\textbf{Situation concrète.} Deux hommes pèsent chacun 85 kg. Le premier mesure 1,90 m : il paraît mince, presque maigre. Le second mesure 1,60 m : il paraît nettement corpulent. Pourtant, la balance affiche la même masse ! Cette observation banale montre que la \cle{masse} seule ne suffit pas à juger de la corpulence : il faut la rapporter à la \cle{taille}.

\textbf{Démarche d'investigation.}
\begin{itemize}
\item \emph{Observation} : à masse égale, plus la taille est grande, plus la corpulence est faible.
\item \emph{Hypothèse} : un rapport entre la masse et la taille devrait permettre de comparer objectivement les corpulences.
\item \emph{Données} : les études épidémiologiques de l'Organisation mondiale de la santé (OMS) ont montré que le rapport masse/taille$^2$ est celui qui se corrèle le mieux avec la masse grasse et avec la fréquence des maladies métaboliques dans les populations.
\item \emph{Conclusion} : ce rapport, appelé IMC, est adopté depuis les années 1990 comme outil international de dépistage.
\end{itemize}

\begin{definition}[Indice de Masse Corporelle]
L'\cle{Indice de Masse Corporelle} (IMC) est le rapport de la masse d'une personne au carré de sa taille :
\[ \text{IMC} = \frac{m}{t^2} \]
avec $m$ la masse en kilogrammes (kg) et $t$ la taille en mètres (m). L'IMC s'exprime en \si{kg/m^2}. C'est un indicateur de \cle{corpulence} utilisé par l'OMS pour évaluer l'état nutritionnel des adultes.
\end{definition}

\begin{attention}[Pourquoi le carré de la taille ?]
On divise par $t^2$ et non par $t$ car la masse d'un corps augmente approximativement avec le volume, donc avec une puissance de la taille voisine de 2. Le carré permet ainsi de « neutraliser » l'effet de la taille : deux personnes de même corpulence mais de tailles différentes obtiennent des IMC comparables.
\end{attention}

\courssub{Calculer un IMC : méthode pas à pas}

\begin{methode}[Calculer et interpréter un IMC]
\begin{enumerate}
\item \textbf{Convertir} la taille en mètres (si elle est donnée en centimètres, diviser par 100 : 158 cm $= 1{,}58$ m).
\item \textbf{Élever au carré} la taille en mètres : calculer $t^2$.
\item \textbf{Diviser} la masse (en kg) par ce carré : $\text{IMC} = m/t^2$.
\item \textbf{Situer} le résultat dans les classes d'interprétation de l'OMS.
\item \textbf{Conclure} sur l'état nutritionnel et le niveau de risque sanitaire associé.
\end{enumerate}
\end{methode}

\begin{exemplebox}[L'IMC d'Amina]
Amina, élève de Première D, pèse 68 kg et mesure 1,58 m. Calculons son IMC en suivant la méthode :
\begin{enumerate}
\item La taille est déjà en mètres : $t = 1{,}58$ m.
\item Carré de la taille : $t^2 = (1{,}58)^2 = 1{,}58 \times 1{,}58 = 2{,}4964$ m$^2$.
\item Division : $\text{IMC} = \dfrac{68}{2{,}4964} \approx 27{,}2$ \si{kg/m^2}.
\item $27{,}2$ se situe entre 25 et 29,9 : Amina est en \cle{surpoids}.
\item Conclusion : Amina présente un risque \emph{modéré} de maladies métaboliques (diabète de type 2, hypertension). Une alimentation équilibrée et une activité physique régulière lui sont recommandées.
\end{enumerate}
\end{exemplebox}

\begin{attention}[Les deux erreurs classiques au Probatoire]
\begin{itemize}
\item \textbf{Oublier le carré} : calculer $68 / 1{,}58 \approx 43$ au lieu de $68 / 2{,}4964 \approx 27{,}2$. Le résultat est faux et conduit à une conclusion aberrante.
\item \textbf{Laisser la taille en centimètres} : $68 / (158)^2 \approx 0{,}0027$... ou pire, diviser par 158 sans carré. En fait, utiliser la taille en cm au lieu de m donne un résultat \textbf{10 000 fois trop petit} (car $100^2 = 10\,000$) ; inversement, oublier la conversion dans l'autre sens donne un résultat 10 000 fois trop grand. Toujours vérifier que la taille est en \textbf{mètres} et qu'elle est \textbf{élevée au carré}.
\end{itemize}
Un bon réflexe : un IMC humain plausible se situe entre 15 et 45 \si{kg/m^2}. Tout résultat hors de cette fourchette doit te faire reprendre le calcul.
\end{attention}

\courssub{Les classes d'interprétation de l'OMS}

L'OMS a défini des seuils internationaux qui relient la valeur de l'IMC à l'état nutritionnel et au niveau de risque pour la santé.

\begin{aretenir}[Classes d'IMC selon l'OMS (adulte)]
\begin{center}
\begin{tabularx}{\linewidth}{|c|X|X|}\hline
\rowcolor{popblueL} \textbf{IMC (kg/m$^2$)} & \textbf{Classification} & \textbf{Niveau de risque} \\ \hline
$< 18{,}5$ & Maigreur (insuffisance pondérale) & Risque accru (carences, infections) \\ \hline
$18{,}5$ à $24{,}9$ & Corpulence normale & Risque minimal \\ \hline
$25$ à $29{,}9$ & Surpoids & Risque modéré \\ \hline
$30$ à $34{,}9$ & Obésité modérée (classe I) & Risque élevé \\ \hline
$35$ à $39{,}9$ & Obésité sévère (classe II) & Risque très élevé \\ \hline
$\geq 40$ & Obésité massive (classe III) & Risque extrêmement élevé \\ \hline
\end{tabularx}
\end{center}
\end{aretenir}

\begin{popfigure}
\begin{tikzpicture}[x=1cm,y=1cm]
% Thermomètre des classes d'IMC
% Zones colorées : échelle de 15 à 45, 1 unité = 0.22 cm
\def\echelle{0.22}
% Maigreur 15-18.5 : bleu
\fill[popblue!60] (0,{15*\echelle}) rectangle (1.6,{18.5*\echelle});
% Normale 18.5-25 : vert
\fill[popgreen!55] (0,{18.5*\echelle}) rectangle (1.6,{25*\echelle});
% Surpoids 25-30 : orange clair
\fill[popgold!70] (0,{25*\echelle}) rectangle (1.6,{30*\echelle});
% Obésité modérée 30-35 : orange
\fill[poporange!75] (0,{30*\echelle}) rectangle (1.6,{35*\echelle});
% Obésité sévère 35-40 : rouge-orangé
\fill[popink!75] (0,{35*\echelle}) rectangle (1.6,{40*\echelle});
% Obésité massive 40-45 : rouge foncé
\fill[popink!100!popdark] (0,{40*\echelle}) rectangle (1.6,{45*\echelle});
% Contour
\draw[popdark,thick] (0,{15*\echelle}) rectangle (1.6,{45*\echelle});
% Graduations et seuils
\foreach \v in {15,18.5,25,30,35,40,45}{
  \draw[popdark] (0,{\v*\echelle}) -- (-0.15,{\v*\echelle});
  \node[left,font=\small] at (-0.2,{\v*\echelle}) {\v};
}
% Étiquettes des zones
\node[right,font=\small\bfseries,text=popblue] at (1.8,{16.75*\echelle}) {Maigreur};
\node[right,font=\small\bfseries,text=popgreen!60!popdark] at (1.8,{21.75*\echelle}) {Corpulence normale};
\node[right,font=\small\bfseries,text=popgold!70!popdark] at (1.8,{27.5*\echelle}) {Surpoids};
\node[right,font=\small\bfseries,text=poporange!80!popdark] at (1.8,{32.5*\echelle}) {Obésité modérée};
\node[right,font=\small\bfseries,text=popink] at (1.8,{37.5*\echelle}) {Obésité sévère};
\node[right,font=\small\bfseries,text=popink!80!popdark] at (1.8,{42.5*\echelle}) {Obésité massive};
% Risques
\node[right,font=\scriptsize,text=popmuted] at (6.3,{16.75*\echelle}) {risque accru};
\node[right,font=\scriptsize,text=popmuted] at (6.3,{21.75*\echelle}) {risque minimal};
\node[right,font=\scriptsize,text=popmuted] at (6.3,{27.5*\echelle}) {risque modéré};
\node[right,font=\scriptsize,text=popmuted] at (6.3,{32.5*\echelle}) {risque élevé};
\node[right,font=\scriptsize,text=popmuted] at (6.3,{37.5*\echelle}) {risque très élevé};
\node[right,font=\scriptsize,text=popmuted] at (6.3,{42.5*\echelle}) {risque extrême};
% Position d'Amina
\draw[popdark,very thick,->] (8.6,{27.2*\echelle}) -- (1.7,{27.2*\echelle});
\node[above,font=\scriptsize,text=popdark] at (5.2,{27.2*\echelle}) {Amina : 27,2};
% Titre axe
\node[rotate=90,font=\small] at (-1.1,{30*\echelle}) {IMC (kg/m$^2$)};
\end{tikzpicture}
\legende{« Thermomètre » des classes d'IMC selon l'OMS. Chaque zone colorée correspond à une classe de corpulence et à un niveau de risque sanitaire croissant, du bleu (maigreur) au rouge foncé (obésité massive). La flèche indique la position d'Amina (IMC = 27,2), dans la zone de surpoids.}
\end{popfigure}

\courssub{IMC élevé et risques sanitaires : un gradient croissant}

L'IMC n'est pas qu'un chiffre esthétique : c'est un \cle{facteur prédictif} de maladies. Les grandes études épidémiologiques montrent que le \cle{risque relatif} de développer un diabète de type 2, une hypertension artérielle ou une maladie cardiovasculaire augmente progressivement — et de façon de plus en plus rapide — avec l'IMC. On parle de \cle{gradient de risque} : il n'y a pas de seuil brutal, mais une aggravation continue.

\begin{popfigure}
\begin{tikzpicture}
\begin{axis}[
  width=\linewidth, height=7cm,
  xlabel={IMC (\si{kg/m^2})},
  ylabel={Risque relatif de diabète de type 2},
  xmin=17, xmax=42, ymin=0, ymax=30,
  xtick={18.5,25,30,35,40},
  ytick={1,5,10,15,20,25,30},
  grid=both, grid style={popline!40},
  legend style={at={(0.03,0.97)},anchor=north west,font=\small},
]
% Courbe exponentielle du risque relatif (référence : IMC 22 -> RR = 1)
\addplot[popcurve,domain=18:41,samples=200]{exp(0.16*(x-22))};
\addlegendentry{Risque relatif (référence : IMC = 22)}
% Ligne de référence RR = 1
\addplot[popmuted,dashed,domain=17:42]{1};
% Zones de seuils
\draw[popgreen!60!popdark,dashed] (axis cs:18.5,0) -- (axis cs:18.5,30);
\draw[popgold!80!popdark,dashed] (axis cs:25,0) -- (axis cs:25,30);
\draw[popink,dashed] (axis cs:30,0) -- (axis cs:30,30);
\node[font=\scriptsize,text=popgreen!50!popdark,rotate=90,anchor=south] at (axis cs:18.5,1) {18,5};
\node[font=\scriptsize,text=popgold!70!popdark,rotate=90,anchor=south] at (axis cs:25,1) {25};
\node[font=\scriptsize,text=popink,rotate=90,anchor=south] at (axis cs:30,1) {30};
\end{axis}
\end{tikzpicture}
\legende{Évolution du risque relatif de diabète de type 2 en fonction de l'IMC (courbe schématique inspirée des grandes études épidémiologiques). Le risque est pris égal à 1 pour un IMC de 22 \si{kg/m^2} ; il croît de façon exponentielle : il est multiplié par environ 5 à un IMC de 32 et par plus de 15 au-delà de 38. Les lignes verticales rappellent les seuils OMS (18,5 ; 25 ; 30).}
\end{popfigure}

\textbf{Interprétation.} La courbe montre une croissance \emph{exponentielle} : passer d'un IMC de 25 à 30 augmente modérément le risque, mais passer de 35 à 40 le fait exploser. C'est pourquoi les obésités sévère et massive sont des urgences de santé publique. Le même gradient existe pour l'\cle{hypertension artérielle} et les \cle{maladies cardiovasculaires} (infarctus, accidents vasculaires cérébraux), car l'excès de masse grasse favorise l'\cle{athérosclérose} (dépôts de graisse sur la paroi des artères) et la \cle{résistance à l'insuline}.

\courssub{Les limites de l'IMC : un outil à manier avec discernement}

L'IMC est un indicateur de \emph{population} et de \emph{dépistage}, pas un diagnostic individuel complet.

\begin{attention}[Quand l'IMC trompe]
\begin{itemize}
\item L'IMC \textbf{ne distingue pas la masse grasse de la masse musculaire}. Or le muscle est plus dense que la graisse.
\item Un \textbf{sportif de haut niveau} (haltérophile, rugbyman) peut afficher un IMC de 28 — classé « surpoids » — alors que sa masse grasse est très faible. Son IMC élevé traduit ses muscles, pas un risque sanitaire.
\item L'IMC n'est \textbf{pas applicable aux femmes enceintes} (la prise de masse est physiologique), ni aux enfants et adolescents sans tables de référence par âge, ni aux personnes très âgées ou très dénutries.
\end{itemize}
\end{attention}

Pour affiner l'évaluation, les médecins utilisent un complément simple et exigible au Probatoire : le \cle{tour de taille}.

\begin{aretenir}[Le tour de taille, indicateur de graisse abdominale]
La graisse localisée autour de l'abdomen (graisse \cle{viscérale}) est la plus dangereuse : elle entoure les organes et sécrète des substances qui favorisent le diabète de type 2 et les maladies cardiovasculaires. On la dépiste en mesurant le tour de taille avec un mètre ruban, à mi-chemin entre la dernière côte et la crête iliaque. Le risque métabolique est \textbf{accru} au-delà de :
\begin{itemize}
\item \textbf{94 cm chez l'homme} ;
\item \textbf{80 cm chez la femme}.
\end{itemize}
Un IMC normal avec un tour de taille élevé signale déjà un danger ; inversement, un léger surpoids avec un tour de taille normal est moins inquiétant.
\end{aretenir}

\begin{exoresolu}[À toi de juger : trois cas cliniques]
\textbf{Énoncé.} Trois adultes consultent au centre de santé :
\begin{itemize}
\item M. Etoundi : 92 kg pour 1,75 m, tour de taille 105 cm ;
\item Mme Fotso : 55 kg pour 1,65 m, tour de taille 76 cm ;
\item M. Abena, joueur de football professionnel : 88 kg pour 1,78 m, tour de taille 82 cm.
\end{itemize}
Pour chacun : calculer l'IMC, le classer selon l'OMS, puis conclure en tenant compte du tour de taille et des limites de l'IMC.
\tcblower
\textbf{Solution détaillée.}

\textbf{M. Etoundi} : $t^2 = (1{,}75)^2 = 3{,}0625$ m$^2$ ; $\text{IMC} = 92 / 3{,}0625 \approx 30{,}0$ \si{kg/m^2}. Il est en \textbf{obésité modérée} (seuil de 30 atteint). Son tour de taille (105 cm $>$ 94 cm) confirme un excès de graisse abdominale : \textbf{risque élevé} de diabète de type 2 et de maladies cardiovasculaires. Une prise en charge associant alimentation équilibrée et activité physique s'impose.

\textbf{Mme Fotso} : $t^2 = (1{,}65)^2 = 2{,}7225$ m$^2$ ; $\text{IMC} = 55 / 2{,}7225 \approx 20{,}2$ \si{kg/m^2}. \textbf{Corpulence normale} (18,5--24,9). Tour de taille 76 cm $<$ 80 cm : pas de graisse abdominale excessive. \textbf{Risque minimal} ; il suffit de maintenir les bonnes habitudes.

\textbf{M. Abena} : $t^2 = (1{,}78)^2 = 3{,}1684$ m$^2$ ; $\text{IMC} = 88 / 3{,}1684 \approx 27{,}8$ \si{kg/m^2}. Classé « \textbf{surpoids} » par l'IMC... mais c'est un sportif de haut niveau : sa masse élevée est musculaire, et son tour de taille (82 cm $<$ 94 cm) montre l'absence de graisse abdominale. \textbf{Conclusion nuancée} : l'IMC seul induit en erreur ; son risque réel est faible. Ce cas illustre parfaitement les limites de l'indicateur.
\end{exoresolu}

\begin{savaistu}[L'obésité, une épidémie qui n'épargne pas le Cameroun]
L'OMS qualifie l'obésité d'\emph{épidémie mondiale} : le nombre de personnes obèses a presque triplé depuis 1975, et plus de 1,9 milliard d'adultes étaient en surpoids en 2016. Contrairement aux idées reçues, ce n'est plus un problème de pays riches : au Cameroun, les enquêtes récentes montrent que \textbf{plus d'un adulte urbain sur quatre est en surpoids ou obèse}, avec une nette prédominance chez les femmes. Les causes ? L'urbanisation, la sédentarité (transport motorisé, écrans), et la transition alimentaire : plats très gras et très sucrés, boissons sucrées, restauration rapide qui côtoient les marchés traditionnels de Douala et de Yaoundé. Le diabète de type 2 qui en résulte surcharge déjà les services hospitaliers. Comprendre et utiliser l'IMC, c'est se donner un outil simple de prévention, pour soi et pour sa communauté.
\end{savaistu}

\begin{exercice}[Pour t'entraîner]
\begin{enumerate}
\item Calcule l'IMC d'un adolescent adulte de 50 kg mesurant 1,70 m. Dans quelle classe se situe-t-il ?
\item Un patient a un IMC de 36,5 \si{kg/m^2}. Classe cette valeur et précise le niveau de risque.
\item Explique pourquoi une femme enceinte de 7 mois ne doit pas être évaluée par l'IMC.
\item Un homme a un IMC de 23 \si{kg/m^2} mais un tour de taille de 98 cm. Son risque est-il nul ? Justifie.
\end{enumerate}
\end{exercice}

\begin{corrige}[Exercice]
\begin{enumerate}
\item $t^2 = (1{,}70)^2 = 2{,}89$ m$^2$ ; $\text{IMC} = 50 / 2{,}89 \approx 17{,}3$ \si{kg/m^2} $< 18{,}5$ : \textbf{maigreur} (insuffisance pondérale), risque de carences.
\item IMC $= 36{,}5$ : entre 35 et 39,9, il s'agit d'une \textbf{obésité sévère} (classe II), à risque très élevé de diabète de type 2, d'hypertension et de maladies cardiovasculaires.
\item La prise de masse pendant la grossesse est physiologique (fœtus, placenta, liquide amniotique) : l'IMC ne reflète plus la corpulence propre de la femme et conduirait à une fausse classification.
\item Non. Malgré un IMC normal, un tour de taille de 98 cm ($> 94$ cm) révèle un excès de \textbf{graisse viscérale abdominale}, facteur de risque métabolique indépendant. Les deux indicateurs doivent être croisés.
\end{enumerate}
\end{corrige}

En résumé, l'IMC est un outil simple, rapide et puissant pour repérer les excès ou les déficits de corpulence, à condition de le calculer correctement (taille en mètres, élevée au carré) et de l'interpréter avec discernement, en le complétant par le tour de taille. Fort de cet indicateur, nous pouvons maintenant examiner dans la section suivante les conséquences précises de la sédentarité et de la mauvaise alimentation sur l'organisme.

\coursec{Conséquences de la sédentarité et de la mauvaise alimentation}

Dans la section précédente, nous avons vu comment l'\cle{IMC} permet de dépister un excès de poids. Mais derrière ce chiffre se cache une réalité physiologique : un déséquilibre durable entre ce que nous mangeons et ce que nous dépensons. Comprenons comment ce déséquilibre, associé à l'inactivité, construit silencieusement les maladies chroniques qui frappent de plus en plus nos villes, de Douala à Yaoundé en passant par Bafoussam.

\courssub{Le bilan énergétique : quand l'apport dépasse la dépense}

\emph{Observation.} Un adolescent qui consomme chaque jour une bouteille de soda (65 cl) et un paquet de chips en plus de ses repas normaux, tout en passant ses journées assis en classe puis devant son téléphone, prend du poids mois après mois. Pourquoi ?

\begin{experience}[Bilan énergétique simplifié sur une journée]
\textbf{Matériel :} Tables de composition des aliments (Ciqual, ANSES), données de dépense énergétique (MET).
\textbf{Protocole :}
\begin{enumerate}
    \item Estimer le \textbf{besoin énergétique journalier (BEJ)} d'un adolescent de 16 ans, 60 kg, 1,70 m (garçon) :
    \begin{itemize}
        \item Métabolisme de base (formule de Mifflin-St Jeor) : $10 \times 60 + 6,25 \times 170 - 5 \times 16 + 5 = 1\,587,5 \approx 1\,590$ kcal.
        \item Niveau d'activité physique (NAP) sédentaire ($1,4$) $\rightarrow$ BEJ $= 1\,590 \times 1,4 \approx 2\,230$ kcal.
    \end{itemize}
    \item Calculer les \textbf{apports} d'une journée type « moderne » :
    \begin{itemize}
        \item Petit-déjeuner (pain + chocolat + lait) : 500 kcal.
        \item Déjeuner (riz + sauce arachide + viande + plantain) : 900 kcal.
        \item Dîner (pâtes + omelette) : 600 kcal.
        \item \textbf{Collations} (1 bouteille soda 65 cl + 1 paquet chips 50 g) $\approx 280 + 270 = 550$ kcal.
    \end{itemize}
    Total apports $\approx 2\,550$ kcal.
    \item Comparer Apports ($2\,550$ kcal) et Dépenses ($2\,230$ kcal).
\end{enumerate}
\textbf{Observations :} Excédent quotidien de $320$ kcal.
\textbf{Interprétation :} Cet excédent, s'il est maintenu, est stocké. 1 kg de tissu adipeux $\approx 7\,700$ kcal (voir Propriété ci-dessous). $320 \times 365 / 7\,700 \approx 15,2$ kg de graisse par an (théorique). En réalité, la prise de poids augmente la dépense de base, limitant la prise, mais la tendance est inexorable.
\textbf{Conclusion :} Un excès modéré mais \textbf{quotidien et durable} suffit à installer un surpoids puis une obésité.
\end{experience}

\begin{propriete}[Équivalence énergie / masse graisseuse -- admise]
Le stockage de \textbf{1 kg de graisse corporelle} (triglycérides dans les adipocytes) correspond à un excès énergétique net d'environ \textbf{7 700 kcal}.
\par\medskip
\textbf{Conséquence pratique :} Un excès constant de seulement \textbf{100 kcal par jour} (l'équivalent d'un petit verre de soda, ou 2 biscuits, ou 10 g d'huile en trop) conduit théoriquement à une prise de \textbf{4,7 kg de graisse par an} ($100 \times 365 / 7\,700 \approx 4,7$). C'est la \textbf{thermodynamique du vivant} : l'énergie ne se perd pas, elle se stocke.
\end{propriete}

\begin{definition}[Sédentarité et inactivité physique]
La \textbf{sédentarité} est un mode de vie caractérisé par une \textbf{dépense énergétique d'activité très faible} (comportements éveillés en position assise/allongée $\le 1,5$ MET), correspondant à moins de \textbf{30 minutes d'activité d'intensité modérée par jour} (marche rapide, vélo, travaux ménagers soutenus). Elle se distingue de l'\textbf{inactivité physique} (non-respect des recommandations d'activité physique, soit $< 150$ min/semaine d'intensité modérée) : on peut être « non sportif » mais actif (marche pour aller à l'école, travaux des champs), ou « sportif » (3h de foot/semaine) mais sédentaire le reste du temps (8 h assis en classe/bureau, écrans).
\end{definition}

\courssub{L'obésité : définition, causes et conséquences multiples}

\begin{definition}[Obésité (adulte et adolescent)]
Excès de masse grasse définie chez l'adulte par un \textbf{IMC $\ge 30$ kg/m$^2$}. Chez l'adolescent, on utilise les courbes de référence IMC-âge (OMS / INSP Cameroun) : obésité = IMC $> +2$ ET (équivalent IMC 30 à 18 ans). L'obésité abdominale (viscérale) est mieux prédite par le \textbf{tour de taille} ($\ge 94$ cm homme / $\ge 80$ cm femme en Afrique subsaharienne).
\end{definition}

\begin{methode}[Mécanisme d'installation de l'obésité et complications métaboliques]
\begin{enumerate}
    \item \textbf{Excès calorique chronique} : apports $>$ dépenses (aliments ultra-transformés riches en sucres ajoutés et graisses cachées, boissons sucrées, portions excessives, grignotage).
    \item \textbf{Sédentarité} : baisse drastique de la dépense énergétique d'activité (thermogénèse d'activité non sportive -- NEAT).
    \item \textbf{Stockage adipocytaire} : les adipocytes grossissent (hypertrophie) puis se multiplient (hyperplasie). Le tissu adipeux viscéral devient un \textbf{organe endocrinien actif} sécrétant des adipokines pro-inflammatoires (leptine, TNF-$\alpha$, IL-6, résistine) et moins d'adiponectine (protectrice).
    \item \textbf{Inflammation chronique de bas grade} : ces molécules perturbent la signalisation de l'insuline dans les muscles, le foie et le tissu adipeux lui-même $\rightarrow$ \textbf{insulinorésistance} systémique.
    \item \textbf{Dyslipidémie athérogène} : résistance à l'insuline hépatique $\rightarrow$ surproduction de VLDL (triglycérides $\uparrow$), baisse du HDL-cholestérol, LDL dense et oxydé.
\end{enumerate}
\end{methode}

\begin{exemplebox}[Conséquences systémiques de l'obésité]
\begin{itemize}
    \item \textbf{Articulaires :} Surcharge mécanique $\rightarrow$ gonarthrose (genou), coxarthrose (hanche), lombalgies chroniques. Fréquent au Cameroun chez les porteurs de charges et les personnes en surpoids.
    \item \textbf{Respiratoires :} Syndrome d'apnées du sommeil (SAS) par compression des voies aériennes supérieures par la graisse cervicale $\rightarrow$ fatigue diurne, hypertension artérielle secondaire.
    \item \textbf{Métaboliques :} Insulinorésistance, dyslipidémie, stéatose hépatique non alcoolique (NASH / « foie gras »), risque de cirrhose.
    \item \textbf{Psychosociales :} Stigmatisation, dépression, isolement, difficultés scolaires/professionnelles. Double fardeau dans nos sociétés où la rondeur est encore parfois culturellement valorisée comme signe de « bonne santé » ou de réussite sociale.
\end{itemize}
\end{exemplebox}

\begin{popfigure}
\begin{tikzpicture}[
    node distance=1.2cm and 1.5cm,
    box/.style={rectangle, draw=popdark, thick, fill=popcream, rounded corners=5pt, minimum width=3.5cm, minimum height=1.2cm, align=center, font=\small},
    arrow/.style={-Stealth, thick, popdark},
    plus/.style={font=\large\bfseries, poporange}
]
\node[box, align=center] (sed){Sédentarité\\ (Dépense $\downarrow$)};
\node[box, right=of sed, align=center] (exc){Excès alimentaire\\ (Apports $\uparrow$)};
\node[plus, below=0.2cm of sed.south east, xshift=1.8cm] {+};
\node[box, below=2cm of $(sed)!0.5!(exc)$, align=center] (obal){Bilan énergétique\\ positif chronique};
\node[box, below=1.8cm of obal, align=center] (obes){Obésité viscérale\\ (Adipokines $\uparrow$\\ Inflammation $\uparrow$)};
\node[box, below=1.8cm of obes, align=center] (ir){Insulinorésistance\\ (Muscle, Foie, Adipeux)};
\node[box, below=1.8cm of ir, align=center] (dt2){Diabète de type 2\\ (Hyperglycémie chronique)};

\draw[arrow] (sed) -- (obal);
\draw[arrow] (exc) -- (obal);
\draw[arrow] (obal) -- (obes);
\draw[arrow] (obes) -- (ir);
\draw[arrow] (ir) -- (dt2);

% Feedback loop
\draw[arrow, poporange] (ir.east) -- ++(1.5,0) |- (obes.east) node[midway, right, font=\tiny, poporange] {Aggravation};
\end{tikzpicture}
\legende{Chaîne causale : la sédentarité et l'excès alimentaire créent un bilan énergétique positif, menant à l'obésité viscérale, source d'inflammation et d'insulinorésistance, aboutissant au diabète de type 2. La flèche orange montre le cercle vicieux (hyperinsulinémie $\rightarrow$ prise de poids).}
\end{popfigure}

\courssub{Le diabète de type 2 : de l'insulinorésistance à l'hyperglycémie chronique}

\begin{definition}[Insulinorésistance et diabète de type 2 (DT2)]
L'\cle{insulinorésistance} est une diminution de la réponse des cellules cibles (muscle squelettique, foie, tissu adipeux) à l'\cle{insuline}. Le glucose n'entre plus efficacement dans les cellules (défaut de translocation du transporteur \cle{GLUT4}). Le pancréas compense en sécrétant \textbf{plus d'insuline} : \textbf{hyperinsulinémie compensatrice}. À terme, les \cle{cellules $\beta$ pancréatiques} s'épuisent (glucotoxicité, lipotoxicité) $\rightarrow$ insuffisance insulinique relative $\rightarrow$ \textbf{diabète de type 2 (DT2)} (hyperglycémie chronique : glycémie à jeun $\ge 1,26$ g/L ou HbA1c $\ge 6,5\%$).
\end{definition}

\begin{experience}[Effet de l'activité physique sur la glycémie post-prandiale -- Données type HGPO]
\textbf{Contexte :} Comparaison de la glycémie après charge orale de glucose (75 g) chez un sujet actif vs sédentaire (même IMC, même repas).
\textbf{Données (simulées d'après littérature) :}
\begin{center}
\begin{tabular}{|c|c|c|}\hline
Temps (min) & Glycémie Sujet Actif (g/L) & Glycémie Sujet Sédentaire (g/L) \\ \hline
0 (à jeun) & 0,85 & 0,95 \\ \hline
30 & 1,40 & 1,70 \\ \hline
60 & 1,60 & 2,10 \\ \hline
120 & 0,95 & 1,45 \\ \hline
\end{tabular}
\end{center}
\textbf{Interprétation :} Le pic est plus haut et le retour à la normale plus lent chez le sédentaire. Le muscle au repos capte mal le glucose (peu de GLUT4 à la membrane). La contraction musculaire active la translocation de GLUT4 \textbf{indépendamment de l'insuline} (via AMPK).
\textbf{Conclusion :} La sédentarité aggrave directement l'hyperglycémie post-prandiale et accélère l'épuisement des cellules $\beta$ pancréatiques.
\end{experience}

\begin{attention}[Piège majeur : « Le diabète de type 2, c'est pour les vieux ou les riches »]
\textbf{FAUX.}
\begin{itemize}
    \item \textbf{Âge :} Au Cameroun, on diagnostique désormais du DT2 chez des adolescents de 14--16 ans, obèses et sédentaires (consommation de boissons sucrées, « fast-food », écrans).
    \item \textbf{Milieu :} Ce n'est plus une maladie « des riches ». Les aliments ultra-transformés bon marché (nouilles instantanées, boissons sucrées, huiles raffinées chauffées) touchent les milieux modestes urbains.
    \item \textbf{Évitabilité :} Le DT2 est \textbf{largement évitable} par l'hygiène de vie (activité physique + alimentation équilibrée). Le diagnostic de \textbf{prédiabète} (glycémie à jeun 1,10--1,25 g/L ou HbA1c 5,7--6,4 \%) est une alerte \textbf{réversible}.
\end{itemize}
\end{attention}

\begin{exemplebox}[Symptômes et complications du DT2 non équilibré]
\begin{itemize}
    \item \textbf{Symptômes cardinaux (souvent discrets au début) :} \textbf{Polyuro-polydipsie} (urines abondantes + soif intense), fatigue majeure, amaigrissement paradoxal (perte de calories dans les urines), infections cutanées récidivantes (mycoses), cicatrisation lente.
    \item \textbf{Complications aiguës :} État hyperosmolaire (déshydratation sévère, coma -- plus fréquent en DT2), acidocétose (rare, plutôt DT1).
    \item \textbf{Complications chroniques (micro et macro-vasculaires) :}
    \begin{itemize}
        \item \textbf{Rétinopathie} $\rightarrow$ \textbf{cécité} (1$^{\text{re}}$ cause de cécité avant 60 ans).
        \item \textbf{Néphropathie} $\rightarrow$ \textbf{insuffisance rénale terminale} (dialyse, greffe ; coût énorme pour les familles et le système de santé camerounais).
        \item \textbf{Neuropathie} + artériopathie des membres inférieurs $\rightarrow$ \textbf{ulcères, gangrène, amputations} (pied diabétique).
        \item \textbf{Macro-vasculaire :} Infarctus du myocarde, AVC ischémique (risque $\times 2$ à $4$).
    \end{itemize}
\end{itemize}
\end{exemplebox}

\courssub{Maladies cardiovasculaires : l'athérosclérose et l'hypertension}

\begin{methode}[Genèse de la plaque d'athérome (athérosclérose)]
\begin{enumerate}
    \item \textbf{Agression de l'endothélium} : Hypertension artérielle (HTA), tabagisme, hyperglycémie, LDL-cholestérol oxydé $\rightarrow$ dysfonction endothéliale (perte du NO vasodilatateur, expression de molécules d'adhésion).
    \item \textbf{Infiltration lipidique} : LDL-cholestérol traverse l'intima $\rightarrow$ piégeage dans la matrice, oxydation (LDL-ox).
    \item \textbf{Réaction inflammatoire} : Monocytes recrutés $\rightarrow$ macrophages $\rightarrow$ phagocytose LDL-ox $\rightarrow$ \textbf{cellules spumeuses} (noyau lipidique central de la plaque).
    \item \textbf{Fibrose et calcification} : Prolifération cellules musculaires lisses (migration de la média vers l'intima), dépôt de collagène, calcification $\rightarrow$ \textbf{plaque fibreuse} (sténose stable) ou \textbf{plaque instable} (riche en lipides, capsule fibreuse fine, infiltrat inflammatoire dense).
    \item \textbf{Complication aiguë :} Rupture de plaque instable $\rightarrow$ exposition du tissu thrombogène $\rightarrow$ thrombose occlusive $\rightarrow$ \textbf{infarctus du myocarde} (coronaire) ou \textbf{AVC ischémique} (cérébrale).
\end{enumerate}
\end{methode}

\begin{popfigure}
\begin{tikzpicture}[scale=1.1, font=\small]
% Artère saine
\begin{scope}[xshift=0cm]
\draw[popblue, thick, fill=popblue!10] (0,0) ellipse (2.5 and 1.2);
\draw[popblue, thick, fill=popblue!30] (0,0) ellipse (2.0 and 0.7);
\draw[popred, thick, fill=popred!20] (0,0) ellipse (1.6 and 0.3);
\node at (0, -1.8) {\textbf{Artère saine}};
\draw[<-, popdark] (0, 1.2) -- (0, 1.7) node[above] {Lumière (sang)};
\draw[<-, popdark] (0, 0.7) -- (0, 1.1) node[above] {Intima (endothélium)};
\draw[<-, popdark] (0, 0.0) -- (0, 0.5) node[above] {Média (muscle lisse)};
\draw[<-, popdark] (0, -0.7) -- (0, -1.2) node[below] {Adventice (conjonctif)};
\draw[<-, popdark] (1.8, 0) -- (2.5, 0) node[right] {Paroi fine, élastique};
\end{scope}

% Artère athérosclérotique
\begin{scope}[xshift=7cm]
\draw[popblue, thick, fill=popblue!10] (0,0) ellipse (2.5 and 1.2);
\draw[popblue, thick, fill=popblue!30] (0,0) ellipse (2.0 and 0.7);
% Plaque d'athérome
\draw[poporange, thick, fill=poporange!30] (-1.2, -0.1) .. controls (-1.2, 0.4) and (0.5, 0.5) .. (0.5, -0.3) .. controls (0.3, -0.4) and (-1.0, -0.3) .. cycle;
\draw[popgold, thick, fill=popgold!50] (-0.8, 0.05) ellipse (0.3 and 0.15); % Nécrose lipidique
\draw[poppurple, thick, fill=poppurple!30] (-1.1, 0.2) .. controls (-0.8, 0.3) and (0.2, 0.3) .. (0.2, 0.1) .. controls (0.0, 0.15) and (-0.9, 0.15) .. cycle; % Cap fibreuse
\draw[popred, thick, fill=popred!20] (0,0) ellipse (1.6 and 0.3); % Média amincie
\node at (0, -1.8) {\textbf{Artère athérosclérotique}};
\draw[<-, poporange] (-0.5, 0.5) -- (-0.5, 1.0) node[above] {Plaque d'athérome};
\draw[<-, popgold] (-0.8, 0.3) -- (-0.8, 0.8) node[above] {Cœur lipidique (nécrose)};
\draw[<-, poppurple] (0.5, 0.4) -- (1.0, 0.8) node[right] {Cap fibreuse (risque rupture)};
\draw[<-, popred] (1.8, 0) -- (2.5, 0) node[right] {Média amincie, rigide};
\draw[<-, popdark] (0, -0.9) -- (0, -1.4) node[below] {Lumière rétrécie (sténose)};
\end{scope}
\end{tikzpicture}
\legende{Comparaison coupe transversale : artère saine (paroi fine, lumière large) vs artère athérosclérotique (plaque d'athérome avec cœur lipidique jaune, cap fibreuse violette, sténose de la lumière, rigidité pariétale).}
\end{popfigure}

\begin{savaistu}[L'hypertension artérielle : le « tueur silencieux » au Cameroun]
L'\cle{hypertension artérielle (HTA)} est définie par une pression artérielle $\ge 140/90$ mmHg (mesurée au cabinet, confirmée). Elle est appelée \textbf{« tueur silencieux »} car \textbf{elle ne donne souvent aucun symptôme} (pas de maux de tête systématiques, pas de vertiges) avant de provoquer un accident grave : AVC, infarctus, insuffisance cardiaque ou rénale.
\par\medskip
Au Cameroun, les enquêtes (STEPS 2003, études récentes) montrent une \textbf{prévalence supérieure à 30 \% chez l'adulte} (souvent 1 adulte sur 3 en zone urbaine). Seule une minorité est dépistée, traitée et contrôlée. L'excès de sel (bouillons cubes, poisson fumé, sel de table), l'obésité, la sédentarité, l'alcool et le stress en sont les moteurs principaux. \textbf{Se faire mesurer la tension artérielle au moins une fois par an est un geste vital.}
\end{savaistu}

\courssub{Le rôle aggravant de la sédentarité : au-delà du bilan calorique}

La sédentarité n'est pas seulement « ne pas bouger ». Elle a des effets physiologiques directs, indépendants de l'IMC (risque résiduel même chez les personnes de poids normal) :

\begin{propriete}[Effets métaboliques et vasculaires directs de la sédentarité (admis)]
\begin{itemize}
    \item \textbf{Baisse de la dépense énergétique de base (DEJ) :} Fonte musculaire (sarcopénie) $\rightarrow$ le muscle est le principal consommateur de glucose au repos. Moins de muscle = moins de « fournaise » métabolique.
    \item \textbf{Mauvaise utilisation du glucose (insulinorésistance musculaire) :} L'inactivité réduit l'expression et la translocation des transporteurs \cle{GLUT4} dans la fibre musculaire. Le glucose reste dans le sang.
    \item \textbf{Dyslipidémie :} Baisse du HDL-cholestérol (« bon » cholestérol, protecteur), hausse des triglycérides.
    \item \textbf{Fonction vasculaire :} Perte de l'effet « cisaillement » (shear stress) du flux sanguin sur l'endothélium $\rightarrow$ moins de NO produit $\rightarrow$ vasoconstriction, HTA, état pro-thrombotique.
    \item \textbf{Inflammation :} Le muscle au repos ne libère pas de \cle{myokines} bénéfiques (ex: IL-6 musculaire à effet anti-inflammatoire lors de l'exercice, irisine favorisant le « brunissement » du tissu adipeux).
\end{itemize}
\end{propriete}

\begin{popfigure}
\begin{tikzpicture}
\begin{axis}[
    ybar,
    bar width=25pt,
    width=\linewidth,
    height=8cm,
    ylabel={Dépense énergétique journalière (kcal)},
    ymin=0, ymax=3200,
    xtick=data,
    xticklabels={Adolescent Actif, Adolescent Sédentaire},
    ymajorgrids=true,
    grid style=dashed,
    legend style={at={(0.5,-0.18)}, anchor=north, legend columns=-1, font=\small},
    nodes near coords,
    nodes near coords align={vertical},
    every node near coord/.append style={font=\small, popdark}
]
\addplot[popblue, fill=popblue!60] coordinates {(1, 2780)};
\addplot[poporange, fill=poporange!60] coordinates {(2, 2230)};
\legend{Actif : MB + Activité (NAP 1,75) $\approx$ 2780 kcal, Sédentaire : MB seul (NAP 1,4) $\approx$ 2230 kcal}
\end{axis}
\end{tikzpicture}
\legende{Comparaison de la dépense énergétique journalière estimée pour un adolescent de 16 ans (60 kg, 1,70 m, MB $\approx$ 1590 kcal). Actif : NAP 1,75 (marche, sport, travaux) $\approx$ 2780 kcal. Sédentaire : NAP 1,4 (école, écrans, transport motorisé) $\approx$ 2230 kcal. \textbf{Écart de 550 kcal/jour} : l'équivalent d'un gros repas ou 1,5 L de soda qu'il « faut » ne pas manger pour ne pas grossir en étant sédentaire.}
\end{popfigure}

\courssub{Facteurs de risque associés : la synergie du danger}

L'obésité, le diabète, l'HTA et la dyslipidémie ne voyagent pas seuls. Ils forment le \textbf{syndrome métabolique} (diagnostic IDF : tour de taille $\uparrow$ \textbf{obligatoire} + 2 critères sur 4 : TG $\uparrow$, HDL $\downarrow$, PA $\uparrow$, glycémie $\uparrow$). S'y ajoutent des facteurs comportementaux qui \textbf{multiplient} les risques (synergie, pas simple addition) :

\begin{center}
\begin{tabularx}{\linewidth}{|l|X|X|}\hline
\rowcolor{popblueL}
\textbf{Facteur} & \textbf{Mécanisme aggravant} & \textbf{Contexte camerounais} \\ \hline
\textbf{Tabac} & Endommage endothélium, oxyde LDL, vasoconstriction, pro-thrombotique. Risque CV $\times 2$--$3$. & Consommation en hausse chez les jeunes (chicha, cigarettes manufacturées). Loi anti-tabac 2006 à faire appliquer. \\ \hline
\textbf{Alcool} & Hypertension (dose-dépendant), cardiomyopathie, hypertriglycéridémie, cancers, accidents. & Bière locale, vins de palme, spiritueux forts. \emph{Binge-drinking} week-end fréquent chez étudiants. \\ \hline
\textbf{Stress chronique} & Cortisol $\uparrow$ $\rightarrow$ hyperglycémie, prise de poids abdominale, HTA, troubles du sommeil $\rightarrow$ cercle vicieux. & Pression scolaire, chômage, insécurité, embouteillages (Douala/Yaoundé), charge mentale. \\ \hline
\textbf{Sommeil insuffisant} ($< 7$ h) & Leptine $\downarrow$ (satiété), Ghréline $\uparrow$ (faim), cortisol $\uparrow$, insulinorésistance. & Écrans tardifs, horaires décalés, bruit urbain, promiscuité. \\ \hline
\end{tabularx}
\end{center}

\begin{aretenir}[Maladies Non Transmissibles (MNT) : l'urgence africaine]
Les \textbf{MNT} (maladies cardiovasculaires, cancers, maladies respiratoires chroniques, diabète) sont la \textbf{1$^{\text{re}}$ cause de mortalité mondiale} (74 \% des décès, OMS 2024).
\par\medskip
En Afrique subsaharienne, longtemps épargnée, la transition épidémiologique est brutale : \textbf{les MNT devraient devenir la 1$^{\text{re}}$ cause de décès vers 2030}. Au Cameroun, elles représentent déjà plus de 35 \% des décès. L'urbanisation rapide (alimentation transformée, transport motorisé, travail sédentaire, pollution, stress) crée un \textbf{environnement obésogène et diabétogène}.
\par\medskip
\textbf{La prévention primaire (mode de vie) est la seule stratégie soutenable} pour nos systèmes de santé fragiles. C'est l'objet des deux prochaines sections.
\end{aretenir}

\begin{exoresolu}[Analyse d'une ration « piège »]
\textbf{Énoncé :} Paul, 17 ans, 75 kg, 1,72 m (IMC 25,3 = surpoids), élève en 1ère D. Il ne fait pas de sport. Sa journée type : Petit-déj : 2 tartines beurre/confiture + 1 bol café au lait sucré (2 sucres). 10h : 1 bouteille « Top » (jus sucré 30 cl) + 1 pain chocolat. Déjeuner : 400 g riz + 200 g sauce arachide (huile rouge) + 100 g poisson frit + 2 plantains mûrs frits. 16h : 1 bouteille Coca-Cola 65 cl + 1 paquet chips. Dîner : 300 g pâtes + 2 œufs durs + sauce tomate industrielle. Coucher : 23h. Écran 4h/jour. Marche 15 min/jour (arrêt bus).
\begin{enumerate}
    \item Estimer ses apports caloriques approximatifs (utiliser : riz cuit 130 kcal/100g, huile 900 kcal/100g, soda 43 kcal/10cl, pain chocolat 300 kcal, plantain frit 200 kcal/100g, pâtes 130 kcal/100g).
    \item Calculer son BEJ (Métabolisme de base $\approx 1\,750$ kcal ; NAP sédentaire 1,4).
    \item Conclure sur l'évolution prévisible de son poids et sa santé à 5 ans.
\end{enumerate}
\textbf{Solution détaillée :}
\begin{enumerate}
    \item \textbf{Apports :} Pdej $\approx 450$ ; 10h $\approx 130 + 300 = 430$ ; Dej : Riz 520 + Sauce (huile 40g $= 360$ + arachide) $\approx 500$ + Poisson frit (huile) 250 + Plantains (200g $= 400$) $= 1\,670$ ; 16h : Soda 280 + Chips 270 $= 550$ ; Dîner : Pâtes 390 + Œufs 160 + Sauce 100 $= 650$. \textbf{Total $\approx 3\,750$ kcal}.
    \item \textbf{BEJ :} MB $1\,750 \times 1,4 = \textbf{2 450 kcal}$.
    \item \textbf{Bilan :} Excédent \textbf{$+1\,300$ kcal/jour}. Stockage potentiel $\approx 1\,300 \times 365 / 7\,700 \approx \textbf{+62 kg de graisse/an}$ (théorique, limité par l'augmentation de la DEJ avec le poids, mais prise massive réelle $\approx 10$--$15$ kg/an). \textbf{Risque imminent :} Obésité morbide, Prédiabète/DT2, HTA, Dyslipidémie, Stéatose hépatique. \textbf{Action urgente :} Réduire boissons sucrées, fritures, portions glucidiques ; augmenter marche (30 min/j) ; consultation médicale.
\end{enumerate}
\end{exoresolu}

\begin{exercice}[Mise en situation : Le diagnostic de risque]
Marie, 16 ans, IMC 28,5. Tour de taille 92 cm. Tension 135/85 mmHg. Glycémie à jeun 1,12 g/L. Triglycérides 1,80 g/L. HDL 0,40 g/L. Elle ne fait aucune activité physique, mange souvent des beignets-haricots le midi, boit 2 sodas/jour. Son père est diabétique type 2.
\begin{enumerate}
    \item Identifier les critères du syndrome métabolique présents chez Marie.
    \item Expliquer le lien physiopathologique entre son tour de taille et sa glycémie à jeun.
    \item Proposer 3 mesures concrètes, réalistes et prioritaires pour inverser la tendance (alimentation, activité, suivi médical).
\end{enumerate}
\end{exercice}

\begin{corrige}[Exercice n°1]
\begin{enumerate}
    \item \textbf{Critères du syndrome métabolique (IDF Afrique subsaharienne) :}
    \begin{enumerate}
        \item Tour de taille $\ge 80$ cm (femme) $\rightarrow$ \textbf{OUI (92 cm)}. \emph{Critère obligatoire rempli.}
        \item Triglycérides $\ge 1,50$ g/L $\rightarrow$ \textbf{OUI (1,80)}.
        \item HDL-cholestérol $< 0,50$ g/L (femme) $\rightarrow$ \textbf{OUI (0,40)}.
        \item Pression artérielle $\ge 130/85$ mmHg $\rightarrow$ \textbf{OUI (135/85)}.
        \item Glycémie à jeun $\ge 1,00$ g/L (ou DT2 traité) $\rightarrow$ \textbf{OUI (1,12 g/L = prédiabète)}.
    \end{enumerate}
    \textbf{5/5 critères : Syndrome métabolique avéré. Risque cardiovasculaire très élevé.}
    \item \textbf{Lien tour de taille / glycémie :} Le tour de taille reflète l'adiposité \textbf{viscérale} (graisse intra-abdominale, péri-viscérale). Ce tissu est hautement lipolytique (insulinorésistant aux effets antilipolytiques de l'insuline) et libère massivement des \textbf{acides gras libres (AGL)} dans la veine porte $\rightarrow$ foie. Les AGL hépatiques : (a) stimulent la néoglucogenèse $\rightarrow$ hyperglycémie à jeun ; (b) favorisent la synthèse et la sécrétion de VLDL (triglycérides $\uparrow$) ; (c) induisent une \textbf{insulinorésistance hépatique} (le foie ne répond plus à l'insuline qui devrait inhiber la production de glucose). D'où la glycémie à jeun élevée et l'hypertriglycéridémie.
    \item \textbf{Mesures prioritaires (réalistes, progressives) :}
    \begin{enumerate}
        \item \textbf{Supprimer les boissons sucrées} (2 sodas $= 560$ kcal vides, $110$ g sucre/j) $\rightarrow$ remplacer par eau / jus de fruits maison non sucrés / bissap/gingembre peu sucré. Gain immédiat $\approx -500$ kcal/j, baisse brutale de la charge glycémique.
        \item \textbf{Marcher 30 min continu / jour} (aller-retour école à pied, marche active le soir en famille) $\rightarrow$ augmente dépense $\approx +150$--$200$ kcal/j, améliore insulinorésistance musculaire (translocation GLUT4), baisse PA, améliore profil lipidique (HDL $\uparrow$).
        \item \textbf{Consultation médicale} pour bilan complet (HbA1c, fonction rénale -- créatininémie, microalbuminurie, ECG, bilan lipidique complet) et prise en charge du prédiabète / HTA limite / dyslipidémie. Suivi à 3 mois (objectif : perte de 5--10 \% du poids initial, normalisation glycémie à jeun $< 1,00$ g/L).
    \end{enumerate}
\end{enumerate}
\end{corrige}

\coursec{Recommandations : alimentation équilibrée}

Nous venons de voir que la mauvaise alimentation et la sédentarité conduisent à l'obésité, au diabète de type 2 et aux maladies cardiovasculaires. La question pratique se pose alors naturellement : \emph{que faut-il manger, et comment, pour rester en bonne santé ?} Cette section répond à cette question en traduisant les recommandations internationales (OMS) en gestes concrets, adaptés à notre alimentation camerounaise quotidienne.

\courssub{Les principes généraux d'une alimentation équilibrée}

\begin{definition}[Alimentation équilibrée]
Une alimentation est dite \cle{équilibrée} lorsqu'elle couvre, sans excès ni carence, les besoins de l'organisme en énergie (glucides, lipides, protéines), en eau, en vitamines, en sels minéraux et en fibres, en quantités adaptées à l'âge, au sexe et à l'activité physique de l'individu.
\end{definition}

Quatre grands principes guident la construction de cette alimentation :

\begin{itemize}
\item \textbf{La variété} : aucun aliment seul ne contient tous les nutriments nécessaires. Le riz apporte surtout des glucides, le poisson des protéines, l'huile rouge de la vitamine A, les légumes des fibres et des vitamines. Il faut donc \cle{varier} les aliments au cours de la journée et de la semaine, en puisant dans tous les groupes alimentaires.
\item \textbf{La modération} : même un bon aliment devient nocif en excès. Un plat d'eru très gras ou une grande quantité de fufu consommée chaque jour finit par déséquilibrer la ration énergétique. C'est la \emph{quantité} qui fait souvent la différence entre un repas sain et un repas nocif.
\item \textbf{La régularité des repas} : l'organisme fonctionne mieux avec trois repas répartis dans la journée (petit-déjeuner, déjeuner, dîner), éventuellement complétés d'une collation. Sauter un repas provoque une hypoglycémie, une fatigue, puis des compulsions alimentaires.
\item \textbf{L'hydratation suffisante} : l'eau est indispensable à toutes les réactions du métabolisme, à l'élimination des déchets par les reins et à la régulation de la température corporelle. Il faut boire \cle{\SI{1,5}{\liter} à \SI{2}{\liter} d'eau par jour}, davantage en cas de forte chaleur ou d'effort physique — une réalité quotidienne sous notre climat.
\end{itemize}

\begin{attention}[L'eau, pas les boissons sucrées]
L'hydratation doit se faire essentiellement avec de l'\textbf{eau potable}. Les sodas et jus industriels, même bus en grande quantité, n'hydratent pas « gratuitement » : ils apportent des sucres rapides qui favorisent la prise de poids et le diabète. Un verre de soda de \SI{33}{\centi\liter} contient environ 7 morceaux de sucre !
\end{attention}

\courssub{L'assiette idéale : une règle simple à retenir}

Pour éviter les calculs compliqués à chaque repas, les nutritionnistes proposent un modèle visuel très simple : l'\cle{assiette idéale}. À chaque repas principal, l'assiette doit être composée de :

\begin{itemize}
\item \textbf{une moitié de légumes} (crus ou cuits) : ils apportent fibres, vitamines, minéraux, et rassasient avec peu de calories ;
\item \textbf{un quart de féculents} (riz, maïs, igname, plantain, patate, miondo...) : source principale d'énergie ;
\item \textbf{un quart de protéines} (poisson, viande maigre, œufs, niébé, haricots, arachides) : construction et renouvellement des tissus ;
\item complétée par \textbf{un fruit} et \textbf{un verre d'eau}.
\end{itemize}

\begin{popfigure}
\begin{center}
\begin{tikzpicture}
% Assiette
\draw[thick, popdark, fill=popcream] (0,0) circle (2.6);
\draw[thick, popdark] (0,0) circle (2.2);
% Moitié légumes (gauche)
\fill[popgreenL] (90:2.2) arc (90:270:2.2) -- cycle;
\draw[thick, popdark] (0,2.2) -- (0,-2.2);
% Quart féculents (haut droit)
\fill[popgoldL] (90:2.2) arc (90:0:2.2) -- cycle;
% Quart protéines (bas droit)
\fill[poporangeL] (0:2.2) arc (0:-90:2.2) -- cycle;
\draw[thick, popdark] (0,0) -- (2.2,0);
% Étiquettes avec lignes de rappel
\node at (-1.1,0) {\small \textbf{1/2 légumes}};
\node at (1.1,1.1) {\small \textbf{1/4 féculents}};
\node at (1.1,-1.1) {\small \textbf{1/4 protéines}};
% Fruit
\draw[thick, popdark, fill=poppinkL] (4.1,0.9) circle (0.55);
\node at (4.1,0.9) {\small fruit};
% Verre d'eau
\draw[thick, popdark] (3.75,-0.5) -- (3.85,-1.6) -- (4.45,-1.6) -- (4.55,-0.5) -- cycle;
\fill[popblueL] (3.79,-0.85) -- (3.87,-1.55) -- (4.43,-1.55) -- (4.51,-0.85) -- cycle;
\node at (4.15,-2.0) {\small verre d'eau};
\end{tikzpicture}
\end{center}
\legende{Schéma de l'assiette idéale : la moitié de l'assiette est occupée par les légumes, un quart par les féculents et un quart par les protéines. Le repas est complété par un fruit et un verre d'eau.}
\end{popfigure}

\begin{propriete}[Recommandations de l'OMS]
L'Organisation mondiale de la santé recommande, pour un adulte :
\begin{itemize}
\item au moins \cle{\SI{400}{\gram} de fruits et légumes par jour} (soit environ 5 portions) ;
\item moins de \cle{\SI{5}{\gram} de sel par jour} (environ une cuillère à café rase, tout sel confondu, y compris celui des cubes d'assaisonnement) ;
\item des \cle{sucres libres} limités à moins de \cle{10 \% de l'apport énergétique total} (idéalement moins de 5 \%), soit moins de \SI{50}{\gram} par jour pour une ration de \SI{2000}{kcal} ;
\item des lipides représentant moins de 30 \% de l'apport énergétique, en limitant les graisses saturées.
\end{itemize}
\end{propriete}

\courssub{Que faut-il augmenter ?}

\textbf{Les fruits et légumes : au moins 5 portions par jour.} Une portion correspond à environ \SI{80}{\gram}–\SI{100}{\gram} : une tomate moyenne, une banane, une poignée de feuilles de ndolé cuites, une goyave. Les fruits et légumes apportent des fibres (qui régulent le transit et ralentissent l'absorption des sucres), des vitamines (A, C, folates) et des minéraux (potassium, protecteur contre l'hypertension). Leur richesse en eau et en fibres permet de se rassasier avec peu de calories : c'est l'allié majeur contre la prise de poids.

\textbf{Les légumineuses} : niébé, haricots, arachides, lentilles, soja. Excellentes sources de protéines végétales peu coûteuses, elles contiennent aussi des fibres et du fer. Associées à une céréale (riz + niébé, maïs + haricots), elles fournissent des protéines de qualité comparable à celle de la viande, pour un coût bien moindre — un atout considérable pour les familles camerounaises.

\textbf{Les céréales complètes et peu raffinées} : maïs en grains, mil, sorgho, riz complet, igname, plantain. Contrairement aux produits très raffinés (pain blanc, pâtes blanches), ils conservent leurs fibres et leurs vitamines, et leur digestion plus lente évite les pics de glycémie.

\textbf{Le poisson} : source de protéines de haute qualité et de bons lipides (acides gras oméga-3 protecteurs du cœur), il devrait figurer au menu au moins deux fois par semaine. Le poisson braisé ou bouilli est préférable au poisson frit dans beaucoup d'huile.

\courssub{Que faut-il réduire ?}

\begin{itemize}
\item \textbf{Les sucres rapides et les boissons sucrées} : sodas, jus industriels, bonbons, beignets sucrés en excès. Ils provoquent des pics de glycémie suivis de fringales, favorisent la carie dentaire, l'obésité et le diabète de type 2.
\item \textbf{Les graisses de friture et l'excès d'huile} : la friture concentre les calories et, si l'huile est réutilisée plusieurs fois, produit des composés toxiques. Préférer la cuisson à la vapeur, au four, braisée ou en sauce modérément huilée.
\item \textbf{Le sel : moins de \SI{5}{\gram} par jour.} L'excès de sel élève la tension artérielle et multiplie le risque d'accident vasculaire cérébral. Attention : le sel ne vient pas seulement de la salière !
\item \textbf{L'alcool} : il apporte des calories « vides » (\SI{7}{kcal.g^{-1}}), sans aucun nutriment utile, et expose à la dépendance, aux maladies du foie (cirrhose) et aux cancers. Pour un adolescent, la recommandation est l'\textbf{abstinence totale}.
\end{itemize}

\begin{popfigure}
\begin{center}
\begin{tikzpicture}
% Barres comparatives : morceaux de sucre par verre de 33 cl
\draw[->, thick, popdark] (0,0) -- (0,4.6) node[above, align=center, font=\small] {morceaux de sucre\\(par verre de \SI{33}{\centi\liter})};
\draw[->, thick, popdark] (0,0) -- (8.6,0);
% Eau
\fill[popblue] (0.6,0) rectangle (2.0,0.05);
\node[below, font=\small, align=center] at (1.3,-0.05) {Eau};
\node[above, font=\small] at (1.3,0.1) {0};
% Jus naturel sans sucre ajouté
\fill[popgreen] (3.4,0) rectangle (4.8,1.6);
\node[below, font=\small, align=center] at (4.1,-0.05) {Jus naturel\\sans sucre ajouté};
\node[above, font=\small] at (4.1,1.65) {$\approx$ 3};
% Soda
\fill[poporange] (6.2,0) rectangle (7.6,3.8);
\node[below, font=\small, align=center] at (6.9,-0.05) {Soda};
\node[above, font=\small] at (6.9,3.85) {$\approx$ 7};
% Graduations
\foreach \y/\lab in {0.55/1, 1.1/2, 1.65/3, 2.2/4, 2.75/5, 3.3/6, 3.85/7}
  \draw[popmuted] (-0.08,\y) -- (0.08,\y) node[left=4pt, font=\scriptsize] at (-0.1,\y) {\lab};
\end{tikzpicture}
\end{center}
\legende{Comparaison de la teneur en sucre de trois boissons (en morceaux de sucre de \SI{5}{\gram} pour un verre de \SI{33}{\centi\liter}). L'eau est la seule boisson sans sucre ; un soda apporte environ \SI{35}{\gram} de sucres rapides, soit près de la limite quotidienne recommandée par l'OMS à lui seul.}
\end{popfigure}

\begin{savaistu}[Le sel caché : l'ennemi invisible]
Une grande partie du sel que nous consommons ne vient pas de la salière : c'est le \cle{sel caché}. Les cubes d'assaisonnement (très riches en sel et en glutamate), le poisson fumé, la viande fumée, les snacks salés et les plats préparés industriels représentent souvent \textbf{plus de la moitié} du sel total consommé. Or l'hypertension artérielle touche une part croissante des adultes camerounais. Réduire d'un cube par jour dans la marmite familiale est un levier simple, gratuit et puissant de prévention cardiovasculaire.
\end{savaistu}

\courssub{Valoriser l'alimentation locale}

Manger équilibré ne signifie pas manger « à l'européenne » : notre cuisine traditionnelle camerounaise offre d'excellents plats complets, à condition de respecter des portions raisonnables et de limiter l'huile et le sel.

\begin{itemize}
\item \textbf{Ndolé + miondo} : les feuilles de ndolé apportent fibres, fer et vitamines ; les arachides et les crevettes apportent protéines et lipides ; le miondo (bâton de manioc) fournit les glucides. C'est un plat presque complet — à condition de ne pas multiplier les bâtons de manioc et de modérer l'huile.
\item \textbf{Eru + water fufu} : l'eru est riche en légumes et en poisson/viande, mais souvent très gras (huile rouge en abondance). Une portion raisonnable de water fufu et une sauce moins huilée en font un bon repas.
\item \textbf{Okok} : plat de feuilles riche en fibres et minéraux, à accompagner de bâton de manioc ou de plantain en quantité modérée.
\item \textbf{Koki} (gâteau de haricot ou de niébé cuit à la vapeur) : excellente source de protéines végétales, cuit sans friture — un modèle de cuisson saine.
\end{itemize}

\begin{exemplebox}[Transformer un repas ordinaire en repas équilibré]
Repas fréquent : une grande montagne de riz blanc avec un peu de sauce grasse. \textbf{Version équilibrée} : réduire le riz à un quart de l'assiette, ajouter une moitié de légumes (sauce feuilles, salade de tomates et concombres), un quart de protéines (un morceau de poisson braisé ou une portion de haricots), puis une banane ou une tranche d'ananas en dessert, et un verre d'eau. Même budget, meilleure santé.
\end{exemplebox}

\courssub{Lire les étiquettes et se méfier des produits ultra-transformés}

Les \cle{produits ultra-transformés} (chips, biscuits industriels, nouilles instantanées, sodas, charcuteries, céréales sucrées du petit-déjeuner) sont fabriqués industriellement à partir d'ingrédients raffinés, avec de nombreux additifs, et sont presque toujours trop sucrés, trop salés et/ou trop gras. Leur consommation régulière est associée à l'obésité et au diabète.

\begin{methode}[Lire une étiquette alimentaire]
\begin{enumerate}
\item Repérer la \textbf{liste des ingrédients} : ils sont classés par ordre de quantité décroissante. Si le sucre (ou le sirop de glucose) figure en première ou deuxième position, le produit est très sucré.
\item Consulter le \textbf{tableau nutritionnel} : comparer les valeurs « pour \SI{100}{\gram} » de sucres, de sel et de lipides entre produits.
\item Vérifier la \textbf{portion réelle} : un paquet peut contenir 2 ou 3 portions alors que les valeurs affichées ne concernent qu'une portion.
\item Se méfier du vocabulaire marketing : « allégé », « source de fibres » ou « vitamines ajoutées » ne garantissent pas un produit sain.
\item Règle simple : plus la liste d'ingrédients est longue et incompréhensible, plus le produit est transformé — préférer les aliments simples et locaux.
\end{enumerate}
\end{methode}

\courssub{Cas particuliers}

\textbf{L'adolescent sportif.} L'effort physique augmente les besoins en énergie et en eau. Il faut renforcer les féculents complets (source de glucides, carburant des muscles), maintenir de bonnes protéines (poisson, œufs, légumineuses) pour la réparation musculaire, et surtout boire avant, pendant et après l'effort. Inutile d'acheter des boissons « énergisantes » sucrées : eau + banane suffisent largement.

\textbf{La jeune fille et le fer.} Les règles entraînent chaque mois des pertes de fer. Or la carence en fer provoque l'\cle{anémie} (fatigue, pâleur, difficultés de concentration, baisse des résultats scolaires). La jeune fille doit consommer régulièrement des aliments riches en fer : feuilles vertes (ndolé, okok, folong), légumineuses, viande, poisson, et les associer à un fruit riche en vitamine C (goyave, orange, citron) qui multiplie l'absorption du fer végétal. À l'inverse, le thé pris pendant le repas freine cette absorption.

\textbf{Jeûne et fêtes : gérer les excès ponctuels.} Lors des fêtes (Noël, Tabaski, mariages) ou des périodes de jeûne (Ramadan, Carême), le rythme alimentaire est bouleversé. Quelques règles de bon sens : ne pas se jeter sur des plats très gras et très sucrés à la rupture du jeûne (commencer par de l'eau, une soupe ou des fruits) ; limiter les sodas qui accompagnent les festins ; boire beaucoup d'eau entre les repas ; et, après un excès, revenir simplement à des repas légers et équilibrés les jours suivants — un excès ponctuel ne ruine pas la santé, c'est la répétition qui est dangereuse.

\courssub{Construire un menu équilibré d'une journée}

\begin{methode}[Construire un menu équilibré d'une journée]
\begin{enumerate}
\item \textbf{Répartir les calories} sur 3 repas (+ éventuellement une collation) : environ \SI{25}{\%} au petit-déjeuner, \SI{35}{\%}–\SI{40}{\%} au déjeuner, \SI{25}{\%}–\SI{30}{\%} au dîner, \SI{5}{\%}–\SI{10}{\%} en collation.
\item \textbf{Inclure les 3 macronutriments à chaque repas} : glucides (féculents), lipides (en quantité modérée), protéines (animales ou végétales).
\item \textbf{Garantir fruits, légumes et eau} : au moins 5 portions de fruits et légumes réparties dans la journée et \SI{1,5}{\liter} à \SI{2}{\liter} d'eau.
\item \textbf{Limiter sel, sucre et fritures} : moins de \SI{5}{\gram} de sel, pas de soda, une seule cuisson à l'huile modérée au maximum.
\end{enumerate}
\end{methode}

\begin{exemplebox}[Menu type à \SI{2500}{kcal} (adolescent actif)]
\begin{center}
\begin{tabularx}{\linewidth}{|l|X|}
\hline
\rowcolor{popblueL} \textbf{Repas} & \textbf{Composition} \\
\hline
Petit-déjeuner & Bouillie de mil enrichie (lait ou arachide) + une banane \\
\hline
Déjeuner & Riz + ndolé au poisson + salade de tomates et concombres + eau \\
\hline
Collation & Une goyave + une poignée d'arachides grillées non salées \\
\hline
Dîner & Haricots + banane plantain vapeur + un fruit + eau \\
\hline
\end{tabularx}
\end{center}
Ce menu couvre les besoins énergétiques d'un adolescent actif, fournit des protéines à chaque repas principal (poisson, arachides, haricots), plus de 5 portions de fruits et légumes, des fibres, du fer, et reste pauvre en sucres rapides et en fritures.
\end{exemplebox}

\begin{attention}[Erreur fréquente : sauter le petit-déjeuner]
Beaucoup d'élèves partent au lycée sans petit-déjeuner, « pour gagner du temps ». Résultat : hypoglycémie en milieu de matinée (fatigue, maux de tête, difficultés de concentration), puis compensation par des grignotages sucrés et gras (beignets, bonbons, biscuits, sodas) achetés à la cantine ou chez les vendeuses ambulantes. Au total, la ration de la journée est \emph{plus} calorique et \emph{moins} équilibrée que si l'élève avait pris un vrai petit-déjeuner. Sauter un repas ne fait jamais maigrir sainement : cela déséquilibre la ration totale.
\end{attention}

\begin{exoresolu}[Analyser et corriger une journée alimentaire]
\textbf{Énoncé.} Rodrigue, 16 ans, prend le matin un beignet et un soda ; à midi, un grand plat de riz blanc avec une sauce très huilée et deux cubes d'assaisonnement ; le soir, des pâtes sautées à l'huile avec des œufs frits ; il boit environ \SI{0,5}{\liter} d'eau par jour et ne mange ni fruit ni légume. Identifier quatre déséquilibres de cette ration et proposer des corrections concrètes.
\tcblower
\textbf{Solution détaillée.}
\begin{enumerate}
\item \textbf{Excès de sucres rapides} (soda dès le matin) : remplacer le soda par de l'eau et le beignet par une bouillie de mil ou du pain complet avec une banane.
\item \textbf{Absence totale de fruits et légumes} (0 portion au lieu de 5) : ajouter une salade ou une sauce feuilles au déjeuner, un fruit au goûter et des légumes au dîner.
\item \textbf{Excès de sel et de graisses de friture} (deux cubes, sauces très huilées, œufs frits) : limiter à un cube, réduire l'huile, préférer les œufs bouillis et les cuissons vapeur ou braisées.
\item \textbf{Hydratation insuffisante} (\SI{0,5}{\liter} au lieu de \SI{1,5}{\liter} à \SI{2}{\liter}) : boire régulièrement de l'eau dans la journée, surtout sous notre climat chaud.
\end{enumerate}
\textbf{Conclusion} : la ration de Rodrigue est trop riche en sucres rapides, en sel et en graisses, carencée en fibres, vitamines et eau. Les corrections proposées ne coûtent pas plus cher : elles reposent sur des aliments locaux simples (bouillie de mil, fruits de saison, légumes feuilles, eau potable).
\end{exoresolu}

\begin{aretenir}[L'essentiel de la section]
\begin{itemize}
\item Une alimentation équilibrée repose sur la \textbf{variété}, la \textbf{modération}, la \textbf{régularité des repas} et une \textbf{hydratation suffisante} (\SI{1,5}{\liter} à \SI{2}{\liter} d'eau par jour).
\item L'\textbf{assiette idéale} : 1/2 légumes, 1/4 féculents, 1/4 protéines, + un fruit et un verre d'eau.
\item Recommandations OMS : au moins \SI{400}{\gram} de fruits et légumes par jour, moins de \SI{5}{\gram} de sel par jour, sucres libres $<$ 10 \% de l'apport énergétique.
\item À augmenter : fruits, légumes, légumineuses, céréales complètes, poisson. À réduire : sucres rapides, sodas, fritures, sel (y compris le sel caché des cubes), alcool.
\item Nos plats traditionnels (ndolé, okok, koki, eru) sont d'excellentes bases, à condition de maîtriser portions, huile et sel.
\item Savoir lire les étiquettes permet de repérer les produits ultra-transformés, trop sucrés, salés et gras.
\end{itemize}
\end{aretenir}

\coursec{Recommandations : activité physique}

Dans la section précédente, nous avons appris à construire une alimentation équilibrée : des portions adaptées, des groupes d'aliments variés, une limitation du sucre, du sel et des graisses. Mais l'équilibre alimentaire ne représente qu'une moitié de l'équation énergétique : il contrôle les \cle{entrées} d'énergie. L'autre moitié, ce sont les \cle{sorties} : la dépense d'énergie. Or, dans nos villes camerounaises, de nombreux adolescents passent leurs journées assis — en classe, dans le taxi-brousse, devant le téléphone ou la télévision. À Yaoundé comme à Douala, les médecins observent une hausse inquiétante du surpoids chez les jeunes, alors même que la nourriture n'a pas fondamentalement changé. Ce qui a changé, c'est le \emph{mouvement}. Cette section te donne les recommandations officielles de l'Organisation mondiale de la Santé (OMS) en matière d'activité physique, et surtout les moyens concrets de les appliquer dans ta vie quotidienne.

\courssub{Activité physique, exercice, sport : trois notions à distinguer}

\textbf{Observation de départ.} Considère trois personnes : une maman qui balaye sa cour et lave le linge à la main pendant une heure ; un lycéen qui fait ses étirements et ses abdominaux chaque matin ; une joueuse de football de l'équipe de quartier qui dispute un match. Les trois bougent, les trois dépensent de l'énergie. Pourtant, on dira spontanément que seule la troisième « fait du sport ». Est-ce que, du point de vue de la santé, seule la joueuse de football bénéficie de son activité ? Nous verrons que non — et c'est une excellente nouvelle.

\begin{definition}[Activité physique]
L'\cle{activité physique} désigne \textbf{tout mouvement corporel produit par la contraction des muscles squelettiques qui entraîne une dépense d'énergie supérieure à celle du repos}. Elle englobe les activités de la vie quotidienne : marcher, monter les escaliers, porter l'eau, jardiner, faire le ménage, danser, jouer.
\end{definition}

\begin{definition}[Exercice et sport]
L'\cle{exercice} est une sous-catégorie de l'activité physique : il s'agit de mouvements \textbf{planifiés, structurés et répétés}, réalisés dans le but d'améliorer ou de maintenir la forme physique (séance d'abdominaux, footing du matin). Le \cle{sport} est un exercice \textbf{codifié par des règles}, souvent pratiqué en compétition (football, handball, athlétisme). On a donc la relation d'inclusion suivante : tout sport est un exercice, tout exercice est une activité physique, mais toute activité physique n'est ni un exercice ni un sport.
\end{definition}

\begin{attention}[Erreur fréquente : « seul le sport en tenue compte »]
Beaucoup d'élèves croient que seule une séance de sport « en tenue », sur un terrain, apporte des bénéfices pour la santé. C'est faux. Du point de vue physiologique, ce qui compte, c'est la \cle{dépense d'énergie} et la sollicitation du cœur, des muscles et des vaisseaux — peu importe le cadre. Les \textbf{activités utilitaires} (marcher jusqu'à l'école, faire le ménage, jardiner au champ, porter l'eau, danser) sont \textbf{pleinement bénéfiques} et comptent dans les recommandations de l'OMS. Au Probatoire, ne réduis jamais l'activité physique au sport de compétition.
\end{attention}

\courssub{Les intensités de l'activité physique : modérée ou soutenue ?}

\textbf{Observation.} Marche tranquillement dans la cour : tu peux discuter et même chanter sans difficulté. Maintenant cours pour attraper le car qui démarre : au bout d'une minute, tu es essoufflé, tu ne peux plus prononcer une phrase entière sans reprendre ton souffle. Ton corps vient de passer d'une intensité \emph{modérée} à une intensité \emph{soutenue}. Cette différence se mesure objectivement (fréquence cardiaque, consommation d'oxygène), mais elle s'évalue très simplement avec le \textbf{test de la parole}.

\begin{definition}[Activité physique d'intensité modérée]
Une activité d'\cle{intensité modérée} est une activité qui \textbf{accélère nettement le rythme cardiaque et la respiration tout en permettant encore de parler} (mais pas de chanter). Exemples : marche rapide (environ 5 à \SI{6}{km/h}), vélo tranquille, ménage actif, danse modérée, jardinage.
\end{definition}

\begin{definition}[Activité physique d'intensité soutenue]
Une activité d'\cle{intensité soutenue} (ou vigoureuse) provoque un \textbf{essoufflement marqué} : on ne peut plus dire que quelques mots sans reprendre son souffle. Exemples : course à pied, football en match, natation rapide, port de charges lourdes en côte. Une minute d'activité soutenue « vaut » environ deux minutes d'activité modérée en termes de bénéfices.
\end{definition}

\begin{methode}[Évaluer l'intensité d'une activité : le test de la parole]
\begin{enumerate}
\item Réalise l'activité pendant au moins 3 minutes.
\item Essaie de tenir une conversation :
\begin{itemize}
\item tu parles \emph{et} chantes sans gêne $\rightarrow$ intensité \textbf{légère} ;
\item tu parles normalement mais tu ne peux plus chanter $\rightarrow$ intensité \textbf{modérée} ;
\item tu ne peux dire que quelques mots entre deux respirations $\rightarrow$ intensité \textbf{soutenue}.
\end{itemize}
\item Note l'activité et sa durée dans ton carnet de suivi.
\end{enumerate}
\end{methode}

La dépense énergétique varie fortement selon l'activité, comme le montre l'histogramme ci-dessous.

\begin{popfigure}
\begin{center}
\begin{tikzpicture}
\begin{axis}[
    ybar,
    width=0.95\linewidth,
    height=7cm,
    bar width=0.55cm,
    ymin=0, ymax=550,
    ylabel={Dépense énergétique (kcal/h)},
    symbolic x coords={Sommeil, Assis, Marche, Ménage, Football, Course},
    xtick=data,
    x tick label style={rotate=20, anchor=east, font=\small},
    ytick={0,100,200,300,400,500},
    axis lines=left,
    ymajorgrids=true,
    grid style={popline!40},
    every node near coord/.append style={font=\small, color=popdark},
    nodes near coords,
]
\addplot[fill=popblue!70, draw=popblue] coordinates {(Sommeil,60) (Assis,80) (Marche,250) (Ménage,200) (Football,450) (Course,600)};
\end{axis}
\end{tikzpicture}
\end{center}
\legende{Dépense énergétique approximative (en kcal par heure) d'un adolescent d'environ \SI{55}{kg} pour diverses activités : sommeil, position assise, marche rapide, ménage actif, football, course à pied. Plus l'activité sollicite de masses musculaires et accélère le cœur, plus la dépense est élevée.}
\end{popfigure}

\courssub{Les recommandations de l'OMS}

\textbf{Problème de santé publique.} Comment les scientifiques savent-ils « combien » d'activité physique il faut pour rester en bonne santé ? La démarche est épidémiologique : on suit de grandes populations pendant des années, on mesure leur niveau d'activité physique (questionnaires, accéléromètres portés au poignet), puis on compare la fréquence d'apparition des maladies (diabète de type 2, infarctus, certains cancers) entre les groupes actifs et inactifs. \textbf{Données} : les études montrent que le risque de maladie diminue fortement à partir d'un certain seuil d'activité, puis se stabilise. \textbf{Interprétation} : il existe une « dose » minimale protectrice d'activité physique. \textbf{Conclusion} : l'OMS a traduit ces données en recommandations chiffrées, valables dans le monde entier.

\begin{propriete}[Recommandations de l'OMS (admises)]
\begin{itemize}
\item \textbf{Enfants et adolescents de 5 à 17 ans} : au moins \cle{60 minutes par jour} d'activité physique d'intensité \textbf{modérée à soutenue}, réparties dans la journée (elles peuvent être fractionnées en plusieurs périodes).
\item \textbf{Adultes de 18 à 64 ans} : au moins \cle{150 à 300 minutes par semaine} d'activité d'intensité \textbf{modérée}, \emph{ou} \cle{75 à 150 minutes par semaine} d'activité d'intensité \textbf{soutenue} (ou une combinaison équivalente), \textbf{plus} des exercices de \cle{renforcement musculaire} au moins \textbf{2 fois par semaine}.
\end{itemize}
Le respect de ces recommandations \textbf{réduit significativement le risque de maladies non transmissibles (MNT)} : obésité, diabète de type 2, maladies cardiovasculaires, certains cancers. (Recommandation admise : sa justification épidémiologique n'est pas exigible au Probatoire, mais les chiffres, eux, doivent être connus.)
\end{propriete}

\begin{exemplebox}[Un exemple chiffré pour bien comprendre]
Marcher 30 minutes à \SI{5}{km/h} (marche rapide) dépense environ \textbf{120 à 150 kcal} ; jouer au football pendant 1 heure en dépense \textbf{400 à 500 kcal}. À titre de comparaison, une bouteille de soda de \SI{33}{cL} apporte environ 140 kcal de sucres : \textbf{l'équivalent du soda quotidien est « effacé » par environ 35 minutes de marche rapide} (selon l'histogramme ci-dessus : 250 kcal/h $\rightarrow$ 140 kcal $\approx$ 0,56 h). Attention toutefois : il est bien plus facile de ne pas boire le soda que de « brûler » ses calories ensuite — l'activité physique et l'alimentation équilibrée travaillent en équipe, l'une ne remplace pas l'autre.
\end{exemplebox}

\courssub{Les bénéfices de l'activité physique régulière}

Pourquoi bouger protège-t-il ? Parce que l'activité physique agit sur plusieurs mécanismes physiologiques précis :

\begin{itemize}
\item \textbf{Contrôle du poids} : la dépense d'énergie augmente, ce qui rétablit l'équilibre entre les apports et les dépenses et limite le stockage des graisses dans le tissu adipeux.
\item \textbf{Amélioration de la sensibilité à l'insuline} : la contraction musculaire permet aux cellules musculaires de capter le glucose sanguin \emph{même avec moins d'insuline} (translocation des transporteurs GLUT4 indépendamment de l'insuline). La glycémie est mieux régulée, ce qui prévient directement le \cle{diabète de type 2}.
\item \textbf{Santé cardiovasculaire} : le cœur, qui est un muscle, se renforce (augmentation du volume d'éjection systolique, baisse de la fréquence cardiaque de repos) ; la tension artérielle diminue ; le taux de « bon » cholestérol (HDL) augmente. Le risque d'hypertension, d'infarctus et d'accident vasculaire cérébral diminue.
\item \textbf{Santé mentale et concentration scolaire} : l'effort physique stimule la libération d'endorphines (molécules du bien-être), réduit le stress et l'anxiété, améliore le sommeil et la concentration en classe. Un élève actif apprend mieux.
\item \textbf{Renforcement des os et des muscles} : les activités avec impact (marche, course, saut) stimulent la fixation du calcium dans les os, essentielle pendant l'adolescence, période de construction du capital osseux.
\end{itemize}

\begin{savaistu}[L'inactivité physique tue plus qu'on ne le croit]
L'OMS estime que l'\cle{inactivité physique} est le \textbf{4\ieme{} facteur de risque de mortalité dans le monde}, responsable d'environ \textbf{3 à 5 millions de décès par an}, un niveau comparable au tabagisme. À l'échelle mondiale, plus de 80 \% des adolescents n'atteignent pas les 60 minutes quotidiennes recommandées. Au Cameroun, l'urbanisation rapide (transports motorisés, écrans, disparition des jeux de plein air) fait progresser cette épidémie silencieuse, alors que la solution est gratuite : bouger.
\end{savaistu}

\courssub{Intégrer l'activité physique au quotidien camerounais}

Pas besoin de salle de sport ni d'équipement coûteux : le quotidien offre de nombreuses occasions de bouger. Voici des stratégies concrètes et réalistes :

\begin{itemize}
\item \textbf{Marcher vers l'école} : 15 minutes de marche rapide le matin et 15 minutes le soir totalisent déjà la moitié des 60 minutes recommandées.
\item \textbf{Travaux domestiques actifs} : balayer, laver le linge à la main, tirer et porter l'eau, aller au marché à pied comptent comme activité d'intensité modérée.
\item \textbf{Jeux et sports de quartier} : football, handball, course, saut à la corde, jeux de poursuite entre amis après les cours.
\item \textbf{Danse traditionnelle} : les danses traditionnelles (bikutsi, makossa, assiko, danses des cérémonies) sont d'excellentes activités d'intensité modérée à soutenue, tout en cultivant notre patrimoine.
\item \textbf{Jardinage et travaux champêtres} : désherber, planter, récolter au champ familial sollicitent l'ensemble du corps.
\item \textbf{Escaliers plutôt que véhicule} : monter les escaliers, descendre du taxi un arrêt plus tôt.
\end{itemize}

\begin{popfigure}
\begin{center}
\begin{tikzpicture}[font=\small, scale=0.7, transform shape]
% En-têtes des jours
\foreach \i/\j in {0/Lundi, 1/Mardi, 2/Mercredi, 3/Jeudi, 4/Vendredi, 5/Samedi, 6/Dimanche} {
  \node[fill=popblue, text=white, minimum width=2.2cm, minimum height=0.7cm, anchor=north west] at (\i*2.3,0) {\j};
}
% Contenus
\node[fill=popgreenL, draw=popgreen, rounded corners, minimum width=2.0cm, minimum height=1.5cm, anchor=north west, align=center] at (0.1,-0.8) {Marche école\\ 30 min\\ + football 30 min};
\node[fill=popgreenL, draw=popgreen, rounded corners, minimum width=2.0cm, minimum height=1.5cm, anchor=north west, align=center] at (2.4,-0.8) {Marche école\\ 30 min\\ + ménage 30 min};
\node[fill=popgreenL, draw=popgreen, rounded corners, minimum width=2.0cm, minimum height=1.5cm, anchor=north west, align=center] at (4.7,-0.8) {Marche école\\ 30 min\\ + danse 30 min};
\node[fill=popgreenL, draw=popgreen, rounded corners, minimum width=2.0cm, minimum height=1.5cm, anchor=north west, align=center] at (7.0,-0.8) {Marche école\\ 30 min\\ + jeux 30 min};
\node[fill=popgreenL, draw=popgreen, rounded corners, minimum width=2.0cm, minimum height=1.5cm, anchor=north west, align=center] at (9.3,-0.8) {Marche école\\ 30 min\\ + corde à sauter\\ 30 min};
\node[fill=poporangeL, draw=poporange, rounded corners, minimum width=2.0cm, minimum height=1.5cm, anchor=north west, align=center] at (11.6,-0.8) {Match de football\\ 60 min\\ (soutenu)};
\node[fill=popblueL, draw=popblue, rounded corners, minimum width=2.0cm, minimum height=1.5cm, anchor=north west, align=center] at (13.9,-0.8) {Jardinage +\\ marche au marché\\ 60 min};
% Bandeau total
\node[fill=poppurple, text=white, rounded corners, minimum width=16.0cm, minimum height=0.8cm, anchor=north west, align=center] at (0.1,-2.6) {Total : 60 min/jour d'activité modérée à soutenue — recommandation OMS 5--17 ans respectée chaque jour};
\end{tikzpicture}
\end{center}
\legende{Exemple de semaine type d'un élève de Première respectant les recommandations de l'OMS : chaque jour totalise au moins 60 minutes d'activité physique d'intensité modérée à soutenue, en combinant déplacements actifs, tâches domestiques, jeux et sports. Aucun équipement particulier n'est nécessaire.}
\end{popfigure}

\courssub{Limiter les comportements sédentaires}

\begin{definition}[Comportement sédentaire]
Un \cle{comportement sédentaire} est toute activité éveillée réalisée en position assise ou allongée, avec une dépense d'énergie très faible (\textless 1,5 MET) : regarder la télévision, utiliser le téléphone ou l'ordinateur, rester assis de longues heures. Attention : on peut être « actif » une heure par jour \emph{et} sédentaire le reste du temps — les deux comportements ont des effets indépendants sur la santé.
\end{definition}

Recommandations pratiques pour réduire la sédentarité :

\begin{itemize}
\item \textbf{Limiter le temps d'écran récréatif} (téléphone, télévision, jeux vidéo) à environ 2 heures par jour au maximum ;
\item \textbf{Faire une pause active toutes les heures} : se lever, marcher 3 à 5 minutes, s'étirer — même pendant les révisions du Probatoire ;
\item Remplacer certains trajets motorisés courts par la marche ;
\item Privilégier les jeux de plein air aux jeux sur écran pendant les loisirs.
\end{itemize}

\begin{exoresolu}[La semaine de Junior est-elle conforme aux recommandations de l'OMS ?]
\textbf{Énoncé.} Junior, 16 ans, relève ses activités d'une semaine : marche rapide vers le lycée 20 min le matin et 20 min le soir, du lundi au vendredi ; ménage actif 30 min le mercredi ; football 45 min le samedi ; le dimanche, il reste assis devant la télévision et son téléphone environ 7 heures, sans autre activité. (a) Calcule sa durée totale d'activité physique hebdomadaire. (b) Respecte-t-il la recommandation de l'OMS pour son âge ? (c) Identifie le comportement à corriger et propose deux solutions concrètes.
\tcblower
\textbf{Solution détaillée.}

(a) \textbf{Durée totale d'activité physique.}
\begin{align*}
\text{Marche} &= 40\ \text{min/j} \times 5\ \text{j} = 200\ \text{min} \\
\text{Ménage} &= 30\ \text{min} \\
\text{Football} &= 45\ \text{min} \\
\text{Total} &= 200 + 30 + 45 = 275\ \text{min/semaine}
\end{align*}

(b) \textbf{Comparaison avec la recommandation OMS.} À 16 ans, Junior doit atteindre \textbf{au moins 60 min par jour}, soit $60 \times 7 = 420$ min/semaine réparties sur \emph{chaque} jour. Or $275 < 420$ : la quantité totale est insuffisante, et surtout la répartition est mauvaise — le dimanche, il n'a aucune activité. \textbf{La recommandation n'est donc pas respectée.}

(c) \textbf{Comportement à corriger.} Le dimanche, Junior cumule 7 heures de comportement sédentaire (écrans) sans pause active. \textbf{Solutions concrètes} : marcher 30 minutes jusqu'au marché ou chez un ami ; participer à une partie de football ou à une danse le dimanche après-midi ; instaurer une pause active de 5 minutes toutes les heures devant les écrans ; limiter le temps d'écran récréatif à 2 heures.
\end{exoresolu}

\begin{aretenir}[L'essentiel de la section]
\begin{itemize}
\item L'\cle{activité physique} est tout mouvement musculaire augmentant la dépense d'énergie ; l'exercice est planifié, le sport est codifié — mais les activités utilitaires (marche, ménage, jardinage) comptent autant que le sport.
\item Intensité \cle{modérée} : on peut parler mais pas chanter ; intensité \cle{soutenue} : essoufflement marqué.
\item OMS : \textbf{60 min/jour} pour les 5--17 ans ; \textbf{150 à 300 min/semaine} d'activité modérée (ou 75 à 150 min soutenue) + renforcement musculaire 2 fois/semaine pour les adultes.
\item Bénéfices : contrôle du poids, meilleure sensibilité à l'insuline (GLUT4), cœur et vaisseaux protégés, bien-être mental et meilleure concentration scolaire.
\item Il faut aussi \textbf{limiter la sédentarité} : moins d'écrans, une pause active toutes les heures.
\end{itemize}
\end{aretenir}

\begin{exercice}[À toi de jouer]
(a) Classe les activités suivantes par intensité (légère, modérée, soutenue) : course à pied ; marche rapide vers l'école ; lecture assise ; vélo tranquille ; match de football ; vaisselle. Justifie avec le test de la parole. (b) Une adulte fait 4 séances de 40 minutes de marche rapide par semaine, sans renforcement musculaire. Respecte-t-elle toutes les recommandations de l'OMS ? Justifie et propose un ajustement.
\end{exercice}

\begin{corrige}[Exercice]
(a) \textbf{Légère} : lecture assise (on peut parler et chanter) ; vaisselle (effort faible, on peut chanter). \textbf{Modérée} : marche rapide, vélo tranquille (on parle mais on ne chante plus). \textbf{Soutenue} : course à pied, match de football (essoufflement marqué, quelques mots seulement entre deux respirations).
(b) Elle totalise $4 \times 40 = 160$ min d'activité modérée : la borne basse (150 min) est atteinte, donc la recommandation d'endurance est respectée. \textbf{Mais} il manque le \cle{renforcement musculaire} au moins 2 fois par semaine. Ajustement : ajouter deux séances courtes d'exercices de renforcement (pompes, squats, gainage) et, idéalement, augmenter progressivement vers 200--300 min hebdomadaires pour un bénéfice maximal.
\end{corrige}

\coursec{Élaborer un programme de prévention}

Après avoir identifié les recommandations nutritionnelles et les lignes directrices en matière d'activité physique, nous disposons désormais des références théoriques nécessaires. Mais la connaissance ne suffit pas : il faut la transformer en action concrète. Comment un adolescent ou un jeune adulte, confronté à la réalité de son emploi du temps, de son budget et de ses habitudes culturelles (le \emph{ndolé} le dimanche, le \emph{suya} au coin de la rue, les boissons sucrées au frais), peut-il construire un projet de santé personnel, réaliste et durable ? C'est l'objet de cette dernière section : passer de la norme à la pratique individualisée.

\courssub{La démarche en quatre étapes : du diagnostic au suivi}

L'élaboration d'un programme de prévention personnalisé suit une démarche scientifique rigoureuse, calquée sur la méthode d'investigation : on observe (diagnostic), on hypothèse des objectifs, on expérimente (plan d'action), on interprète les résultats (suivi) et on conclut (ajustement). Cette boucle itérative garantit l'adaptation aux spécificités de chaque individu.

\begin{propriete}[Rappel : Estimation des besoins énergétiques journaliers (BEJ) -- Admis au Probatoire]
Pour un adolescent/adulte, le BEJ se calcule en deux temps :
\begin{enumerate}
    \item \textbf{Métabolisme de base (MB)} (kcal/j) : énergie au repos strict.
        \begin{itemize}
            \item Garçon/Homme : $\text{MB} \approx 1 \times \text{masse (kg)} \times 24$ (ou formule de Harris-Benedict/Black et al. plus précise).
            \item Fille/Femme : $\text{MB} \approx 0,9 \times \text{masse (kg)} \times 24$.
        \end{itemize}
    \item \textbf{BEJ} = $\text{MB} \times \text{Niveau d'Activité Physique (NAP)}$.
        \begin{center}
        \begin{tabular}{|l|c|}\hline
        \textbf{Niveau d'activité} & \textbf{NAP} \\\hline
        Sédentaire (bureau, école, pas de sport) & 1,4 -- 1,5 \\\hline
        Actif (marche 30--60 min/j, travail manuel léger) & 1,6 -- 1,7 \\\hline
        Très actif (sport quotidien intense, travail manuel lourd) & 1,8 -- 2,0 \\\hline
        \end{tabular}
        \end{center}
\end{enumerate}
\emph{Exemple Amina (71,5 kg, sédentaire) : MB $\approx 0,9 \times 71,5 \times 24 \approx 1544$ kcal ; BEJ $\approx 1544 \times 1,45 \approx 2240$ kcal/j.}
\end{propriete}

\begin{methode}[Élaborer un programme combinant alimentation équilibrée et activité physique]
\begin{enumerate}
    \item \textbf{Calculer l'IMC, le tour de taille et analyser la ration habituelle} : noter tout ce qui est mangé et bu pendant 3 jours (dont un week-end), estimer les portions (poids ou mesures usuelles), calculer l'apport énergétique moyen (via table de composition Ciqual/ANSES ou application validée) et le comparer aux BEJ calculés ci-dessus. Calculer l'IMC = $\frac{\text{masse (kg)}}{\text{taille (m)}^2}$ et mesurer le tour de taille (niveau du dernier bord costal, en fin d'expiration).
    \item \textbf{Identifier 2 à 3 excès majeurs à corriger} : repérer les écarts les plus fréquents et impactants (ex. : boisson sucrée quotidienne, absence de légumes au déjeuner, collations grasses/sucrées, temps d'écran $> 4$ h/jour sans activité physique structurée).
    \item \textbf{Fixer des objectifs SMART} : Spécifiques, Mesurables, Atteignables, Réalistes, Temporels (voir encadré ci-dessous).
    \item \textbf{Planifier menus et séances d'activité sur 7 jours} : proposer des substitutions concrètes (eau $\leftarrow$ soda, fruit $\leftarrow$ gâteau, marche $\leftarrow$ taxi-moto pour 1 km) et répartir au minimum 60 min d'activité modérée à soutenue par jour (adolescent) ou 30 min (adulte), en variant endurance, renforcement musculaire et souplesse.
    \item \textbf{Prévoir des indicateurs de suivi simples} : pesée hebdomadaire (même jour, mêmes conditions), recalcul de l'IMC mensuel, mesure du tour de taille mensuel, carnet d'activité (durée, type, ressenti), compteur de pas (objectif progressif vers 8 000--10 000 pas/jour).
    \item \textbf{Réévaluer après 4 semaines} : comparer les indicateurs aux objectifs. Si progression : maintenir ou augmenter légèrement la charge. Si stagnation : identifier l'obstacle (portions cachées, sédentarité compensatoire, stress, sommeil) et ajuster une seule variable à la fois.
\end{enumerate}
\end{methode}

\begin{popfigure}
\begin{tikzpicture}[
    node distance=1.2cm and 2.2cm,
    box/.style={draw=popblue, thick, rounded corners, fill=popblueL, minimum width=2.6cm, minimum height=1.3cm, align=center, font=\small},
    arrow/.style={-Stealth, thick, popdark},
    decision/.style={diamond, draw=poporange, thick, fill=poporangeL, minimum width=2.5cm, minimum height=1.2cm, align=center, font=\small, aspect=2}
]
% Nodes
\node[box, align=center] (diag){1. DIAGNOSTIC\\IMC, Tour de taille,\\Ration 3 jours,\\Niveau d'activité (NAP)};
\node[box, right=of diag, align=center] (obj){2. OBJECTIFS SMART\\Perte 0,5 kg/sem,\\60 min act./jour en 4 sem,\\Supprimer 1 soda/jour};
\node[box, right=of obj, align=center] (plan){3. PLAN HEBDOMADAIRE\\Menus équilibrés +\\  Séances d'activité\\ (Endurance/Force/Souplesse)};
\node[box, right=of plan, align=center] (suivi){4. SUIVI \& AJUSTEMENT\\Pesée hebdo, Carnet,\\ Compteur pas, IMC mensuel};
\node[decision, below=of plan, align=center] (eval){Résultats\\ à 4 semaines ?};
\node[box, below=of eval, align=center] (ajust){AJUSTER\\ (Modifier 1 variable)\\ OU CONSOLIDER};

% Arrows
\draw[arrow] (diag) -- (obj) node[midway, above, font=\small] {Analyse des écarts\par vs Besoins};
\draw[arrow] (obj) -- (plan) node[midway, above, font=\small] {Traduction concrète};
\draw[arrow] (plan) -- (suivi) node[midway, above, font=\small] {Mise en œuvre};
\draw[arrow] (suivi) -- (eval);
\draw[arrow] (eval) -- node[left, font=\small, align=center]{Objectifs\\ atteints} (ajust);
\draw[arrow] (eval) -| node[near start, above, font=\small, align=center]{Stagnation /\\ Aggravation} (diag);
\draw[arrow] (ajust) |- (plan) node[near start, below, font=\small] {Nouveau cycle};
\end{tikzpicture}
\legende{Organigramme de la démarche itérative en 4 étapes pour l'élaboration et le suivi d'un programme de prévention personnalisé. La flèche de retour vers le diagnostic symbolise la réévaluation complète en cas d'échec partiel.}
\end{popfigure}

\courssub{Fixer des objectifs réalistes et mesurables : la méthode SMART}

Un objectif vague (\og je veux maigrir \fg{}, \og je vais faire du sport \fg{}) est voué à l'échec car il ne permet ni évaluation ni motivation durable. La méthode \textbf{SMART} (\cle{Spécifique}, \cle{Mesurable}, \cle{Atteignable}, \cle{Réaliste}, \cle{Temporel}) structure la contractualisation avec soi-même.

\begin{propriete}[Perte de poids durable et sans danger -- Admise au Probatoire]
Une perte de poids durable et sans danger pour la santé se situe entre \textbf{0,5 et 1 kg par semaine}. Au-delà, le risque de reprise rapide (\emph{effet yo-yo}), de perte de masse musculaire, de carences en micronutriments (fer, calcium, vitamines) et de troubles du comportement alimentaire augmente significativement. Un déficit énergétique modéré de \textbf{300 à 500 kcal/jour} (combinant restriction modérée et dépense supplémentaire par l'activité physique) est la cible physiologique de référence.
\end{propriete}

\begin{attention}[Interprétation de l'IMC chez l'adolescent : vigilance !]
Chez l'adolescent (jusqu'à 18--19 ans), l'IMC \textbf{ne s'interprète pas avec les seuils fixes de l'adulte} (25 / 30). On utilise les \textbf{courbes de référence par âge et sexe} (OMS 2007 / INSERM) en percentiles :
\begin{itemize}
    \item \textbf{Surpoids :} IMC $>$ percentile 85 (ou $> +1$ écart-type).
    \item \textbf{Obésité :} IMC $>$ percentile 97 (ou $> +2$ écarts-types).
\end{itemize}
À 17 ans (comme Amina), les valeurs seuils approchent celles de l'adulte (surpoids $\approx$ 25--26, obésité $\approx$ 30), mais la rigueur scientifique impose de consulter la courbe de croissance. Le tour de taille s'interprète aussi par percentiles (risque métabolique accru $> P90$).
\end{attention}

\begin{exemplebox}[Programme chiffré d'Amina -- Application de la méthode SMART]
\textbf{Profil :} Amina, 17 ans, élève en 1ère D à Yaoundé. Taille : 1,62 m. Masse : 71,5 kg. IMC = $71,5 / 1,62^2 \approx 27,2$. Selon la courbe OMS fille 17 ans, cela la place au \textbf{percentile 90 (surpoids)}. Tour de taille : 88 cm ($> P90$, risque métabolique élevé). Ration habituelle : sauté petit-déjeuner, cantine midi (riz + sauce grasse + viande), goûter (beignet + boisson gazeuse 33 cl), dîner familial (pâte de maïs + sauce arachide). Activité : 15 min marche aller-retour école, 4 h écrans/jour. NAP estimé 1,45. BEJ $\approx$ 2240 kcal/j. Apport estimé $\approx$ 2600 kcal/j (excès $\approx$ 360 kcal/j).

\textbf{Diagnostic :} Excès identifiés : (1) Boisson gazeuse quotidienne ($\approx 140$ kcal), beignet ($\approx 250$ kcal) $\rightarrow$ $\approx 390$ kcal/jour en sucres ajoutés/gras. (2) Sauter le petit-déjeuner $\rightarrow$ fringale goûter. (3) Sédentarité marquée (NAP 1,45).

\textbf{Objectifs SMART à 4 semaines :}
\begin{itemize}
    \item Perdre $2$ kg (soit $0,5$ kg/semaine) $\rightarrow$ Masse cible $69,5$ kg (IMC $\approx 26,5$, percentile $\approx 85$).
    \item Atteindre 45 min d'activité modérée 5 j/7 + 2 séances de 30 min renforcement/danse (NAP visé 1,65).
    \item Supprimer la boisson gazeuse quotidienne, remplacer le beignet par un fruit de saison.
    \item Prendre un petit-déjeuner complet 6 j/7.
\end{itemize}

\textbf{Calcul du déficit visé :}
\begin{align*}
\text{Suppression soda (33 cl) + beignet} &\approx -390 \text{ kcal/j} \\
\text{Ajout petit-déj équilibré (pain + œuf + fruit)} &\approx +300 \text{ kcal/j} \quad (\text{mais réparti, coupe la fringale}) \\
\text{Bilan alimentaire net} &\approx -90 \text{ kcal/j} \\
\text{Marche rapide 45 min (5 j/sem, MET 4,5)} &\approx 5 \times (4,5 \times 71,5 \times 0,75) \approx 1200 \text{ kcal/sem} \approx -171 \text{ kcal/j (moy.)} \\
\text{Danse/Renforcement 30 min (2 j/sem, MET 5,5)} &\approx 2 \times (5,5 \times 71,5 \times 0,5) \approx 390 \text{ kcal/sem} \approx -56 \text{ kcal/j (moy.)} \\
\textbf{Déficit moyen total (aliment + activité)} &\approx \mathbf{317 \text{ kcal/j}} \quad (\text{dans la cible } 300\text{--}500 \text{ kcal/j}) \\
\text{Ajustement portions dîner (moins sauce/huile)} &\approx -100 \text{ kcal/j} \rightarrow \textbf{Total } \approx 417 \text{ kcal/j}
\end{align*}
\emph{Résultat attendu :} $\approx 0,5$ kg/semaine. Après 6 mois : IMC $\approx 24$ (percentile $\approx 50$, normale).
\end{exemplebox}

\begin{attention}[Piège fréquent : le régime drastique sans activité physique]
Proposer ou suivre un régime très restrictif ($< 1200$ kcal/j pour un adolescent) \textbf{sans activité physique} entraîne :
\begin{itemize}
    \item Une perte de \textbf{masse musculaire} (le corps puise dans les protéines musculaires pour maintenir la glycémie via la néoglucogenèse), ce qui abaisse le métabolisme de base.
    \item Une \textbf{frustration} psychologique majeure, source de craquages compulsifs (hyperphagie boulimique).
    \item Un \textbf{effet yo-yo} quasi certain : à la reprise d'une alimentation normale, le métabolisme ralenti favorise un stockage rapide des graisses, souvent au-delà du poids initial.
    \item Des \textbf{carences} (fer $\rightarrow$ anémie, calcium $\rightarrow$ fragilité osseuse future, vitamines B $\rightarrow$ fatigue cognitive).
\end{itemize}
\textbf{Règle d'or :} On ne maigrit pas \emph{en mangeant moins}, on maigrit \emph{en mangeant mieux ET en bougeant plus}. L'activité physique préserve la masse musculaire, augmente la dépense au repos et régule l'appétit par la sécrétion de peptides (PYY, GLP-1) et la sensibilité à l'insuline.
\end{attention}

\courssub{Construire le plan hebdomadaire : menus et activités physiques}

Le plan hebdomadaire traduit les objectifs en actes quotidiens. Il doit respecter la culture alimentaire camerounaise (coût, disponibilité, goût) et l'emploi du temps scolaire. L'équilibre se joue sur la semaine, pas sur un seul repas.

\begin{center}
\begin{tabularx}{\linewidth}{|>{\bfseries}l|X|X|}\hline
\textbf{Jour} & \textbf{Alimentation (Exemple type Amina)} & \textbf{Activité Physique} \\\hline
Lundi & Pdej: Pain complet + œuf brouillé + mangue + eau.\\
Midi: Riz + légumes (feuilles) + poisson grillé + eau.\\
Goûter: Banane + arachides (30 g).\\
Soir: Pâte maïs + sauce okro (peu d'huile) + viande maigre. & Marche rapide 45 min (soir).\\
Renforcement: Gainage 3$\times$30s, Squats 3$\times$15. \\\hline
Mardi & Pdej: Bouillie maïs/soja (peu sucrée) + orange.\\
Midi: Haricots + plantain bouilli + salade tomate/oignon.\\
Goûter: Yaourt nature + fruit.\\
Soir: Légumes sautés + blanc de poulet + petit pain. & \textbf{Repos actif} : Étirements 15 min + marche 20 min. \\\hline
Mercredi & Pdej: Pain + beurre d'arachide (fine couche) + papaye.\\
Midi: Cantine : choisir riz + sauce légumes + viande (retirer peau/gras visible).\\
Goûter: Carotte/concombre + fromage blanc.\\
Soir: Soupe de poisson + manioc (modéré). & Danse moderne / Zumba 45 min (club école ou maison). \\\hline
Jeudi & Pdej: Œufs + pain + ananas.\\
Midi: Pâtes complètes + sauce tomate/viande hachée maigre + salade.\\
Goûter: Pomme + amandes.\\
Soir: Igname bouilli + sauce feuille (éru/ndolé allégé). & Marche rapide 45 min (soir).\\
Pompes (genoux) 3$\times$10. \\\hline
Vendredi & Pdej: Bouillie millet + banane.\\
Midi: Maïs + haricots (kwem) + salade.\\
Goûter: Jus de fruit MAISON (sans sucre ajouté) + galette maïs.\\
Soir: Poisson braisé + légumes vapeur + riz (petite portion). & \textbf{Sortie nature / Sport collectif} : Football, handball, randonnée Mont Fébé (60--90 min). \\\hline
Samedi & Pdej: Crêpes millet (peu d'huile) + miel + citron.\\
Midi: Repas familial : se servir 1/2 assiette légumes, 1/4 féculents, 1/4 protéines.\\
Goûter: Fruit de saison.\\
Soir: Léger : Salade composée + œufs durs. & Danse / Renforcement circuit training 30 min.\\
Marche détente 30 min (famille). \\\hline
Dimanche & Pdej: Libre (raisonnable) + fruit.\\
Midi: Repas fête (ndolé, koki) : \textbf{portions contrôlées}, pas de resservice, beaucoup d'eau.\\
Soir: Très léger (potage + fruit) si midi copieux. & \textbf{Repos complet} ou balade douce 30 min.\\
Préparation repas semaine (batch cooking). \\\hline
\end{tabularx}
\end{center}

\courssub{Outils de suivi et ajustement du programme}

La motivation fluctue ; seul le suivi objectif permet de maintenir le cap. Les outils doivent être simples, accessibles et non intrusifs.

\begin{itemize}
    \item \textbf{Carnet alimentaire (papier ou application type MyFitnessPal / Yazio / FatSecret) :} Noter \emph{avant} de manger. Prendre en photo l'assiette aide à la prise de conscience des portions.
    \item \textbf{Pesée hebdomadaire :} Une fois par semaine, à jeun, après miction, même balance. Ne pas se peser tous les jours (variations hydriques $\pm 1$ kg masquant la tendance).
    \item \textbf{Recalcul de l'IMC et tour de taille mensuels :} L'IMC ne distingue pas gras/muscle. Le tour de taille (mesuré au niveau du dernier bord costal, milieu entre la dernière côte et la crête iliaque) reflète la graisse viscérale (risque métabolique). Objectif : $< P90$ percentile (ado) ou $< 80$ cm (femme), $< 94$ cm (homme) adulte.
    \item \textbf{Compteur de pas (smartphone/montre) :} Objectif progressif : 5 000 $\rightarrow$ 7 000 $\rightarrow$ 10 000 pas/jour.
    \item \textbf{Journal de bord ressenti :} Énergie, qualité du sommeil, humeur, douleurs articulaires. La santé ne se résume pas au poids.
\end{itemize}

\begin{popfigure}
\begin{tikzpicture}
\begin{axis}[
    width=\linewidth, height=8cm,
    xlabel={Mois}, ylabel={IMC (kg/m$^2$)},
    xmin=0, xmax=6, ymin=22, ymax=28,
    xtick={0,1,2,3,4,5,6},
    ytick={22,23,24,25,26,27,28},
    grid=major, grid style={popline},
    axis lines=left,
    tick label style={font=\small},
    label style={font=\small},
    legend style={at={(0.98,0.98)}, anchor=north east, font=\small, fill=popcream, draw=popmuted},
    domain=0:6, samples=100,
]
% Zones de référence (Corrigé : remplissage entre bornes)
% Zone Normale (18.5 - 25) -> on affiche 22-25
\addplot[popgreenL, fill opacity=0.2, draw=none, domain=0:6] {25} \closedcycle;
% Zone Surpoids (25 - 30) -> on affiche 25-27.5 (max graph)
\addplot[poporangeL, fill opacity=0.15, draw=none, domain=0:6] {27.5} \closedcycle;
% On remet le fond blanc sur la zone normale pour que le vert soit visible (superposition)
% Astuce : tracer la zone surpoids D'ABORD, puis la zone normale. Ou utiliser fill between.
% Solution simple ici : tracer zone surpoids (25 à 27.5) via coordonnées.
\addplot[poporangeL, fill opacity=0.15, draw=none] coordinates {(0,25) (6,25) (6,27.5) (0,27.5)} \closedcycle;
\addplot[popgreenL, fill opacity=0.2, draw=none] coordinates {(0,22) (6,22) (6,25) (0,25)} \closedcycle;

\node[popgreen, font=\tiny, anchor=north east] at (axis cs: 5.8, 24.5) {Normale ($<25$)};
\node[poporange, font=\tiny, anchor=south east] at (axis cs: 5.8, 26.5) {Surpoids ($25\text{--}30$)};

% Courbe théorique : décroissance exponentielle vers 24
\addplot[popcurve, popblue, very thick, domain=0:6] {24 + 3.2*exp(-0.45*x)};
\addlegendentry{Évolution IMC Amina (modèle)}

% Points de mesure mensuels
\addplot[only marks, mark=*, mark size=2.5, popblue] coordinates {
    (0, 27.2)
    (1, 26.2)
    (2, 25.4)
    (3, 24.9)
    (4, 24.5)
    (5, 24.2)
    (6, 24.0)
};
\addlegendentry{Pesées mensuelles réelles}
\end{axis}
\end{tikzpicture}
\legende{Évolution attendue de l'IMC d'Amina sur 6 mois avec le programme (déficit $\approx$ 400--500 kcal/j). La courbe montre une décroissance progressive, rapide au début (perte d'eau/glycogène), puis plus lente, se stabilisant vers la normale (IMC 24). Les points correspondent aux pesées mensuelles. La zone verte indique la normale (IMC $< 25$), la zone orange le surpoids (IMC $25\text{--}30$).}
\end{popfigure}

\courssub{Le rôle de l'entourage, de l'école et la communication de prévention}

La réussite individuelle est fortement conditionnée par l'environnement. L'approche \emph{écologique} de la santé (Bronfenbrenner) rappelle que l'individu est au centre de systèmes imbriqués : famille, école, communauté, politiques publiques.

\begin{itemize}
    \item \textbf{Famille/Entourage :} Soutien logistique (achats, cuisine), modèle parental (les parents actifs ont des enfants actifs), ambiance repas (sans écrans, conviviale). Éviter les commentaires sur le poids (\og tu as grossi \fg{}), privilégier l'encouragement sur les comportements (\og bravo pour ta marche hier \fg{}).
    \item \textbf{École (Lycée/Collège) :} Rôle pivot.
        \begin{itemize}
            \item \textbf{Cantine :} Menus validés par un nutritionniste, portionnement adapté, eau potable gratuite, interdiction boissons sucrées/industrielles.
            \item \textbf{EPS :} 2 à 3 h/sem obligatoires, variées, évaluées sur la progression et l'engagement, pas seulement la performance. Clubs sportifs (mercredi après-midi, vacances).
            \item \textbf{Éducation par les pairs :} \og Ambassadeurs santé \fg{} élèves formés pour animer causeries, défis \og 0 soda \fg{}, marches collectives.
        \end{itemize}
    \item \textbf{Communauté/Quartier :} Aménagement urbain (trottoirs, pistes cyclables, éclairage pour marche sécurisée soir), marchés de produits frais accessibles, événements \og Journée du sport pour tous \fg{} (ex. : Journée Mondiale de l'Activité Physique le 6 avril, Journée Mondiale de la Santé le 7 avril).
\end{itemize}

\begin{savaistu}[Efficacité des programmes communautaires]
Des études menées au Cameroun (ex. : programme \emph{« École en Santé »} appuyé par le MINSANTE/MINESEC et l'OMS) et dans d'autres pays d'Afrique subsaharienne montrent que les interventions \textbf{multi-composantes} (cantine + EPS + éducation nutritionnelle + implication parents + environnement scolaire sain) réduisent l'incidence du surpoids de 15 à 20 \% en 2 ans, là où les conseils individuels isolés n'ont qu'un effet transitoire. La \emph{marche collective du samedi matin} (très populaire à Yaoundé, Douala, Bafoussam) crée une norme sociale positive : \og tout le monde marche, c'est normal, c'est gratuit, c'est convivial \fg{}. C'est le levier le plus puissant pour l'adolescent en quête d'appartenance.
\end{savaistu}

\textbf{Communiquer pour prévenir :} Construire un message de prévention efficace (affiche, sketch, spot radio, causerie éducative) suit des règles simples :
\begin{enumerate}
    \item \textbf{Positif et non culpabilisant :} \og Bouge pour ton énergie ! \fg{} plutôt que \og Ne reste pas assis, tu vas grossir. \fg{}
    \item \textbf{Concret et actionnable :} \og Remplace ton jus de fruit industriel par un fruit entier + un verre d'eau \fg{} (action précise) vs \og Mange sain \fg{} (vague).
    \item \textbf{Adapté au public :} Pour les 1ère D : insister sur le lien \emph{performance cognitive / alimentation / sommeil / activité} (réussite au Probatoire). Utiliser l'humour, les codes culturels (argot, musiques urbaines, influenceurs locaux).
    \item \textbf{Visuel et mémorable :} Une affiche \og Ton assiette idéale : 1/2 légumes, 1/4 féculents, 1/4 protéines \fg{} avec des aliments locaux (feuilles, plantain, poisson, arachide) parle plus qu'une pyramide alimentaire occidentale.
\end{enumerate}

\begin{exoresolu}[Conception d'une action de prévention : « Défi 30 Jours Zéro Soda Lycée »]
\textbf{Énoncé :} En tant que délégué santé de ta classe de 1ère D, tu dois proposer au bureau du Proviseur une action de prévention sur 4 semaines pour réduire la consommation de boissons sucrées au lycée. Détaille la démarche selon les 4 étapes (Diagnostic $\rightarrow$ Objectifs $\rightarrow$ Plan $\rightarrow$ Suivi) et propose un support de communication (affiche ou slogan).

\textbf{Solution détaillée :}

\textbf{1. Diagnostic (Semaine 0) :}
\begin{itemize}
    \item Enquête flash (Google Forms / papier) auprès de 200 élèves : \og Combien de boissons sucrées (soda, jus industriel, boisson énergisante) bois-tu par semaine ? \fg{} $\rightarrow$ Résultat moyen : 5,2 canettes/semaine.
    \item Observation cantine/buvette : 3 frigos pleins de sodas, 1 distributeur d'eau en panne.
    \item IMC moyen classe : 24,5 (limite surpoids). 30 \% IMC $> 25$ (seuil adulte, indicateur de tendance).
\end{itemize}

\textbf{2. Objectifs SMART (4 semaines) :}
\begin{itemize}
    \item Réduire la consommation moyenne de 5,2 à $\leq$ 1 canette/semaine/élève ($-80\%$).
    \item Réparer/mettre en service 3 fontaines à eau filtrée (budget APE/Partenaires).
    \item Sensibiliser 100 \% des élèves de 1ère via causeries de 15 min en classe.
\end{itemize}

\textbf{3. Plan d'action (Semaines 1 à 4) :}
\begin{center}
\begin{tabularx}{\linewidth}{|l|X|}\hline
\textbf{Semaine} & \textbf{Actions} \\\hline
1 (Lancement) & Assemblée générale 10 min : annonce défi. Distribution gourdes « Lycée Zéro Soda » (recherche sponsors). Affichage dans couloirs/cantine. Causeries classe par classe (binômes élèves formés). \\\hline
2 (Renforcement) & Défi « Eau infusée » : concours recettes (citron/menthe, concombre/gingembre, ananas/basilic) -- vote Instagram lycée. Atelier « Lire l'étiquette : sucre caché » (SVT/Chimie). \\\hline
3 (Consolidation) & Journée « Zéro Franc pour le Soda » : l'argent économisé mis en cagnotte pour achat ballons/paniers. Témoignages élèves « Je dors mieux, je cours plus vite ». \\\hline
4 (Clôture) & Enquête post-défi. Remise prix (gourde inox, inscription club sport) aux classes les plus performantes. Restitution au Proviseur + proposition pérennisation (interdiction vente soda dans l'enceinte). \\\hline
\end{tabularx}
\end{center}

\textbf{4. Suivi \& Indicateurs :}
\begin{itemize}
    \item Hebdo : relevé ventes buvette (baisse volume sodas, hausse eau).
    \item Hebdo : sondage rapide 5 questions (auto-déclaration consommation).
    \item Final : Comparaison IMC moyen classe (indicateur long terme).
\end{itemize}

\textbf{Support de communication -- Affiche (Format A3, couleurs vives) :}
\begin{center}
\begin{tikzpicture}
\node[draw=popblue, thick, rounded corners, fill=popcream, inner sep=10pt, align=center, text width=0.85\linewidth] {
    \textbf{\Huge \textcolor{popblue}{DÉFI 30 JOURS}}\\[4pt]
    \textbf{\Large \textcolor{poporange}{ZÉRO SODA}}\\[8pt]
    % Remplacement de l'image par un dessin TikZ simple
    \begin{tikzpicture}[scale=0.8, baseline=-0.5ex]
        \draw[popblue, thick] (0,0) circle (1.5);
        \draw[popblue, thick] (0,0.5) -- (0,-0.5); % straw
        \draw[popblue, thick] (-0.5,1) -- (0.5,1); % lid
        \draw[popgreen, thick, fill=popgreenL] (-0.8,-1.5) rectangle (0.8,-1.8) node[midway, font=\small, text=popdark] {EAU};
        \draw[->, poporange, thick] (1.8,0) -- (3,0);
        \draw[popblue, thick] (4,0) circle (1.5);
        \draw[popblue, thick] (4,0.5) -- (4,-0.5);
        \draw[popblue, thick] (3.5,1) -- (4.5,1);
        \draw[popblue, thick, fill=popblueL] (3.2,-1.5) rectangle (4.8,-1.8) node[midway, font=\small, text=popdark] {SODA};
        \draw[popred, very thick] (3.5,1.5) -- (4.5,-0.5); % Croix rouge
    \end{tikzpicture}
    \\[8pt]
    \textcolor{popdark}{\large \textbf{1 canette = 7 morceaux de sucre \\ Tu ne les mangerais pas à la cuillère !}}\\[6pt]
    \textcolor{popgreen}{\large \textbf{Choisis l'eau, choisis ta forme.}}\\[6pt]
    \textcolor{poppurple}{\large \textbf{\#LyceeZeroSoda \#FormeProbatoire}}\\[6pt]
    \small Partage ta gourde sur Insta @SanteLyceeBiyemAssi
};
\end{tikzpicture}
\end{center}
\legende{Exemple de support de communication (affiche) conçu par les élèves pour le défi. Message positif, visuel choc (équivalence sucre), appel à l'action concret (gourde), hashtag pour viralité.}
\end{exoresolu}

\begin{aretenir}[Synthèse : Clés d'un programme de prévention réussi]
\begin{itemize}
    \item \textbf{Personnalisé :} Parti du diagnostic réel (IMC percentiles, habitudes, environnement, BEJ), pas d'un modèle standard.
    \item \textbf{Progressif et SMART :} Objectifs mesurables (0,5 kg/sem, +15 min act./sem), pas de révolution brutale.
    \item \textbf{Combiné :} Déficit énergétique 300--500 kcal/j = Alimentation (qualité/portions/densité énergétique) + Activité (plaisir/régularité/NAP). L'un ne va pas sans l'autre.
    \item \textbf{Suivi objectif et bienveillant :} Indicateurs simples (poids, tour de taille, pas, carnet, ressenti), réévaluation mensuelle, ajustement sans culpabilité.
    \item \textbf{Ancré dans l'environnement :} Famille, école, quartier comme alliés. La prévention collective (marche, cantine, défis pairs) démultiplie l'efficacité individuelle.
    \item \textbf{Communication positive :} Messages clairs, visuels, valorisants, adaptés à la culture jeune camerounaise.
\end{itemize}
\end{aretenir}

\newpage

\coursec{Activité d'intégration}

\courssub{Objectifs de l'activité}

À l'issue de cette activité d'intégration, tu seras capable de :

\begin{itemize}
\item \cle{calculer} l'indice de masse corporelle (IMC) d'une personne et l'\cle{interpréter} à l'aide de la classification de l'Organisation mondiale de la santé (OMS) ;
\item \cle{évaluer} une ration alimentaire quotidienne à partir d'une table simplifiée de composition des aliments et la comparer aux besoins recommandés ;
\item \cle{identifier} les excès (sucres rapides, lipides) et les carences (fruits, légumes frais, eau) d'une alimentation réelle ;
\item \cle{élaborer} un programme hebdomadaire combinant alimentation équilibrée et activité physique, avec un objectif de perte de masse raisonnable ;
\item \cle{argumenter} et \cle{présenter} un conseil de santé à l'oral, en adoptant une attitude de prévention et non de jugement.
\end{itemize}

Cette activité mobilise les trois savoir-faire de la leçon : évaluer une ration alimentaire, interpréter l'IMC, élaborer un programme de prévention.

\courssub{Situation}

Amina, 17 ans, est élève en classe de Première au lycée de Bafoussam. Lors d'une visite médicale scolaire, l'infirmier relève ses mensurations : \textbf{taille = 1,58 m ; masse = 68 kg}. Il s'inquiète discrètement et demande à Amina de tenir un journal alimentaire et d'activité pendant une journée ordinaire. Voici ce qu'elle a noté.

\textbf{Journal alimentaire d'une journée d'Amina :}

\begin{center}
\rowcolors{1}{popcream}{white}
\begin{tabularx}{\linewidth}{|l|X|l|}
\hline
\rowcolor{popblueL} \textbf{Repas} & \textbf{Contenu} & \textbf{Boisson} \\
\hline
Petit-déjeuner (7 h) & 4 beignets haricots (environ 200 g) & 1 soda de 33 cL \\
\hline
Déjeuner (13 h) & 1 grande assiette de riz à l'huile rouge (300 g de riz cuit, 3 cuillères à soupe d'huile de palme) & 1 soda de 33 cL \\
\hline
Dîner (20 h) & Fufu maïs (250 g) + eru (150 g, préparé avec huile de palme et eau feuilles) & 1 verre d'eau (25 cL) \\
\hline
\end{tabularx}
\end{center}

\textbf{Activité physique et mode de vie :} Amina se rend au lycée en moto-taxi, ne pratique aucun sport, et passe environ \textbf{5 heures par jour} devant un écran (téléphone et télévision), assise ou allongée. Elle ne boit presque pas d'eau en dehors des repas.

\textbf{Table simplifiée de composition des aliments (fournie par l'infirmier) :}

\begin{center}
\rowcolors{1}{popcream}{white}
\begin{tabularx}{\linewidth}{|X|c|c|}
\hline
\rowcolor{popblueL} \textbf{Aliment / boisson} & \textbf{Portion} & \textbf{Énergie (kcal)} \\
\hline
Beignets haricots & 100 g & 350 \\
\hline
Soda sucré & 33 cL & 140 \\
\hline
Riz blanc cuit & 100 g & 130 \\
\hline
Huile de palme & 1 cuillère à soupe (10 g) & 90 \\
\hline
Fufu maïs & 100 g & 160 \\
\hline
Eru (avec huile) & 100 g & 180 \\
\hline
Eau & 1 L & 0 \\
\hline
\end{tabularx}
\end{center}

\textbf{Rappels de référence :} une adolescente de 17 ans modérément active a besoin d'environ \SI{2200}{kcal} par jour ; l'infirmier du lycée recommande \SI{2500}{kcal} comme seuil à ne pas dépasser pour le profil d'Amina. L'OMS recommande aux adolescents au moins \textbf{60 minutes d'activité physique modérée à soutenue par jour} et au moins \textbf{1,5 L d'eau par jour}. Une perte de masse raisonnable et durable est de \textbf{0,5 kg par semaine} au maximum.

\textbf{Problème :} tu es membre du club santé du lycée. Amina accepte que ton groupe l'aide. Comment établir un diagnostic objectif de sa situation, puis construire avec elle un programme réaliste d'alimentation équilibrée et d'activité physique ?

\courssub{Consignes}

Par groupes de 4, et en utilisant les données de la situation :

\begin{enumerate}
\item Calcule l'IMC d'Amina (formule : $IMC = \dfrac{m}{t^2}$, avec $m$ en kg et $t$ en m). Arrondis au dixième.
\item À l'aide de la classification de l'OMS (maigreur : $IMC < 18,5$ ; corpulence normale : $18,5 \leq IMC < 25$ ; surpoids : $25 \leq IMC < 30$ ; obésité : $IMC \geq 30$), conclus sur la corpulence d'Amina et cite deux risques sanitaires encourus si rien ne change.
\item À l'aide de la table de composition, estime l'apport énergétique total de la journée d'Amina, repas par repas, puis au total. Compare ce total aux \SI{2500}{kcal} recommandées.
\item Identifie dans son alimentation : (a) les principales sources d'excès (sucres rapides, lipides) ; (b) les carences (fruits, légumes frais, eau). Justifie chaque réponse par un aliment précis du journal.
\item Évalue le mode de vie d'Amina au regard des recommandations de l'OMS (60 min d'activité physique par jour ; temps d'écran ; hydratation).
\item Sur la grille vierge fournie, élabore un programme hebdomadaire comprenant : \textbf{trois menus types équilibrés} (petit-déjeuner, déjeuner, dîner) à base d'aliments locaux disponibles au marché de Bafoussam (maïs, haricots, ndolé, banane plantain, arachides, fruits de saison, légumes feuilles...), et un \textbf{planning d'activité physique} respectant les 60 min/jour, compatible avec les horaires d'une lycéenne.
\item Justifie que ton programme permet une perte d'environ 0,5 kg par semaine (déficit énergétique modéré + dépense par l'activité physique), sans supprimer aucun repas.
\item Présente ton programme à la classe en 3 minutes. La classe vote le programme le plus réaliste et le mieux argumenté. La restitution collective au tableau donnera une « fiche conseil » validée par tous.
\end{enumerate}

\courssub{Ressources à mobiliser}

\begin{aretenir}[Les outils de la leçon à réactiver]
\begin{itemize}
\item \textbf{Formule de l'IMC :} $IMC = \dfrac{m}{t^2}$ (kg/m$^2$). Classification OMS : maigreur $< 18,5$ ; normale $[18,5\ ; 25[$ ; surpoids $[25\ ; 30[$ ; obésité $\geq 30$.
\item \textbf{Besoins énergétiques :} environ \SI{2200}{kcal/j} pour une adolescente ; une ration équilibrée répartit l'énergie en glucides (50--55 \%), lipides (30--35 \%), protéines (12--15 \%), avec fibres, vitamines, sels minéraux et au moins 1,5 L d'eau.
\item \textbf{Conséquences de la sédentarité et de la mauvaise alimentation :} obésité, diabète de type 2, maladies cardiovasculaires (hypertension, infarctus, AVC).
\item \textbf{Recommandations OMS :} au moins 60 min d'activité physique modérée à soutenue par jour pour les adolescents ; limiter sucres rapides et graisses ; consommer fruits et légumes quotidiennement ; perte de masse visée : 0,5 kg/semaine maximum.
\end{itemize}
\end{aretenir}

\begin{methode}[Évaluer une ration alimentaire et interpréter l'IMC]
\begin{enumerate}
\item \textbf{Calculer l'IMC :} mesurer la masse $m$ (kg) et la taille $t$ (m) ; calculer $m/t^2$ ; arrondir au dixième ; situer dans la classification OMS.
\item \textbf{Estimer l'apport énergétique :} pour chaque aliment consommé, multiplier la quantité (en g ou en portions) par la valeur énergétique pour 100 g (ou par portion) ; sommer par repas puis sur la journée.
\item \textbf{Analyser la qualité nutritionnelle :} repérer les sources de sucres rapides (boissons sucrées, confiseries), de lipides (fritures, huiles ajoutées), l'absence de fruits/légumes frais, le volume d'eau bu.
\item \textbf{Comparer aux recommandations :} confronter l'apport total aux besoins (ici \SI{2500}{kcal} max) ; confronter l'activité physique aux 60 min/jour OMS ; confronter l'hydratation à 1,5 L/jour.
\item \textbf{Élaborer le programme :} proposer des menus variés, locaux, riches en fibres et protéines maigres, pauvres en sucres ajoutés et graisses saturées ; planifier 60 min d'activité physique quotidienne fractionnables ; viser un déficit de 300--500 kcal/jour par l'alimentation + 200--300 kcal/jour par l'activité pour perdre 0,5 kg/semaine.
\end{enumerate}
\end{methode}

\courssub{Démarche et corrigé commenté}

\begin{exoresolu}[Coach santé pour Amina : résolution complète]
Énoncé : à partir de la fiche d'Amina (1,58 m ; 68 kg ; journal alimentaire ; 0 min de sport ; 5 h d'écran), établir le diagnostic et élaborer un programme de prévention.

\tcblower

\textbf{Étape 1 — Calcul et interprétation de l'IMC.}

\[
IMC = \frac{m}{t^2} = \frac{68}{(1,58)^2} = \frac{68}{2,4964} \approx 27,2 \ \text{kg/m}^2.
\]

Or $25 \leq 27,2 < 30$ : Amina est en \cle{surpoids} selon la classification de l'OMS. Si rien ne change, elle s'expose notamment au \cle{diabète de type 2} (résistance à l'insuline favorisée par l'adiposité viscérale et l'inactivité) et aux \cle{maladies cardiovasculaires} (hypertension artérielle, athérosclérose, risque d'infarctus du myocarde ou d'accident vasculaire cérébral), ainsi qu'à l'aggravation vers l'obésité.

\begin{attention}[Piège fréquent]
Ne pas oublier d'élever la taille au carré : $\frac{68}{1,58} \approx 43$ serait une erreur grave de calcul. On calcule d'abord $1,58 \times 1,58 = 2,4964$, puis on divise. L'unité de l'IMC est kg/m$^2$ (souvent omise à l'oral, mais exigée à l'écrit).
\end{attention}

\textbf{Étape 2 — Estimation de la ration quotidienne.}

\begin{center}
\rowcolors{1}{popcream}{white}
\begin{tabularx}{\linewidth}{|X|c|}
\hline
\rowcolor{popblueL} \textbf{Aliment consommé} & \textbf{Énergie (kcal)} \\
\hline
Beignets haricots : $200 \times 3,5$ & 700 \\
\hline
Soda du matin (33 cL) & 140 \\
\hline
Riz blanc cuit : $300 \times 1,3$ & 390 \\
\hline
Huile de palme : $3 \times 90$ & 270 \\
\hline
Soda du midi (33 cL) & 140 \\
\hline
Fufu maïs : $250 \times 1,6$ & 400 \\
\hline
Eru (avec huile) : $150 \times 1,8$ & 270 \\
\hline
\textbf{Total de la journée} & \textbf{2 310} \\
\hline
\end{tabularx}
\end{center}

Le total (environ \SI{2310}{kcal}) reste sous le seuil de \SI{2500}{kcal} préconisé par l'infirmier, mais il est \textbf{trompeur} : la ration est de \cle{mauvaise qualité nutritionnelle}. Les \cle{excès} portent sur les sucres rapides (2 sodas = 280 kcal « vides », sans vitamines, ni fibres, ni sels minéraux) et les lipides (beignets frits + 3 cuillères à soupe d'huile de palme = environ 60 g de matières grasses, soit plus de 540 kcal provenant presque exclusivement de graisses saturées). Les \cle{carences} concernent les fruits (aucun), les légumes frais (l'eru est cuit longuement dans l'huile, ce qui dégrade les vitamines thermosensibles, et il n'y a aucun autre légume), l'eau (environ 0,5 L au lieu de 1,5 L minimum) et les protéines maigres (pas de poisson, viande maigre, œufs, légumineuses non frites). De plus, sans aucune dépense physique, même 2 300 kcal dépassent les besoins réels d'Amina (métabolisme de base + activités sédentaires $\approx$ 1 800--1 900 kcal), d'où une prise de masse progressive.

\textbf{Étape 3 — Diagnostic du mode de vie.}

Activité physique : 0 min/j au lieu de 60 min/j recommandées ; temps d'écran loisir : 5 h/j, ce qui favorise la sédentarité, le grignotage et les troubles du sommeil ; hydratation insuffisante (0,5 L/j). Le triptyque « excès de sucres rapides et lipides + carences en fibres/vitamines/eau + sédentarité marquée » explique le surpoids installé.

\textbf{Étape 4 — Programme hebdomadaire proposé.}

\emph{Menus types équilibrés (aliments du marché de Bafoussam, portions adaptées) :}
\begin{itemize}
\item \textbf{Petit-déjeuner (environ 500 kcal) :} bouillie de maïs peu sucrée (150 g) + 1 banane plantain bouillie (100 g) + 1 verre d'eau (25 cL) \textbf{OU} pain complet (100 g) + omelette (1 œuf + tomate/oignon) + 1 orange + eau.
\item \textbf{Déjeuner (environ 700 kcal) :} riz blanc cuit (200 g) + ndolé aux arachides (150 g, huile réduite à 1 c. à soupe) + 1 fruit de saison (mangue, ananas, papaye) ; eau à volonté.
\item \textbf{Dîner (environ 500 kcal) :} banane plantain bouillie (200 g) + légumes feuilles (folong ou kpwem, 150 g, peu huilés, 1 c. à soupe) + poisson braisé ou grillé (100 g) ; eau.
\end{itemize}
\textbf{Principes appliqués :} suppression des sodas (remplacés par de l'eau ou jus de fruits maison sans sucre ajouté) ; friture limitée à une fois par semaine maximum ; huile réduite à 1 cuillère à soupe par repas cuisiné ; portion de féculents modérée ; protéines animales maigres (poisson, volaille sans peau) ou végétales (haricots, arachides) quotidiennes ; fruits et légumes à chaque repas principal.

\emph{Planning d'activité physique (60 min/j minimum, fractionnables) :}
\begin{itemize}
\item \textbf{Lundi, mardi, jeudi, vendredi :} marche rapide domicile--lycée aller-retour (2 $\times$ 15 min = 30 min) + 20 min de corde à sauter ou danse active le soir + 10 min d'étirements/renforcement (gainage, squats).
\item \textbf{Mercredi :} 60 min de sport collectif (football, handball, basketball) au lycée ou au quartier, ou travaux champêtres actifs (bêchage, sarclage).
\item \textbf{Samedi :} 60 min de randonnée, vélo, ou aide aux travaux ménagers vigoureux (lessive à la main, nettoyage cour).
\item \textbf{Dimanche :} repos actif (promenade 30 min) + étirements ; révision du journal alimentaire de la semaine.
\end{itemize}
Réduction du temps d'écran loisir à 2 h maximum par jour.

\textbf{Étape 5 — Justification de la perte de 0,5 kg/semaine.}

Le nouveau programme ramène la ration à environ \SI{1700}{}--\SI{1900}{kcal/j} (déficit modéré de 400--600 kcal/j par rapport aux 2 310 kcal actuels, et de 300--500 kcal/j par rapport aux besoins de maintien \SI{2200}{kcal}). L'activité physique ajoutée (60 min modérées) augmente la dépense énergétique journalière d'environ 200--300 kcal. Le déficit énergétique quotidien total est donc de l'ordre de 600--800 kcal. Sur une semaine, le déficit cumulé est de 4 200--5 600 kcal. Sachant que 1 kg de tissu adipeux correspond à environ 7 700 kcal, un déficit de 3 850 kcal correspond à 0,5 kg. Le programme vise donc une perte de \textbf{0,5 kg par semaine}, rythme sûr, durable, sans suppression de repas, sans carence, et compatible avec la croissance et la scolarité d'une adolescente.

\textbf{Étape 6 — Fiche conseil validée (restitution collective).}

\begin{itemize}
\item Boire de l'eau au lieu des sodas ; viser au moins 1,5 L par jour (une gourde de 50 cl $\times$ 3).
\item Consommer au moins 1 fruit et des légumes à chaque jour ; limiter fritures et huile de cuisson (1 c. à soupe max/repas).
\item Pratiquer 60 minutes d'activité physique par jour (marche, sport, danse, travaux) ; réduire le temps d'écran loisir à moins de 2 h.
\item Objectif réaliste : perdre 0,5 kg par semaine, avec suivi mensuel de l'IMC et de la taille de taille.
\item Ne sauter aucun repas ; fractionner si besoin (collation saine : fruit, arachides grillées non salées).
\end{itemize}
\end{exoresolu}

\courssub{Bilan de l'activité}

Cette activité t'a permis de \cle{mobiliser ensemble} les trois savoir-faire de la leçon sur un cas authentique. D'abord, le calcul de l'IMC ($27,2$ kg/m$^2$) a transformé une impression subjective (« Amina est un peu forte ») en un \cle{diagnostic objectif} : un surpoids qui expose au diabète de type 2 et aux maladies cardiovasculaires. Ensuite, l'évaluation de la ration a montré un point essentiel : \textbf{une ration peut rester sous le seuil calorique tout en étant déséquilibrée} — la qualité des aliments (sucres rapides, lipides saturés, fibres, vitamines, eau) compte autant que la quantité. Enfin, l'élaboration du programme a mis en évidence que la prévention ne consiste pas à supprimer des repas ni à imposer des aliments introuvables ou coûteux, mais à \cle{rééquilibrer} l'alimentation avec les produits locaux accessibles (maïs, haricots, ndolé, banane plantain, poisson, fruits de saison) et à \cle{réintroduire} une activité physique quotidienne compatible avec la vie d'un lycéen camerounais (marche active, corde à sauter, sport collectif le mercredi). Tu as aussi développé des savoir-être indispensables : écoute active, absence de jugement moralisateur, argumentation scientifique rigoureuse et esprit d'équipe — les qualités d'un véritable acteur de santé communautaire au service du bien-être de ses pairs.

\newpage

\coursec{Méthodes et exercices résolus}

\coursec{Lutte contre la mauvaise alimentation et l'inactivité physique}

\courssub{Équilibre alimentaire et besoins nutritionnels}

\begin{definition}[Équilibre alimentaire]
L'\cle{équilibre alimentaire} est un régime qui couvre, sans excès ni carence, les besoins de l'organisme en énergie et en nutriments (macronutriments : glucides, lipides, protides ; micronutriments : vitamines, sels minéraux, oligo-éléments ; eau et fibres). Il repose sur la diversité des aliments consommés, leur répartition au cours de la journée (3 à 4 prises) et l'adaptation des quantités aux dépenses énergétiques.
\end{definition}

\begin{propriete}[Les sept groupes d'aliments et leurs apports principaux]
Une ration équilibrée doit puiser quotidiennement dans les \cle{sept groupes d'aliments} :
\begin{center}
\begin{tabularx}{\linewidth}{|l|X|l|}\hline
\rowcolor{popblueL} \textbf{Groupe} & \textbf{Aliments représentatifs (contexte camerounais)} & \textbf{Apports majeurs} \\ \hline
1. Céréales, tubercules, dérivés & Maïs, riz, mil, sorgho, manioc (bâton, miondo, foufou), plantain, pain, pâtes & Glucides complexes (amidon), énergie, fibres, vitamines B \\ \hline
2. Légumes & Feuilles de ndolé, gombo, folong, épinards, tomate, carotte, chou & Fibres, vitamines (A, C, K, B9), sels minéraux (K, Mg), eau \\ \hline
3. Fruits & Mangue, banane, papaye, ananas, agrumes, pastèque, safou & Fibres, vitamines (C, B9, provitamine A), eau, sucres simples \\ \hline
4. Viandes, poissons, œufs & Poulet, bœuf, poisson (capitaine, carpe), crevettes, œufs, viande de brousse & Protides de haute valeur biologique, fer héminique, zinc, vitamine B12 \\ \hline
5. Lait et produits laitiers & Lait caillé, yaourt, fromage, lait en poudre & Calcium, phosphore, protides, vitamines A, D, B2, B12 \\ \hline
6. Corps gras & Huile rouge, huile d'arachide, huile de coton, beurre, margarine, avocat & Lipides (acides gras), vitamines A, D, E, K, énergie dense \\ \hline
7. Eau & Eau de source, eau traitée, bouillons, tisanes non sucrées & Hydratation, apport minéral, transport des nutriments \\ \hline
\end{tabularx}
\end{center}
\legende{Les 7 groupes d'aliments : la diversité quotidienne est la clé de la couverture des besoins en micronutriments.}
\end{propriete}

\begin{propriete}[Besoins énergétiques journaliers (Apports Nutritionnels Conseillés -- ANC)]
Les besoins varient selon l'\cle{âge}, le \cle{sexe}, la \cle{stature}, la \cle{composition corporelle} et surtout le \cle{niveau d'activité physique (NAP)}. Valeurs de référence pour la population camerounaise (inspirées OMS/FAO) :
\begin{center}
\begin{tabularx}{\linewidth}{|l|X|X|}\hline
\rowcolor{popblueL} \textbf{Population} & \textbf{NAP Faible (Sédentaire)} & \textbf{NAP Modéré à Élevé (Actif)} \\ \hline
Adolescent garçon (13--19 ans) & \SIrange{2500}{2800}{kcal/j} & \SIrange{3000}{3500}{kcal/j} \\ \hline
Adolescente fille (13--19 ans) & \SIrange{2200}{2400}{kcal/j} & \SIrange{2600}{3000}{kcal/j} \\ \hline
Homme adulte (18--60 ans) & \SIrange{2100}{2500}{kcal/j} & \SIrange{2600}{3200}{kcal/j} \\ \hline
Femme adulte (18--60 ans) & \SIrange{1800}{2200}{kcal/j} & \SIrange{2200}{2600}{kcal/j} \\ \hline
\end{tabularx}
\end{center}
\legende{Le Métabolisme de Base (MB) représente environ 60--70\% de la dépense énergétique totale. L'effet thermique des aliments ~10\%. L'activité physique est la composante la plus variable (15--30\% voire plus).}
\end{propriete}

\begin{propriete}[Répartition qualitative de l'apport énergétique total (AET)]
Pour une alimentation équilibrée, les macronutriments doivent contribuer à l'AET selon les proportions suivantes (repères ANC) :
\begin{center}
\begin{tabularx}{\linewidth}{|l|X|X|}\hline
\rowcolor{popblueL} \textbf{Macronutriment} & \textbf{\% de l'AET} & \textbf{Rôle principal} \\ \hline
Glucides (dont sucres ajoutés $< 10\%$) & 50 -- 55 \% & Carburant principal (cerveau, muscles, globules rouges) \\ \hline
Lipides (AG saturés $< 12\%$, AG trans $< 1\%$) & 30 -- 35 \% & Structure membranaire, réserves, vitamines liposolubles \\ \hline
Protides (animaux + végétaux) & 12 -- 15 \% & Structure, enzymes, transport, immunité, renouvellement \\ \hline
\end{tabularx}
\end{center}
\legende{Les fibres (25--30 g/j) et l'eau (1,5--2 L/j) sont indispensables bien qu'apportant peu ou pas d'énergie.}
\end{propriete}

\begin{methode}[Calculer la valeur énergétique d'un aliment ou d'un repas à partir de sa composition]
\begin{enumerate}
\item \textbf{Identifier} les masses (en grammes) de glucides ($m_G$), protides ($m_P$) et lipides ($m_L$) fournis par l'aliment ou le repas (étiquette nutritionnelle ou table de composition Ciqual/ANSES).
\item \textbf{Appliquer les équivalents énergétiques} (coefficients d'Atwater) :
\[ E_G = m_G \times 4\ \mathrm{kcal/g} \quad ; \quad E_P = m_P \times 4\ \mathrm{kcal/g} \quad ; \quad E_L = m_L \times 9\ \mathrm{kcal/g} \]
\item \textbf{Calculer l'énergie totale} : $E_{\text{totale}} = E_G + E_P + E_L$ (en kcal).
\item \textbf{Calculer la répartition en pourcentage} pour chaque macronutriment :
\[ \%_G = \frac{E_G}{E_{\text{totale}}} \times 100 \quad ; \quad \%_P = \frac{E_P}{E_{\text{totale}}} \times 100 \quad ; \quad \%_L = \frac{E_L}{E_{\text{totale}}} \times 100 \]
\item \textbf{Comparer} les pourcentages obtenus aux normes de l'équilibre (G 50--55\%, L 30--35\%, P 12--15\%).
\end{enumerate}
\end{methode}

\begin{exoresolu}[Valeur énergétique et équilibre d'un repas de cantine]
Un repas servi au réfectoire d'un lycée de Douala apporte : 85 g de glucides, 28 g de protides, 22 g de lipides.
\begin{enumerate}
\item Calculez l'apport énergétique total de ce repas.
\item Déterminez la répartition énergétique en \% et dites si ce repas est équilibré au regard des normes.
\end{enumerate}
\tcblower
\textbf{1. Apports énergétiques par macronutriment.}
\begin{align*}
E_G &= 85 \times 4 = 340\ \mathrm{kcal} \\
E_P &= 28 \times 4 = 112\ \mathrm{kcal} \\
E_L &= 22 \times 9 = 198\ \mathrm{kcal} \\
E_{\text{totale}} &= 340 + 112 + 198 = \SI{650}{kcal}
\end{align*}

\textbf{2. Répartition en pourcentages.}
\begin{align*}
\%_G &= \frac{340}{650} \times 100 \approx 52,3\% \\
\%_P &= \frac{112}{650} \times 100 \approx 17,2\% \\
\%_L &= \frac{198}{650} \times 100 \approx 30,5\%
\end{align*}

\textbf{3. Interprétation.}
\begin{itemize}
\item Glucides : 52,3\% $\in [50; 55]$ $\rightarrow$ \textbf{Conforme}.
\item Lipides : 30,5\% $\in [30; 35]$ $\rightarrow$ \textbf{Conforme}.
\item Protides : 17,2\% $> 15\%$ $\rightarrow$ \textbf{Légèrement excédentaire} (fréquent si portion de viande/poisson importante).
\end{itemize}
\textbf{Conclusion} : Ce repas est globalement bien équilibré sur le plan énergétique, avec une répartition lipidique et glucidique correcte. L'apport protidique est un peu élevé mais acceptable ponctuellement.
\end{exoresolu}

\begin{methode}[Évaluer une ration alimentaire journalière par rapport aux besoins]
\begin{enumerate}
\item \textbf{Déterminer les besoins énergétiques} du sujet (âge, sexe, NAP) à l'aide des ANC.
\item \textbf{Lister tous les aliments} consommés dans la journée (pesées ou estimations de portions usuelles).
\item \textbf{Calculer l'apport énergétique total (AET)} de la journée (somme des repas + collations) via les tables de composition ou les coefficients d'Atwater.
\item \textbf{Analyser la couverture qualitative} : vérifier la présence des 7 groupes d'aliments, la répartition G/L/P, l'apport en fibres, eau, sel, sucres ajoutés.
\item \textbf{Conclure} : \emph{Quantitativement} (AET vs Besoins : déficitaire / équilibré / excédentaire) ET \emph{Qualitativement} (groupes manquants, excès de gras/sucre/sel, fibres). Proposer des corrections concrètes.
\end{enumerate}
\end{methode}

\begin{exoresolu}[Évaluation de la ration de Sandrine, 16 ans, élève de Première D]
Besoins estimés (NAP modéré) : \SI{2600}{kcal/j}. Ration déclarée :
\begin{center}
\begin{tabularx}{\linewidth}{|l|X|l|}\hline
\rowcolor{popblueL} \textbf{Repas} & \textbf{Aliments} & \textbf{Énergie estimée (kcal)} \\ \hline
Matin & Beignets-haricots + bouillie de maïs sucrée & 650 \\ \hline
Midi & Ndolé (feuilles, arachides, viande/poisson) + Miondo (manioc) & 950 \\ \hline
Soir & Riz sauté (huile, légumes) + Cuisse de poulet + 1/2 Avocat & 900 \\ \hline
\end{tabularx}
\end{center}
\begin{enumerate}
\item Calculez l'AET et comparez aux besoins.
\item Analysez la couverture des 7 groupes. Proposez deux corrections précises.
\end{enumerate}
\tcblower
\textbf{1. Bilan énergétique.}
\[ \text{AET} = 650 + 950 + 900 = \SI{2500}{kcal} \]
Écart : $2600 - 2500 = \SI{-100}{kcal}$ (déficit léger, $< 5\%$). \textbf{Quantitativement : ration quasi-équilibrée.}

\textbf{2. Bilan qualitatif (Groupes présents / Absents).}
\begin{itemize}
\item \textbf{Présents} : G1 (Maïs, Miondo, Riz), G2 (Feuilles ndolé, légumes riz), G4 (Viande/Poisson ndolé, Poulet), G6 (Huile beignets, Arachides, Huile riz, Avocat).
\item \textbf{Absents ou insuffisants} : \textbf{G3 (Fruits)} -- aucun fruit frais. \textbf{G5 (Lait/Produits laitiers)} -- aucun. \textbf{G2} : Légumes peu variés (surtout feuilles cuites), pas de crudités.
\end{itemize}

\textbf{Corrections proposées :}
\begin{enumerate}
\item Ajouter \textbf{un fruit de saison} (mangue, banane, papaye, ananas) au petit-déjeuner ou en collation ($\approx$ 80--100 kcal, vit C, fibres, eau). Comble le déficit énergétique et couvre G3.
\item Introduire \textbf{un produit laitier} (yaourt nature, lait caillé, verre de lait) en collation ou au soir ($\approx$ 100--150 kcal, Ca, protides, vit B2, B12). Couvre G5.
\end{enumerate}
\textbf{Nouvel AET} $\approx 2650$ kcal $\rightarrow$ Parfaitement adapté. Ration diversifiée, couvrantes en micronutriments.
\end{exoresolu}

\courssub{L'IMC : calcul et interprétation}

\begin{definition}[Indice de Masse Corporelle (IMC)]
L'\cle{IMC} (ou indice de Quetelet) est un indicateur simple, validé par l'OMS, pour estimer la corpulence d'un \textbf{adulte} (18--65 ans) et le risque sanitaire associé au poids. Il ne distingue pas masse grasse et masse musculaire.
\[ \text{IMC} = \frac{m}{T^2} \]
avec $m$ la masse en \textbf{kilogrammes (kg)} et $T$ la taille en \textbf{mètres (m)}. Unité : $\mathrm{kg/m^2}$.
\end{definition}

\begin{attention}[Conditions d'utilisation et limites de l'IMC]
\begin{itemize}
\item \textbf{Valide pour} : Adultes 18--65 ans, non sportifs de haut niveau, non femmes enceintes/allaitantes, non personnes âgées $> 65$ ans (perte musculaire), non amputés.
\item \textbf{Ne s'applique PAS directement aux enfants/adolescents} : utiliser les \cle{courbes de référence IMC-âge} (carnets de santé) avec percentiles (surpoids $> P97$, obésité $> P99,5$).
\item \textbf{Ne distingue pas} : Un rugbyman de 100 kg pour 1,80 m (IMC 30,8) n'est pas obèse (masse musculaire). Compléter par le \cle{tour de taille} (risque métabolique si $> 94$ cm H / $> 80$ cm F) ou l'impédancemétrie.
\end{itemize}
\end{attention}

\begin{methode}[Calculer et interpréter l'IMC d'un adulte]
\begin{enumerate}
\item \textbf{Relever} : Masse $m$ (kg) et Taille $T$ (m). \textbf{Convertir impérativement cm $\rightarrow$ m} (ex: 172 cm = 1,72 m).
\item \textbf{Calculer} : $\text{IMC} = m / T^2$.
\item \textbf{Arrondir} à \textbf{un chiffre après la virgule} et noter l'unité $\mathrm{kg/m^2}$.
\item \textbf{Classer} selon la nomenclature OMS 2004 :
\begin{center}
\begin{tabularx}{\linewidth}{|l|X|}\hline
\rowcolor{popblueL} \textbf{IMC (kg/m$^2$)} & \textbf{Classification OMS / Risque} \\ \hline
$< 16,5$ & Dénutrition / Famine \\ \hline
$16,5 \leq \text{IMC} < 18,5$ & Maigreur / Insuffisance pondérale \\ \hline
$18,5 \leq \text{IMC} < 25$ & \textbf{Corpulence normale} (Référence) \\ \hline
$25 \leq \text{IMC} < 30$ & \textbf{Surpoids} (Risque accru) \\ \hline
$30 \leq \text{IMC} < 35$ & \textbf{Obésité classe I} (Modérée, Risque élevé) \\ \hline
$35 \leq \text{IMC} < 40$ & \textbf{Obésité classe II} (Sévère, Risque très élevé) \\ \hline
$\geq 40$ & \textbf{Obésité classe III} (Morbide / Massive, Risque extrême) \\ \hline
\end{tabularx}
\end{center}
\item \textbf{Conclure} sur les risques : Diabète type 2, HTA, Dyslipidémie, Maladies cardiovasculaires (IDM, AVC), Apnée du sommeil, Ostéoarthrite, Certains cancers.
\end{enumerate}
\end{methode}

\begin{popfigure}
\begin{tikzpicture}[scale=0.9, every node/.style={font=\small}]
% Axes
\draw[->, thick] (0,0) -- (10,0) node[right] {IMC (kg/m$^2$)};
\draw[->, thick] (0,0) -- (0,3.5) node[above] {Risque relatif};
% Zones colorées
\fill[popgreenL] (1.85,0) rectangle (5,3);
\fill[popgoldL] (5,0) rectangle (6,3);
\fill[poporangeL] (6,0) rectangle (7,3);
\fill[poppinkL] (7,0) rectangle (8,3);
\fill[poppink] (8,0) rectangle (9.5,3);
\fill[popblueL] (0.5,0) rectangle (1.85,3);
\fill[popblueL] (9.5,0) rectangle (10,3);
% Graduations et labels
\foreach \x/\label in {1.85/18.5, 5/25, 6/30, 7/35, 8/40} {
    \draw (\x,0.1) -- (\x,-0.1) node[below] {\label};
}
% Labels zones
\node[align=center, text=popdark] at (1.2,2.5) {Insuffisance\\pondérale};
\node[align=center, text=popdark] at (3.4,2.5) {\textbf{Normale}\\(Risque ref.)};
\node[align=center, text=popdark] at (5.5,2.5) {Surpoids\\(Risque +)};
\node[align=center, text=popdark] at (6.5,2.5) {Obésité I\\(Risque ++)};
\node[align=center, text=popdark] at (7.5,2.5) {Obésité II\\(Risque +++)};
\node[align=center, text=white] at (8.75,2.5) {Obésité III\\(Risque ++++)};
% Courbe risque schématique (U inversé)
\draw[popcurve, very thick, domain=1.85:9.5, samples=100] plot (\x, {0.5 + 0.03*(\x-3.4)^2});
\end{tikzpicture}
\legende{Relation en U entre IMC et mortalité/risque morbide : risque minimal pour IMC $\in [20; 25]$.}
\end{popfigure}

\begin{exoresolu}[Calcul et interprétation : Monsieur Ngoa, 42 ans, Yaoundé]
Données : $m = \SI{95}{kg}$, $T = \SI{1,75}{m}$.
\begin{enumerate}
\item Calculez son IMC.
\item Interprétez et citez 3 risques majeurs.
\end{enumerate}
\tcblower
\textbf{1. Calcul.}
\[ \text{IMC} = \frac{95}{(1,75)^2} = \frac{95}{3,0625} \approx 31,0\ \mathrm{kg/m^2} \]

\textbf{2. Interprétation.}
$30 \leq 31,0 < 35$ $\rightarrow$ \textbf{Obésité classe I (Modérée)}.

\textbf{3. Risques sanitaires majeurs (étiologie multifactorielle, l'obésité est facteur favorisant) :}
\begin{enumerate}
\item \cle{Diabète de type 2} (insulinorésistance).
\item \cle{Hypertension artérielle} (HTA) et \cle{Maladies cardiovasculaires} (Infarctus du myocarde, AVC ischémique).
\item \cle{Dyslipidémie} (hypertriglycéridémie, HDL bas) $\rightarrow$ \cle{Athérosclérose} accélérée.
\end{enumerate}
\textbf{Prise en charge} : Rééquilibrage alimentaire (déficit 500 kcal/j), Activité physique adaptée (150 min/semaine endurance), Suivi médical (glycémie, lipidogramme, tension).
\end{exoresolu}

\courssub{Conséquences de la sédentarité et de la mauvaise alimentation}

\begin{experience}[Modélisation de la balance énergétique]
\begin{center}
\begin{tikzpicture}[node distance=1.5cm, >=Stealth, auto, thick]
\node[draw, rectangle, rounded corners, fill=popblueL, minimum width=3cm, align=center] (apports){Apports Énergétiques\\ (Alimentation)};
\node[draw, rectangle, rounded corners, fill=poporangeL, minimum width=3cm, right=of apports, align=center] (stock){Stockage \\ (Tissu adipeux)};
\node[draw, rectangle, rounded corners, fill=popgreenL, minimum width=3cm, below=of stock, align=center] (depenses){Dépenses Énergétiques\\ (MB + Thermique + AP)};
\draw[->, popblue] (apports) -- node[above] {Entrées} (stock);
\draw[->, popgreen] (stock) -- node[right] {Mobilisation} (depenses);
\draw[->, poporange, dashed] (apports) to[bend left=30] node[above] {Si Apports > Dépenses} (stock);
\draw[->, poporange, dashed] (depenses) to[bend right=30] node[below] {Si Dépenses > Apports} (stock);
\end{tikzpicture}
\end{center}
\legende{La balance énergétique : l'excès d'apports (sédentarité + alimentation riche) remplit le stock adipeux. L'activité physique augmente les dépenses et mobilise les stocks.}
\end{experience}

\begin{propriete}[Chaîne causale : Sédentarité + Surcharge alimentaire $\rightarrow$ Maladies non transmissibles (MNT)]
\begin{enumerate}
\item \textbf{Déséquilibre énergétique} : Sédentarité $\downarrow$ Dépenses ; Alimentation riche (gras, sucres) $\uparrow$ Apports $\rightarrow$ \cle{Balance énergétique positive}.
\item \textbf{Stockage adipocytaire} : Excès d'énergie $\rightarrow$ Lipogenèse hépatique $\rightarrow$ Stockage sous forme de \cle{triglycérides} dans les \cle{adipocytes} (hypertrophie puis hyperplasie) $\rightarrow$ \cle{Surpoids} puis \cle{Obésité} (IMC $\geq 30$).
\item \textbf{Insulinorésistance (Mécanisme central du DT2)} :
\begin{itemize}
\item L'adipocyte hypertrophié sécrète des \cle{adipokines pro-inflammatoires} (TNF-$\alpha$, IL-6, résistine) et moins d'adiponectine.
\item Ces facteurs perturbent la \cle{voie de signalisation de l'insuline} (récepteur tyrosine kinase $\rightarrow$ IRS-1 $\rightarrow$ PI3K $\rightarrow$ Akt $\rightarrow$ translocation \cle{GLUT4}) dans le muscle, le foie, le tissu adipeux.
\item $\rightarrow$ \cle{Captation du glucose} diminuée $\rightarrow$ \cle{Hyperglycémie} post-prandiale puis à jeun.
\item Pancréas $\rightarrow$ \cle{Hyperinsulinémie} compensatoire $\rightarrow$ Épuisement des \cle{cellules $\beta$ pancréatiques} $\rightarrow$ \cle{Diabète de type 2} (Glycémie à jeun $\geq 1,26$ g/L ou 7 mmol/L).
\end{itemize}
\item \textbf{Athérosclérose (Mécanisme central des MCV)} :
\begin{itemize}
\item Excès d'apports lipidiques (AG saturés, cholestérol) + Insulinorésistance $\rightarrow$ \cle{Dyslipidémie athérogène} : LDL-c $\uparrow$, Triglycérides $\uparrow$, HDL-c $\downarrow$, LDL dense (petit, oxydable).
\item \cle{LDL oxydé} capté par macrophages $\rightarrow$ \cle{Cellules spumeuses} $\rightarrow$ \cle{Stries graisseuses} $\rightarrow$ \cle{Plaque d'athérome} (lipides, cellules inflammatoires, matrice fibreuse).
\item \cle{Complications} : Rétrécissement lumière (ischémie), Rupture plaque + Thrombose $\rightarrow$ \cle{Infarctus du myocarde} (coronaire) ou \cle{Accident Vasculaire Cérébral} (carotide/cérébrale).
\item Hypertension artérielle (raideur artérielle, hyperinsulinémie, activation système rénine-angiotensine) aggrave le risque.
\end{itemize}
\end{enumerate}
\end{propriete}

\begin{popfigure}
\begin{tikzpicture}[scale=0.85, every node/.style={font=\small}]
% Insulinorésistance
\begin{scope}[xshift=0cm]
\node[draw, rectangle, rounded corners, fill=poppinkL, minimum width=5cm, minimum height=3.5cm] (box1) {};
\node[align=center, font=\bfseries] at (0, 1.5) {Insulinorésistance (Muscle / Foie / Adipocyte)};
% Normal
\node[draw, circle, fill=popgreenL, minimum size=0.8cm] (ins) at (-2, 0.5) {Insuline};
\node[draw, rectangle, rounded corners, fill=white, minimum width=1.5cm] (rec) at (0, 0.5) {Récepteur Insuline};
\node[draw, diamond, fill=popblueL, minimum size=1cm] (glut) at (2, 0.5) {GLUT4};
\node[circle, fill=poporange, minimum size=0.5cm] (gluc) at (3.5, 0.5) {Glucose};
\draw[->, thick, popgreen] (ins) -- (rec);
\draw[->, thick, popgreen] (rec) -- (glut);
\draw[->, thick, popgreen] (glut) -- (gluc);
\node[align=left, font=\tiny, text=popgreen] at (0.75, 1.2) {Voie PI3K/Akt\\Translocation GLUT4};
% Resistant
\node[draw, circle, fill=poppink, minimum size=0.8cm] (ins2) at (-2, -0.8) {Insuline};
\node[draw, rectangle, rounded corners, fill=white, minimum width=1.5cm] (rec2) at (0, -0.8) {Récepteur Insuline};
\node[draw, diamond, fill=popgray, minimum size=1cm] (glut2) at (2, -0.8) {GLUT4 (intracell.)};
\node[circle, fill=poporange, minimum size=0.5cm] (gluc2) at (3.5, -0.8) {Glucose (sang)};
\draw[->, thick, poppink] (ins2) -- (rec2);
\draw[-, thick, dashed, poppink] (rec2) -- (glut2);
\draw[-, thick, dashed, poppink] (glut2) -- (gluc2);
\node[align=left, font=\tiny, text=poppink] at (0.75, -1.5) {Signal bloqué\\GLUT4 non translocé};
\node[align=left, font=\tiny, text=popdark] at (-2, -1.8) {TNF-$\alpha$, IL-6, \\Acides gras libres};
\end{scope}
% Athérosclérose
\begin{scope}[xshift=7cm]
\node[draw, rectangle, rounded corners, fill=poppinkL, minimum width=5cm, minimum height=3.5cm] (box2) {};
\node[align=center, font=\bfseries] at (7, 1.5) {Athérosclérose (Artère coronaire / carotide)};
% Healthy
\draw[popgreen, line width=3mm] (-2, 0.5) -- (2, 0.5);
\draw[popgreen, line width=1mm] (-2, 0.5) -- (2, 0.5);
\node[align=center, font=\tiny] at (0, 1.1) {Endothelium sain\\Lumière large};
% Plaque
\draw[poppink, line width=3mm] (-2, -0.8) -- (2, -0.8);
\draw[poppink, line width=1mm] (-2, -0.8) -- (2, -0.8);
\fill[poporange] (-0.5, -1.1) rectangle (0.5, -0.5);
\node[align=center, font=\tiny, text=white] at (0, -0.8) {Plaque d'athérome};
\node[align=center, font=\tiny] at (0, -1.5) {LDL oxydé, Cellules spumeuses,\\Matrice fibreuse, Calcification};
\draw[->, thick, poppink] (1.5, -0.8) -- (1.5, -1.8) node[below] {Rupture $\rightarrow$ Thrombus};
\end{scope}
\end{tikzpicture}
\legende{Gauche : Insulinorésistance (blocage translocation GLUT4). Droite : Plaque d'athérome rétrécissant la lumière artérielle.}
\end{popfigure}

\begin{methode}[Expliquer le lien « Mode de vie $\rightarrow$ Maladie » (Rédaction attendue au Probatoire)]
Pour répondre à « Expliquez comment la sédentarité favorise le diabète de type 2 » :
\begin{enumerate}
\item \textbf{Énoncer le point de départ} : Sédentarité $\downarrow$ Dépenses + Alimentation déséquilibrée $\uparrow$ Apports $\rightarrow$ Balance positive.
\item \textbf{Décrire le stockage} : Excès $\rightarrow$ Triglycérides dans adipocytes $\rightarrow$ Obésité abdominale (adiposité viscérale).
\item \textbf{Expliquer le mécanisme cellulaire} : Adiposité viscérale $\rightarrow$ Sécrétion cytokines pro-inflammatoires (TNF-$\alpha$, IL-6) + Acides Gras Libres $\rightarrow$ \textbf{Inhibition de la voie de signalisation de l'insuline} (phosphorylation IRS-1 sur sérine) $\rightarrow$ \textbf{Défaut de translocation des transporteurs GLUT4} $\rightarrow$ \textbf{Diminution de l'entrée du glucose} dans myocytes/adipocytes/hepatocytes.
\item \textbf{Nommer la conséquence} : \textbf{Hyperglycémie chronique} $\rightarrow$ \textbf{Diabète de type 2}.
\item \textbf{Citer la prévention} : Activité physique $\rightarrow$ Voie AMPk $\rightarrow$ Translocation GLUT4 \emph{indépendante de l'insuline} $\rightarrow$ Amélioration sensibilité.
\end{enumerate}
\end{methode}

\begin{exoresolu}[Sédentarité, obésité et diabète de type 2 -- Contexte Hôpital Bafoussam]
\textit{Question : Expliquez le lien physiologique entre sédentarité, obésité et apparition du diabète de type 2.}
\tcblower
\textbf{Raisonnement causal structuré :}
\begin{enumerate}
\item La \cle{sédentarité} réduit la dépense énergétique d'activité (DEA). Sans diminution corrélative des apports, la \cle{balance énergétique} devient positive.
\item L'énergie excédentaire est convertie en acides gras puis stockée sous forme de \cle{triglycérides} dans les \cle{adipocytes} du tissu adipeux (hypertrophie cellulaire). L'accumulation de masse grasse, notamment \cle{viscérale}, définit l'\cle{obésité} (IMC $\geq 30\ \mathrm{kg/m^2}$).
\item L'adipocyte hypertrophié devient dysfonctionnel : il sécrète des \cle{adipokines pro-inflammatoires} (TNF-$\alpha$, IL-6, résistine) et libère massivement des \cle{acides gras libres} (AGL) dans la veine porte vers le foie.
\item Ces facteurs induisent une \cle{insulinorésistance} périphérique (muscle, tissu adipeux) et hépatique : la cascade de signalisation de l'\cle{insuline} (Récepteur $\rightarrow$ IRS-1 $\rightarrow$ PI3K $\rightarrow$ Akt) est inhibée $\rightarrow$ les vésicules \cle{GLUT4} ne fusionnent pas avec la membrane plasmique.
\item Le glucose sanguin ne pénètre plus efficacement dans les cellules insulinodépendantes $\rightarrow$ \cle{hyperglycémie} post-prandiale puis à jeun. L'\cle{hyperinsulinémie} compensatoire masque initialement le défaut, puis l'épuisement des \cle{cellules $\beta$ pancréatiques} survient.
\item Seuil diagnostique franchi : Glycémie à jeun $\geq \SI{1,26}{g/L}$ (ou $\geq \SI{7}{mmol/L}$) $\rightarrow$ \cle{Diabète de type 2}.
\end{enumerate}
\textbf{Conclusion} : La sédentarité, par l'obésité viscérale qu'elle favorise, déclenche une inflammation chronique de bas grade et une lipotoxicité qui bloquent l'action de l'insuline (insulinorésistance), aboutissant au diabète de type 2. L'activité physique régulière restaure la sensibilité à l'insuline via la voie AMPK/GLUT4.
\end{exoresolu}

\courssub{Recommandations : Alimentation équilibrée}

\begin{aretenir}[Repères alimentaires pour la population camerounaise (Adaptation PNNS / OMS)]
\begin{itemize}
\item \textbf{Fruits \& Légumes} : \textbf{$\geq$ 5 portions/jour} (400 g minimum). Varier couleurs/cuissons. Privilégier crus ou cuisson vapeur/braisée (ndolé, folong peu cuits).
\item \textbf{Céréales complètes / Tubères locaux} : À chaque repas selon appétit. Préférer \textbf{complets} (maïs grain, mil, sorgho, riz complet, manioc/foufou) aux raffinés (pain blanc, pâtes blanches, riz blanc très poli). Source fibres, Mg, vitamines B.
\item \textbf{Produits laitiers} : 2 à 3 portions/jour (Lait caillé, yaourt nature, fromage). Source Ca, Vit D, B12. Attention au lait concentré sucré (limiter).
\item \textbf{Viandes, Poissons, Œufs, Légumineuses} : 1 à 2 portions/jour. Alterner. \textbf{Poisson} (capitaine, carpe, sardines) $\geq$ 2 fois/sem (Oméga-3). \textbf{Légumineuses} (haricots, niébé, soja, arachides) : source protides végétaux + fibres + fer (associer vit C pour absorption). Limiter viandes rouges/transformées (saucisses, viande séchée type \textit{kilishi}).
\item \textbf{Corps gras} : \textbf{Varier les huiles} (Rouge = vit A/E + AG saturés modérés ; Arachide/Coton/Tournesol = Oméga-6 ; Colza/Soja/Lin = Oméga-3). \textbf{Limiter} fritures, beurre, margarines hydrogénées (AG trans). Cuissons : vapeur, braise, papillote, mijoté peu huilé.
\item \textbf{Produits sucrés / Boissons sucrées} : \textbf{Limiter drastiquement} (jus industriels, sodas, \textit{foléré} sucré, beignets, bonbons). Sucres ajoutés $< 10\%$ AET (idéal $< 5\%$). Eau = seule boisson recommandée ($\geq \SI{1,5}{L/j}$).
\item \textbf{Sel / Sodium} : $< \SI{5}{g/j}$ (1 c. à café). Limiter cubes bouillon (Maggi/Knorr), poisson séché/salé, sauces salées. Utiliser épices, herbes, ail, oignon, gingembre, piment.
\item \textbf{Alcool} : Zéro risque = Zéro consommation. Sinon $\leq$ 2 verres/j (H) / 1 verre/j (F), pas tous les jours.
\end{itemize}
\end{aretenir}

\courssub{Recommandations : Activité physique}

\begin{aretenir}[Recommandations OMS 2020 -- Activité Physique (AP) pour la santé]
\begin{center}
\begin{tabularx}{\linewidth}{|l|X|X|}\hline
\rowcolor{popblueL} \textbf{Groupe} & \textbf{Endurance (Modérée à Soutenue)} & \textbf{Renforcement Musculaire / Souplesse} \\ \hline
\textbf{Enfants / Adolescents (5--17 ans)} & \textbf{60 min/jour} (majoritairement aérobie, intensité modérée à soutenue) & \textbf{3 fois/sem} minimum (jeux, port de charges, sport) \\ \hline
\textbf{Adultes (18--64 ans)} & \textbf{150--300 min/sem} modérée \textbf{OU} \textbf{75--150 min/sem} soutenue (ou combinaison) & \textbf{2 fois/sem} minimum (grands groupes musculaires) \\ \hline
\textbf{Personnes âgées ($\geq$ 65 ans)} & Idem adultes + \textbf{Equilibre / Prévention chutes} 3 fois/sem & \textbf{2--3 fois/sem} (force, équilibre, souplesse) \\ \hline
\end{tabularx}
\end{center}
\legende{Modérée = marche rapide, vélo lent, danse, travaux champêtres légers (on peut parler pas chanter). Soutenue = footing, football, vélo rapide, nage, port charges lourdes (on ne peut pas parler). \textbf{Sédentarité} : temps assis/allongé $\neq$ inactivité. Casser la sédentarité : lever toutes les 30--60 min.}
\end{aretenir}

\begin{savaistu}[Activité physique au Cameroun : Enquête STEPS 2017--2018]
\begin{itemize}
\item \textbf{Prévalence de l'inactivité physique} : $\approx 22\%$ des adultes (Femmes 27\% > Hommes 17\%).
\item \textbf{Urbanisation} : Inactivité plus forte en ville (Douala, Yaoundé) qu'en milieu rural (travaux champêtres, marche).
\item \textbf{Transition épidémiologique} : Les MNT (Maladies Non Transmissibles) représentent $> 35\%$ des décès. L'inactivité physique est le 4ème facteur de risque de mortalité mondial.
\item \textbf{Opportunités} : Culture de la danse (makossa, bikutsi, coupés-décalés), football populaire, marche pour aller à l'école/marché, agriculture. Levier pour la prévention primaire.
\end{itemize}
\end{savaistu}

\courssub{Élaborer un programme de prévention combiné}

\begin{methode}[Construire un programme de prévention personnalisé « Alimentation + Activité Physique »]
\begin{enumerate}
\item \textbf{Diagnostic initial (Bilan) :}
\begin{itemize}
\item Anthropométrie : IMC, Tour de taille (critère métabolique : H $> 94$ cm, F $> 80$ cm).
\item Alimentaire : Rappel 24h ou 3 jours (méthode des portions) $\rightarrow$ AET, Répartition G/L/P, Groupes manquants, Excès (Sel, Sucre, AG saturés, Alcool).
\item Activité physique : Questionnaire (type GPAQ/STEPS) $\rightarrow$ NAP (METs-min/sem), Temps sédentaire.
\item Biologique (si possible) : Glycémie à jeun, Lipidogramme, Créaninémie, TAS/TAD.
\end{itemize}
\item \textbf{Objectifs SMART (Spécifiques, Mesurables, Atteignables, Réalistes, Temporels) :}
\begin{itemize}
\item Perte de masse : \textbf{0,5 à 1 kg/semaine} (déficit $\approx 500$ kcal/j). \textbf{Jamais $> 1$ kg/sem}.
\item IMC cible : $< 25$ (ou $- 5$ à $10\%$ de la masse initiale si obésité sévère).
\item AP : Atteindre recommandations OMS (150 min modérée/sem) \textbf{progressivement}.
\item Alimentation : Couverture 7 groupes, Répartition G/L/P, Eau 1,5--2 L, Sel $< 5$ g.
\end{itemize}
\item \textbf{Plan d'action Alimentaire (Exemple type) :}
\begin{itemize}
\item 3 repas structurés + 1 collation (fruit/lait) si besoin. \textbf{Pas de grignotage}.
\item Assiette type : 1/2 Légumes, 1/4 Féculents (complets), 1/4 Protéines (poisson/légumineuses/viande maigre). 1 Fruit. 1 Laitage. Eau.
\item Cuissons : Vapeur, Braise, Four, Mijoté (peu d'huile ajoutée en fin de cuisson).
\item Suppression : Boissons sucrées, Fritures quotidiennes, Charcuteries, Cubes bouillon excessifs.
\end{itemize}
\item \textbf{Plan d'action Activité Physique (Progressivité \textbf{obligatoire}) :}
\begin{itemize}
\item \textbf{Phase 1 (Sem 1--2) :} Marche rapide 20--30 min, 3--4 j/sem. Intensité : "Peut parler, pas chanter". Échauffement 5 min / Retour au calme 5 min.
\item \textbf{Phase 2 (Sem 3--8) :} Augmenter durée (30--45 min) et fréquence (5 j/sem). Introduire 1 séance renforcement (poids de corps : squats, pompes genoux, gainage, montées d'escaliers).
\item \textbf{Phase 3 (Entretien) :} 150+ min modérée ou 75 min soutenue + 2 séances renforcement. Varier plaisirs (Danse, Football, Natation, Jardinage, Vélo).
\item \textbf{Intégration vie quotidienne :} Marche marché/école/travail, Escaliers, Pause active au bureau.
\end{itemize}
\item \textbf{Suivi et Évaluation (Contractualisation) :}
\begin{itemize}
\item Carnet de bord (Aliments, AP, Ressenti, Poids 1 fois/sem même conditions).
\item Réévaluation mensuelle (Poids, Tour de taille, Tolérance à l'effort).
\item Bilan biologique à 3 mois (Glycémie, Lipides, Foie, Rein).
\item \textbf{Alerte} : Douleur thoracique, essoufflement anormal, malaise $\rightarrow$ Arrêt + Avis médical immédiat.
\end{itemize}
\end{enumerate}
\end{methode}

\begin{exoresolu}[Programme de prévention pour Madame Etoundi, 38 ans, Ebolowa]
\textit{Données :} $T = \SI{1,60}{m}$, $m = \SI{77}{kg}$, Sédentaire (commerçante assise), Alimentation : Fritures quotidiennes (beignets, poisson frit), Boissons sucrées (sodas, jus industriels), Peu de légumes/fruits/laitages.
\begin{enumerate}
\item Établissez le diagnostic anthropométrique et nutritionnel.
\item Proposez un programme sur 3 mois (Objectifs, Alimentation, AP, Suivi).
\end{enumerate}
\tcblower
\textbf{1. Diagnostic.}
\begin{itemize}
\item \textbf{IMC} = $77 / 1,60^2 = 77 / 2,56 = \mathbf{30,1\ kg/m^2}$ $\rightarrow$ \cle{Obésité classe I}.
\item \textbf{Tour de taille} (estimé) $> 88$ cm $\rightarrow$ \cle{Obésité androïde / viscérale} (Facteur de risque métabolique majeur).
\item \textbf{Alimentation} : Excès Lipides (fritures, AG saturés/trans), Excès Sucres rapides (boissons), Déficit Fibres (légumes/fruits), Déficit Ca (laitages), Excès Sel (cubes, poisson séché).
\item \textbf{AP} : NAP $\approx 1,2$ (Sédentaire strict). Dépense $\approx$ MB ($\sim 1500$ kcal) + faible AP. Apports estim. $> 2500$ kcal $\rightarrow$ Balance $> +1000$ kcal/j.
\item \textbf{Risques immédiats} : Pré-diabète / DT2, HTA, Dyslipidémie, Stéatose hépatique (NASH).
\end{itemize}

\textbf{2. Objectifs 3 mois (Réalistes).}
\begin{itemize}
\item Perte de masse : \textbf{4 à 6 kg} (cible $\sim 71-73$ kg, IMC $\sim 27,7-28,5$ : sortie obésité vers surpoids).
\item AP : Atteindre \textbf{150 min marche rapide/sem} (5 $\times$ 30 min).
\item Alimentation : 7 groupes couverts, Suppression sodas/fritures, Eau 2 L/j.
\end{itemize}

\textbf{3. Programme Alimentaire (Exemple de journée type).}
\begin{center}
\begin{tabularx}{\linewidth}{|l|X|}\hline
\rowcolor{popblueL} \textbf{Moment} & \textbf{Composition} \\ \hline
\textbf{Matin} & Bouillie maïs/mil \textbf{peu sucrée} (ou pain complet + beurre léger) + \textbf{1 Fruit (banane/papaye)} + \textbf{1 Yaourt nature / Lait caillé} + Eau/Thé non sucré \\ \hline
\textbf{Midi} & \textbf{Ndolé} (feuilles ++, arachides modéré, poisson fumé/viande maigre) + \textbf{Miondo/Foufou} (portion contrôlée $\approx$ 200g cuit) + \textbf{Salade crue} (tomate, concombre, carotte râpée, vinaigrette huile rouge/citron) + Eau \\ \hline
\textbf{Collation (16h)} & \textbf{1 Fruit de saison} + \textbf{Qq amandes/arachides non grillées salées} (30g) ou 1 Yaourt \\ \hline
\textbf{Soir} & \textbf{Poisson braisé / Poulet grillé (sans peau)} + \textbf{Légumes vapeur/sauté peu huilé} (Gombo, Folong, Épinards) + \textbf{Riz complet / Plantain bouilli} (portion modérée) + Eau \\ \hline
\end{tabularx}
\end{center}
\textit{Principes :} Réduction huile (cuisson braise/vapeur), Suppression beignets/sodas, Introduction Fruits/Laitages/Crudités, Fractionnement.

\textbf{4. Programme Activité Physique (Progression 12 semaines).}
\begin{center}
\begin{tabularx}{\linewidth}{|l|X|X|}\hline
\rowcolor{popblueL} \textbf{Période} & \textbf{Endurance (Marche rapide / Danse active)} & \textbf{Renforcement / Vie Quotidienne} \\ \hline
\textbf{Sem 1--2} & 20 min, 4 j/sem (Allure : discute aisément) & Étirements quotidiens 5 min. Monter escaliers 1 étage. \\ \hline
\textbf{Sem 3--6} & 30 min, 5 j/sem (Allure : discute difficilement) & + 10 min Renfo (Squats chaise, Pompes mur, Gainage 20s, Mollets) 2 j/sem. \\ \hline
\textbf{Sem 7--12} & 40--45 min, 5 j/sem (Inclure 2 $\times$ 10 min soutenue) & + 15 min Renfo complet 2--3 j/sem. Marche marché/aller-retour domicile. \\ \hline
\end{tabularx}
\end{center}
\textit{Note :} Certificat médical d'aptitude recommandé avant reprise (HTA, obésité). Surveillance : Fréquence cardiaque cible (FCmax théo = 220 - 38 = 182 bpm ; Cible 60--70\% FCmax = 110--127 bpm).

\textbf{5. Suivi.}
\begin{itemize}
\item Pesée hebdo (vendredi matin, à jeun, après miction). Courbe de poids.
\item Carnet alimentaire (photos repas / application) 3 j/sem.
\item Contrôle tensionnel mensuel (pharmacie/CSI).
\item RDV Centre de Santé Mois 3 : Poids, TA, Glycémie à jeun, Lipidogramme, Créaninémie. Ajustement programme.
\end{itemize}
\textbf{Conclusion} : Ce programme, progressif, encadré, associant restriction modérée des apports (qualité + quantité) et augmentation progressive des dépenses, vise un amaigrissement durable ($-5\%$ masse initiale) et la rémission des facteurs de risque métaboliques.
\end{exoresolu}

\courssub{Exploitation de données épidémiologiques (Savoir-faire transversal)}

\begin{methode}[Analyser un tableau ou graphique de données de santé publique]
\begin{enumerate}
\item \textbf{Identifier} : Titre, Source, Année, Population, Variables (Indicateurs : Prévalence, Incidence, Mortalité ; Facteurs : IMC, AP, Alim, Âge, Sexe, Milieu), Unités.
\item \textbf{Décrire les tendances} : Variations (augmentation/diminution), Comparaisons (Groupes, Territoires, Temps). \textbf{Chiffrer} (écarts absolus, relatifs, ratios).
\item \textbf{Interpréter} : Mobiliser les connaissances du cours (Mécanismes physiologiques, Transition nutritionnelle, Déterminants sociaux, Environnement built).
\item \textbf{Conclure} : Réponse directe à la question, lien cause-effet argumenté, perspective de prévention.
\end{enumerate}
\end{methode}

\begin{exoresolu}[Analyse comparative : Douala vs Milieu Rural Extrême-Nord]
\textit{Données : Enquête adultes 25--50 ans.}
\begin{center}
\begin{tabularx}{\linewidth}{|l|X|X|}\hline
\rowcolor{popblueL} \textbf{Localité} & \textbf{Obésité (\% IMC $\geq$ 30)} & \textbf{AP Régulière (\% $\geq$ 150 min/sem)} \\ \hline
Ville de Douala (Urbain) & 24 & 18 \\ \hline
Zone Rurale Extrême-Nord & 7 & 62 \\ \hline
\end{tabularx}
\end{center}
\begin{enumerate}
\item Comparez la prévalence de l'obésité. Calculez le risque relatif.
\item Expliquez cet écart par les déterminants du mode de vie.
\end{enumerate}
\tcblower
\textbf{1. Comparaison chiffrée.}
\begin{itemize}
\item Prévalence Obésité Douala = 24\% ; Rurale = 7\%.
\item \textbf{Risque Relatif (RR)} = $24 / 7 \approx \mathbf{3,4}$.
\item \textbf{Différence absolue} = $24 - 7 = \mathbf{17}$ points de pourcentage.
\item Inversement, AP régulière : Douala 18\% vs Rurale 62\% (RR = 0,29). Fort gradient inverse.
\end{itemize}

\textbf{2. Explication étiologique (Modèle éco-social).}
\begin{itemize}
\item \textbf{Environnement physique (Built environment) :} Douala = Ville dense, motorisée (taxis, motos), insécurité routière, manque d'espaces verts/sportifs $\rightarrow$ \cle{Sédentarité contrainte} (Déplacements passifs, Travail bureau/commerce assis). Rural = Déplacements actifs (marche champs/marché/eau), Travaux champêtres manuels $\rightarrow$ \cle{Dépenses énergétiques élevées} structurelles.
\item \textbf{Environnement alimentaire :} Douala = \cle{Transition nutritionnelle} : Disponibilité/Accessibilité/Prix bas aliments ultra-transformés (pâtes, riz blanc, huile raffinée, sodas, snacks, fritures rue), Marketing agressif. Rural = Alimentation \cle{traditionnelle} : Céréales complètes (mil, sorgho, maïs), Légumineuses (niébé), Légumes feuilles, Poisson séché/fumé, Peu de sucres ajoutés, Portions adaptées à l'activité.
\item \textbf{Facteurs socio-économiques/culturels :} Douala = Stress chronique, horaires décalés, sommeil court, statut social associé à la corpulence (perception positive du "poids de la réussite"). Rural = Rythmes circadiens respectés, norme corporelle mince.
\end{itemize}
\textbf{Synthèse} : L'écart de prévalence (x3,4) illustre l'impact majeur de la \cle{transition épidémiologique et nutritionnelle} urbaine. La sédentarité et l'hyper-disponibilité alimentaire déséquilibrée sont les moteurs de l'épidémie d'obésité/MNT. La prévention nécessite des politiques publiques (urbanisme actif, fiscalité soda, régulation marketing, éducation) et non seulement des conseils individuels.
\end{exoresolu}

\begin{attention}[Les 5 erreurs fatales au Probatoire (SVT Éducation à la Santé)]
\begin{enumerate}
\item \textbf{IMC : Unités et conversion.} Oublier de convertir cm en m ($T=175$ au lieu de $1,75$) $\rightarrow$ IMC divisé par 10 000. \textbf{Perte de tous les points de calcul.} Oublier l'unité $\mathrm{kg/m^2}$ ou la classe OMS $\rightarrow$ Réponse incomplète.
\item \textbf{Équivalents énergétiques.} Confondre \textbf{9 kcal/g (Lipides)} et 4 kcal/g (Glucides/Protides). Inverser les coefficients $\rightarrow$ Calcul totalement faux. Ne pas savoir calculer les \%\ de répartition.
\item \textbf{Mécanismes physiologiques : Vocabulaire imprécis.} Écrire « \emph{le sucre donne le diabète} », « \emph{la graisse bouche les artères} », « \emph{l'insuline détruit le sucre} ». \textbf{Sanctionné.} Vocabulaire exigé : \cle{Insuline}, \cle{Récepteur}, \cle{GLUT4}, \cle{Insulinorésistance}, \cle{Hyperglycémie}, \cle{Triglycérides}, \cle{Adipocyte}, \cle{LDL-c}, \cle{HDL-c}, \cle{Athérosclérose}, \cle{Plaque d'athérome}, \cle{Thrombose}.
\item \textbf{Programme de prévention irréaliste / Dangereux.} Perte $> 1$ kg/sem, Jeûne, Régime mono-aliment, Sport intense d'emblée (risque cardiaque, tendinites), Suppression totale lipides/glucides. \textbf{Note zéro sur la partie proposition.} Exigé : Progressivité, Équilibre, Plaisir, Suivi médical.
\item \textbf{Analyse de données : Affirmation sans chiffre.} « Il y a plus d'obésité en ville » sans citer les 24\% et 7\%, sans calculer le RR ou l'écart. \textbf{Points perdus.} Toute conclusion doit être \textbf{chiffrée et sourcée} du document.
\end{enumerate}
\end{attention}

\newpage

\coursec{Exercices}

\courssub{Je vérifie mes connaissances}

\begin{exercice}[QCM : les fondamentaux de l'équilibre alimentaire]
Pour chaque question, une seule réponse est exacte. Recopie la lettre correspondante.

\textbf{1.} L'énergie fournie par 1 g de lipides est environ :
\begin{itemize}
\item[a)] 4 kcal ; \quad b) 7 kcal ; \quad c) 9 kcal ; \quad d) 17 kcal.
\end{itemize}

\textbf{2.} Selon l'OMS, un adulte devrait consommer chaque jour au moins :
\begin{itemize}
\item[a)] 100 g de fruits et légumes ; \quad b) 400 g de fruits et légumes ; \quad c) 1 kg de fruits et légumes ; \quad d) 50 g de fruits et légumes.
\end{itemize}

\textbf{3.} Un IMC de 27 kg/m$^2$ correspond à :
\begin{itemize}
\item[a)] une dénutrition ; \quad b) une corpulence normale ; \quad c) un surpoids ; \quad d) une obésité.
\end{itemize}

\textbf{4.} L'OMS recommande aux adultes de limiter le sel à moins de :
\begin{itemize}
\item[a)] 5 g par jour ; \quad b) 15 g par jour ; \quad c) 25 g par jour ; \quad d) 50 g par jour.
\end{itemize}

\textbf{5.} Pour les adolescents (5--17 ans), la durée minimale d'activité physique quotidienne recommandée est :
\begin{itemize}
\item[a)] 15 minutes ; \quad b) 30 minutes ; \quad c) 60 minutes ; \quad d) 3 heures.
\end{itemize}
\end{exercice}

\begin{exercice}[Vrai ou faux, à justifier]
Réponds par vrai ou faux, puis justifie chaque réponse en une ou deux phrases.

\textbf{1.} Les glucides, les protéines et les lipides fournissent tous la même énergie par gramme.

\textbf{2.} Le diabète de type 2 est une maladie infectieuse transmise par un virus.

\textbf{3.} Un adolescent a des besoins énergétiques plus élevés qu'un adulte sédentaire, en raison de sa croissance.

\textbf{4.} L'obésité abdominale est un facteur de risque des maladies cardiovasculaires.

\textbf{5.} Seule une activité sportive intense en salle permet de lutter contre la sédentarité.
\end{exercice}

\begin{exercice}[Définitions et notions clés]
Définis en deux ou trois phrases chacun des termes suivants :
\begin{enumerate}
\item l'équilibre alimentaire ;
\item l'indice de masse corporelle (IMC) ;
\item la sédentarité ;
\item le diabète de type 2 ;
\item la ration alimentaire.
\end{enumerate}
\end{exercice}

\begin{exercice}[Calculs directs]
\textbf{1.} Un aliment contient 12 g de glucides, 5 g de lipides et 8 g de protéines. Sachant que 1 g de glucides ou de protéines fournit 4 kcal et 1 g de lipides 9 kcal, calcule l'énergie totale fournie par cet aliment.

\textbf{2.} Calcule l'IMC d'une personne de masse 60 kg et de taille 1,65 m (arrondis à 0,1 près). Rappel : IMC = masse (kg) / taille$^2$ (m$^2$).

\textbf{3.} Un adulte dépense 2 200 kcal par jour. Il consomme 2 500 kcal par jour. Calcule son excès énergétique quotidien.
\end{exercice}

\courssub{Je m'entraîne}

\begin{exercice}[Classer des corpulences]
Voici les IMC de cinq élèves de Première D : A : 16,8 ; B : 21,5 ; C : 26,4 ; D : 31,2 ; E : 38,5 (en kg/m$^2$).

\textbf{1.} À l'aide des seuils de l'OMS (dénutrition : IMC $<$ 18,5 ; corpulence normale : 18,5 à 24,9 ; surpoids : 25 à 29,9 ; obésité modérée : 30 à 34,9 ; obésité sévère : 35 à 39,9 ; obésité morbide : $\geq$ 40), classe chaque élève.

\textbf{2.} Quels élèves présentent un risque sanitaire accru ? Justifie.

\textbf{3.} Propose pour l'élève D deux conseils alimentaires et deux conseils d'activité physique.
\end{exercice}

\begin{exercice}[Évaluer une ration alimentaire]
Voici le menu journalier de Rodrigue, élève de 16 ans, et la table de composition simplifiée des aliments consommés.

\begin{center}
\begin{tabularx}{\linewidth}{|l|X|c|}
\hline
\rowcolor{popblueL} \textbf{Repas} & \textbf{Aliments} & \textbf{Énergie (kcal)} \\
\hline
Petit-déjeuner & Beignets (3), bouillie sucrée & 650 \\
\hline
Déjeuner & Riz blanc (300 g cuit), sauce arachide, boisson gazeuse (33 cL) & 1 150 \\
\hline
Goûter & Biscuits (paquet) & 450 \\
\hline
Dîner & Fufu, eru, peu d'eau & 900 \\
\hline
\end{tabularx}
\end{center}

\textbf{1.} Calcule l'apport énergétique total de la journée.

\textbf{2.} Les besoins d'un adolescent de 16 ans sont d'environ 2 800 kcal par jour. Compare l'apport de Rodrigue à ses besoins.

\textbf{3.} Ce menu contient très peu de fruits, de légumes et de produits laitiers, et beaucoup de sucres rapides. Identifie deux déséquilibres nutritionnels et propose pour chacun une correction concrète adaptée au contexte local (marché camerounais).
\end{exercice}

\begin{exercice}[Bilan énergétique d'un adulte]
Monsieur N., employé de bureau à Yaoundé, consomme chaque jour 2 900 kcal alors qu'il n'en dépense que 2 500 kcal. On rappelle que 1 kg de masse grasse correspond à environ 7 700 kcal.

\textbf{1.} Calcule son excès énergétique quotidien.

\textbf{2.} Estime sa prise de masse grasse au bout de 3 mois (90 jours), en supposant que tout l'excès est stocké.

\textbf{3.} Sachant qu'une marche rapide de 30 minutes dépense environ 150 kcal, propose une combinaison « réduction alimentaire + activité physique » permettant d'annuler exactement l'excès quotidien.
\end{exercice}

\begin{exercice}[Programme d'activité physique]
Sandrine, 17 ans, est très sédentaire : elle passe plus de 8 heures par jour assise (cours, téléphone, télévision) et ne pratique aucun sport.

\textbf{1.} Rappelle les recommandations de l'OMS en matière d'activité physique pour les adolescents de 5 à 17 ans.

\textbf{2.} Élabore pour Sandrine un programme hebdomadaire réaliste (jours, durées, types d'activités adaptées à une élève : marche, danse, football, corde à sauter, tâches actives...) respectant ces recommandations.

\textbf{3.} Cite deux bénéfices attendus de ce programme sur sa santé.
\end{exercice}

\courssub{Je me prépare au Probatoire}

\begin{exercice}[Sujet type Probatoire : IMC et risques sanitaires]
Lors d'une visite médicale scolaire au lycée de Bafoussam, on mesure les paramètres suivants chez trois élèves de Première D :

\begin{center}
\begin{tabularx}{\linewidth}{|l|X|X|X|}
\hline
\rowcolor{popblueL} \textbf{Élève} & \textbf{Masse (kg)} & \textbf{Taille (m)} & \textbf{Tour de taille (cm)} \\
\hline
Élève 1 & 85 & 1,70 & 102 \\
\hline
Élève 2 & 52 & 1,70 & 74 \\
\hline
Élève 3 & 96 & 1,65 & 110 \\
\hline
\end{tabularx}
\end{center}

Seuils OMS : dénutrition : IMC $<$ 18,5 ; normal : 18,5 à 24,9 ; surpoids : 25 à 29,9 ; obésité : IMC $\geq$ 30.

\textbf{1.} Calcule l'IMC de chaque élève (arrondis à 0,1 près) et classe sa corpulence.

\textbf{2.} Quel élève présente le risque sanitaire le plus élevé ? Justifie en utilisant deux indicateurs du tableau.

\textbf{3.} Cite deux maladies liées au mode de vie auxquelles cet élève est exposé, en précisant pour chacune le mécanisme simplifié qui la relie à l'excès de masse grasse.

\textbf{4.} Propose à cet élève deux mesures correctives alimentaires et deux mesures d'activité physique, en les justifiant brièvement.
\end{exercice}

\begin{exercice}[Sujet type Probatoire : de la sédentarité au diabète de type 2]
Dans un quartier de Douala, le service de santé communautaire constate une augmentation des cas de diabète de type 2 chez des adultes de 35 à 50 ans. L'enquête révèle que la majorité des personnes atteintes travaille assise toute la journée, consomme quotidiennement des boissons sucrées et ne pratique aucune activité physique.

\textbf{1.} Définis le diabète de type 2 et précise la valeur de la glycémie à jeun qui permet de le diagnostiquer.

\textbf{2.} Explique, en 8 à 10 lignes, la chaîne de mécanismes par laquelle l'excès chronique de sucres associé à la sédentarité conduit au diabète de type 2 (on attend les notions de glycémie, d'insuline, de pancréas et d'insulinorésistance).

\textbf{3.} Schématise cette chaîne sous forme d'un diagramme de cause à effet annoté (au moins quatre étapes).

\textbf{4.} Le service de santé souhaite lancer une campagne de prévention. Rédige trois recommandations chiffrées (alimentation et activité physique) que tu afficherais au centre de santé, en t'appuyant sur les recommandations de l'OMS.

\textbf{5.} Explique pourquoi la pratique régulière d'une activité physique améliore l'action de l'insuline et contribue ainsi à prévenir cette maladie.
\end{exercice}

\begin{exercice}[Mini-projet noté : affiche de prévention des maladies cardiovasculaires]
Dans le cadre de la semaine de la santé de ton lycée, le club de Biologie te charge de concevoir une affiche de prévention des maladies cardiovasculaires destinée à toute la communauté scolaire (élèves, enseignants, parents).

\textbf{1.} Rappelle ce qu'est une maladie cardiovasculaire et cite-en deux exemples.

\textbf{2.} Identifie quatre facteurs de risque liés au mode de vie (alimentation et inactivité physique) que ton affiche doit cibler.

\textbf{3.} Rédige le contenu de ton affiche : un titre accrocheur, un slogan, au moins trois recommandations chiffrées (par exemple : quantité de fruits et légumes, limite de sel, durée d'activité physique, limite de boissons sucrées) et un message adapté au contexte camerounais (aliments locaux à privilégier ou à limiter).

\textbf{4.} Justifie scientifiquement chacune de tes trois recommandations chiffrées en une ou deux phrases (effet sur le cœur, les vaisseaux, la tension artérielle ou le poids).

\textbf{5.} Propose une activité concrète que le club de Biologie pourrait organiser pendant la semaine de la santé pour compléter l'affiche (dépistage, séance de sport, dégustation de fruits locaux...) et explique son intérêt préventif.
\end{exercice}



\newpage

\coursec{Corrigés des exercices}

\courssub{Je vérifie mes connaissances}

\begin{corrige}[Exercice 1 --- QCM : les fondamentaux de l'équilibre alimentaire]
\textbf{1. c)} 1 g de lipides fournit environ 9 kcal (contre environ 4 kcal pour 1 g de glucides ou de protéines). C'est pourquoi les aliments gras sont les plus « denses » en énergie : à masse égale, ils apportent plus du double de calories.

\textbf{2. b)} L'OMS recommande au moins 400 g de fruits et légumes par jour, soit environ 5 portions, pour prévenir les maladies non transmissibles (diabète de type 2, maladies cardiovasculaires, certains cancers).

\textbf{3. c)} Un IMC compris entre 25 et 29,9 kg/m$^2$ définit le surpoids ; l'obésité ne commence qu'à partir de 30 kg/m$^2$. La corpulence normale correspond à un IMC compris entre 18,5 et 24,9 kg/m$^2$.

\textbf{4. a)} L'OMS recommande de limiter le sel à moins de 5 g par jour (environ une cuillère à café rase) pour prévenir l'hypertension artérielle, premier facteur de risque d'accident vasculaire cérébral.

\textbf{5. c)} Les 5--17 ans doivent pratiquer au moins 60 minutes par jour d'activité physique d'intensité modérée à soutenue.

\emph{Point de vigilance :} ne pas confondre les seuils « surpoids » (25 kg/m$^2$) et « obésité » (30 kg/m$^2$), ni les recommandations « adolescents » (au moins 60 min/j) et « adultes » (au moins 150 min/semaine).
\end{corrige}

\begin{corrige}[Exercice 2 --- Vrai ou faux, à justifier]
\textbf{1. Faux.} Glucides et protéines fournissent environ 4 kcal/g, tandis que les lipides fournissent environ 9 kcal/g, soit plus du double.

\textbf{2. Faux.} Le diabète de type 2 est une maladie \emph{non transmissible}, liée au mode de vie (excès alimentaire, sédentarité) et à des facteurs génétiques ; il n'est causé par aucun microbe et ne se transmet pas d'une personne à l'autre.

\textbf{3. Vrai.} La croissance (allongement des os, développement des muscles et des organes) exige un supplément d'énergie et de nutriments : un adolescent actif peut avoir besoin de 2 500 à 3 000 kcal/j, davantage qu'un adulte sédentaire.

\textbf{4. Vrai.} La graisse abdominale (viscérale) sécrète des substances qui favorisent l'inflammation, l'athérosclérose et l'hypertension artérielle ; un tour de taille élevé est donc un indicateur de risque cardiovasculaire, indépendamment de l'IMC.

\textbf{5. Faux.} Toute activité d'intensité modérée compte : marche rapide, vélo, danse, travaux champêtres, corde à sauter, jeux de ballon. L'essentiel est de bouger régulièrement et de réduire le temps passé assis.
\end{corrige}

\begin{corrige}[Exercice 3 --- Définitions et notions clés]
\textbf{1. Équilibre alimentaire :} état dans lequel les apports alimentaires (énergie, protéines, lipides, glucides, vitamines, sels minéraux, eau) couvrent exactement les besoins de l'organisme, ni en déficit ni en excès, grâce à une alimentation variée et à des quantités adaptées.

\textbf{2. Indice de masse corporelle (IMC) :} rapport de la masse (en kg) sur le carré de la taille (en m$^2$) : IMC $= m/t^2$. Il permet de classer la corpulence (dénutrition, corpulence normale, surpoids, obésité) et d'estimer le risque sanitaire.

\textbf{3. Sédentarité :} mode de vie caractérisé par une très faible dépense physique et de longues périodes passées assis ou allongé (absence d'activité physique régulière).

\textbf{4. Diabète de type 2 :} maladie chronique caractérisée par une hyperglycémie permanente (glycémie à jeun $\geq$ 1,26 g/L, soit 7 mmol/L, à deux reprises) due à une résistance des cellules à l'insuline (insulinorésistance), souvent associée à un épuisement progressif du pancréas endocrine.

\textbf{5. Ration alimentaire :} ensemble des aliments consommés par un individu en une journée. Elle doit être évaluée en quantité (énergie apportée) et en qualité (diversité des nutriments) par rapport aux besoins recommandés.

\emph{Point de vigilance :} dans la définition de l'IMC, citer obligatoirement la formule et les unités (kg et m$^2$) ; dans celle du diabète, donner le seuil de glycémie et le mécanisme (insulinorésistance).
\end{corrige}

\begin{corrige}[Exercice 4 --- Calculs directs]
\textbf{1.} Énergie $= (12 \times 4) + (5 \times 9) + (8 \times 4) = 48 + 45 + 32 = 125$ kcal.

\textbf{2.} IMC $= \dfrac{60}{1,65^2} = \dfrac{60}{2,7225} \approx 22,0$ kg/m$^2$ : corpulence normale (valeur comprise entre 18,5 et 24,9 kg/m$^2$).

\textbf{3.} Excès $= 2\,500 - 2\,200 = 300$ kcal par jour. Cet excès, s'il se répète chaque jour, sera stocké sous forme de graisse dans le tissu adipeux.

\emph{Point de vigilance :} penser à élever la taille au carré \emph{avant} de diviser ; donner le résultat avec son unité (kg/m$^2$) et conclure par la classe de corpulence correspondante.
\end{corrige}

\courssub{Je m'entraîne}

\begin{corrige}[Exercice 5 --- Classer des corpulences]
\textbf{1.} A : dénutrition (16,8 $<$ 18,5) ; B : corpulence normale (18,5 $\leq$ 21,5 $\leq$ 24,9) ; C : surpoids (25 $\leq$ 26,4 $\leq$ 29,9) ; D : obésité modérée (30 $\leq$ 31,2 $\leq$ 34,9) ; E : obésité sévère (35 $\leq$ 38,5 $\leq$ 39,9).

\textbf{2.} Tous sauf B présentent un risque accru : A risque carences, anémie et baisse des défenses immunitaires ; C, D et E risquent diabète de type 2, hypertension artérielle et maladies cardiovasculaires, le risque croissant avec l'IMC (E est le plus exposé).

\textbf{3.} Pour D : \emph{alimentation} --- réduire les boissons sucrées et les fritures, augmenter les légumes (ndolé, okok, feuilles de manioc) et les fruits locaux ; \emph{activité physique} --- au moins 60 min/j de marche rapide ou de sport, et limiter le temps assis devant les écrans.

\emph{Point de vigilance :} la dénutrition est aussi un risque sanitaire ; ne pas oublier de la mentionner dans la question 2.
\end{corrige}

\begin{corrige}[Exercice 6 --- Évaluer une ration alimentaire]
\textbf{1.} Total $= 650 + 1\,150 + 450 + 900 = 3\,150$ kcal.

\textbf{2.} L'apport (3 150 kcal) dépasse les besoins (2 800 kcal) de 350 kcal/j : cet excès, à long terme, favorise la prise de masse grasse.

\textbf{3.} Déséquilibres possibles : \emph{(a)} excès de sucres rapides (boisson gazeuse, biscuits, bouillie sucrée) $\to$ remplacer la boisson gazeuse par de l'eau ou un jus de fruit local sans sucre ajouté (foléré non sucré, jus de goyave) ; \emph{(b)} déficit en fruits, légumes et produits laitiers $\to$ ajouter un fruit de saison (mangue, papaye, banane, avocat) chaque jour, une portion de légumes feuilles (ndolé, okok) et un lait caillé ou du kossam. On note aussi que l'excès calorique coexiste avec une mauvaise \emph{qualité} nutritionnelle : c'est la « double peine » de la malnutrition.

\emph{Point de vigilance :} une ration peut être à la fois excessive en énergie et carencée en micronutriments (vitamines, sels minéraux) ; l'évaluation porte sur la quantité \emph{et} sur la qualité.
\end{corrige}

\begin{corrige}[Exercice 7 --- Bilan énergétique d'un adulte]
\textbf{1.} Excès $= 2\,900 - 2\,500 = 400$ kcal/j.

\textbf{2.} Excès sur 90 jours : $400 \times 90 = 36\,000$ kcal. Prise de masse grasse : $\dfrac{36\,000}{7\,700} \approx 4,7$ kg en 3 mois (en admettant qu'environ 7 700 kcal stockées correspondent à 1 kg de graisse corporelle).

\textbf{3.} Il faut supprimer 400 kcal/j. Exemple : réduire l'alimentation de 250 kcal (supprimer la boisson gazeuse quotidienne et une portion de friture) et ajouter 30 min de marche rapide (150 kcal) : $250 + 150 = 400$ kcal. Toute combinaison dont la somme vaut 400 kcal est correcte (par exemple 100 kcal en moins + 60 min de marche, soit $100 + 2 \times 150 = 400$ kcal).

\emph{Point de vigilance :} vérifier que la combinaison proposée atteint \emph{exactement} 400 kcal et présenter le calcul de vérification.
\end{corrige}

\begin{corrige}[Exercice 8 --- Programme d'activité physique]
\textbf{1.} Pour les 5--17 ans, l'OMS recommande au moins 60 minutes par jour d'activité physique d'intensité modérée à soutenue, en limitant le temps sédentaire (écrans) ; des activités de renforcement musculaire et osseux (sauts, jeux de ballon) au moins 3 fois par semaine.

\textbf{2.} Exemple de programme : lundi : 30 min de marche rapide (trajet école) + 30 min de danse ; mardi : 60 min de football ou handball ; mercredi : 30 min de corde à sauter + 30 min de tâches actives (jardinage, courses) ; jeudi : 60 min de marche rapide ; vendredi : 60 min de danse ou basket ; samedi : 90 min d'activités libres (natation, vélo, match) ; dimanche : repos actif (promenade en famille). Chaque jour atteint ou dépasse 60 min, et les activités de renforcement (corde à sauter, ballon) apparaissent au moins 3 fois par semaine.

\textbf{3.} Bénéfices : amélioration de la sensibilité à l'insuline (prévention du diabète de type 2), renforcement du cœur et des vaisseaux, contrôle du poids, meilleure santé osseuse, meilleur sommeil et meilleure concentration scolaire.

\emph{Point de vigilance :} le programme doit être chiffré (durées en minutes) et vérifier explicitement le seuil de 60 min/j.
\end{corrige}

\courssub{Je me prépare au Probatoire}

\begin{corrige}[Exercice 9 --- Sujet type Probatoire : IMC et risques sanitaires]
\textbf{1.} Élève 1 : IMC $= \dfrac{85}{1,70^2} = \dfrac{85}{2,89} \approx 29,4$ kg/m$^2$ : surpoids. Élève 2 : IMC $= \dfrac{52}{2,89} \approx 18,0$ kg/m$^2$ : dénutrition (légère). Élève 3 : IMC $= \dfrac{96}{1,65^2} = \dfrac{96}{2,7225} \approx 35,3$ kg/m$^2$ : obésité (sévère).

\textbf{2.} L'élève 3 présente le risque le plus élevé : son IMC (35,3 kg/m$^2$) est le plus haut du groupe et atteint l'obésité sévère ; son tour de taille (110 cm) traduit une obésité abdominale, facteur de risque cardiovasculaire et métabolique majeur.

\textbf{3.} \emph{Diabète de type 2 :} l'excès de graisse, surtout abdominale, rend les cellules résistantes à l'insuline (insulinorésistance) ; le pancréas s'épuise et la glycémie reste durablement élevée. \emph{Maladies cardiovasculaires (hypertension, athérosclérose) :} l'excès de lipides circulants se dépose sur la paroi des artères (plaques d'athérome), ce qui rétrécit les vaisseaux, augmente la tension artérielle et fatigue le cœur.

\textbf{4.} \emph{Alimentation :} supprimer les boissons sucrées (réduction des sucres rapides qui aggravent l'insulinorésistance) ; augmenter légumes et fruits locaux et réduire les fritures (diminution de l'apport calorique et des graisses). \emph{Activité physique :} 60 min/j de marche rapide ou de sport (augmente la dépense énergétique et la sensibilité à l'insuline) ; réduire le temps assis devant les écrans (lutte contre la sédentarité).

\emph{Barème indicatif (sur 10) :} Q1 : 3 pts (1 pt par IMC calculé et classé) ; Q2 : 2 pts (choix + double justification IMC/tour de taille) ; Q3 : 2 pts (deux maladies avec mécanisme) ; Q4 : 3 pts (quatre mesures justifiées).

\emph{Points de vigilance :} écrire la formule de l'IMC avant le calcul ; arrondir à 0,1 ; justifier la question 2 avec \emph{deux} indicateurs chiffrés du tableau ; chaque mesure corrective doit être accompagnée d'une justification.
\end{corrige}

\begin{corrige}[Exercice 10 --- Sujet type Probatoire : de la sédentarité au diabète de type 2]
\textbf{1.} Le diabète de type 2 est une maladie chronique non transmissible caractérisée par une hyperglycémie permanente due à une résistance des cellules à l'insuline. Le diagnostic est posé lorsque la glycémie à jeun est supérieure ou égale à 1,26 g/L (7 mmol/L) à deux reprises.

\textbf{2.} La consommation chronique de sucres rapides provoque des élévations répétées et importantes de la glycémie. En réponse, le pancréas (îlots de Langerhans) sécrète de grandes quantités d'insuline, l'hormone qui fait entrer le glucose dans les cellules. La sédentarité aggrave la situation : les muscles, peu sollicités, consomment peu de glucose et l'excès énergétique est stocké en graisse, notamment abdominale. Cette graisse viscérale rend les cellules de moins en moins sensibles à l'insuline : c'est l'insulinorésistance. Le glucose pénètre alors mal dans les cellules et s'accumule dans le sang. Le pancréas, sollicité en permanence, finit par s'épuiser et produit moins d'insuline : la glycémie reste durablement élevée, c'est le diabète de type 2 installé.

\textbf{3.} Diagramme de cause à effet :

\begin{popfigure}
\begin{tikzpicture}[node distance=0.55cm, every node/.style={font=\small},
box/.style={draw=popblue, fill=popblueL, rounded corners, align=center, text width=4.6cm, inner sep=4pt},
fleche/.style={-{Stealth[length=2.5mm]}, thick, popdark}]
\node[box] (a) {Excès chronique de sucres rapides + sédentarité};
\node[box, below=of a] (b) {Hyperglycémies répétées et stockage de graisse abdominale};
\node[box, below=of b] (c) {Sécrétion excessive d'insuline par le pancréas};
\node[box, below=of c] (d) {Insulinorésistance des cellules (muscles, foie, tissu adipeux)};
\node[box, below=of d] (e) {Épuisement du pancréas $\Rightarrow$ hyperglycémie permanente : diabète de type 2};
\draw[fleche] (a) -- (b);
\draw[fleche] (b) -- (c);
\draw[fleche] (c) -- (d);
\draw[fleche] (d) -- (e);
\end{tikzpicture}
\legende{Chaîne de mécanismes conduisant de l'excès de sucres et de la sédentarité au diabète de type 2.}
\end{popfigure}

\textbf{4.} Trois recommandations chiffrées : \emph{(a)} consommer au moins 400 g de fruits et légumes par jour (5 portions) ; \emph{(b)} limiter le sel à moins de 5 g/j et les sucres libres à moins de 10 \% de l'apport énergétique (idéalement moins de 5 \%), en supprimant les boissons sucrées quotidiennes ; \emph{(c)} pratiquer au moins 150 minutes d'activité physique d'intensité modérée par semaine (par exemple 30 min de marche rapide 5 jours sur 7).

\textbf{5.} Pendant l'effort, les muscles consomment du glucose directement, ce qui fait baisser la glycémie. L'exercice régulier augmente aussi la sensibilité des cellules musculaires à l'insuline : le glucose y pénètre plus facilement, le pancréas est moins sollicité et la glycémie est mieux régulée. L'activité physique réduit enfin la graisse abdominale, cause majeure de l'insulinorésistance.

\emph{Barème indicatif (sur 10) :} Q1 : 1,5 pt (définition + seuil) ; Q2 : 3 pts (chaîne logique avec les quatre notions exigées) ; Q3 : 2 pts (schéma à au moins quatre étapes, fléché et annoté) ; Q4 : 2 pts (trois recommandations chiffrées conformes à l'OMS) ; Q5 : 1,5 pt.

\emph{Points de vigilance :} les quatre notions (glycémie, insuline, pancréas, insulinorésistance) doivent toutes apparaître dans la question 2 ; le schéma doit être fléché dans le sens cause $\to$ effet ; les recommandations doivent être \emph{chiffrées}, une recommandation vague (« manger équilibré ») ne rapporte pas de point.
\end{corrige}

\begin{corrige}[Exercice 11 --- Mini-projet noté : affiche de prévention des maladies cardiovasculaires]
\textbf{1.} Une maladie cardiovasculaire est une maladie qui atteint le cœur et/ou les vaisseaux sanguins. Exemples : l'hypertension artérielle, l'infarctus du myocarde, l'accident vasculaire cérébral (AVC), l'athérosclérose.

\textbf{2.} Quatre facteurs de risque liés au mode de vie : excès de sel ; excès de graisses (fritures) et de sucres rapides ; consommation insuffisante de fruits et légumes ; sédentarité et inactivité physique (on peut ajouter le tabac et l'alcool).

\textbf{3.} Exemple de contenu d'affiche. \emph{Titre :} « Protège ton cœur, il n'en existe qu'un ! ». \emph{Slogan :} « Moins de sel, plus de pas : un cœur qui dure. » \emph{Recommandations chiffrées :} mange au moins 400 g de fruits et légumes par jour ; limite le sel à moins de 5 g par jour ; bouge au moins 30 minutes par jour (adultes) ou 60 minutes (jeunes) ; remplace les boissons sucrées par de l'eau. \emph{Message local :} privilégie les légumes feuilles (ndolé, okok, folon), les fruits de saison (mangue, papaye, avocat), le poisson braisé plutôt que les fritures ; va mollo avec le cube d'assaisonnement et le sel de cuisine.

\textbf{4.} Justifications : \emph{400 g de fruits et légumes} apportent potassium, fibres et antioxydants qui aident à réguler la tension artérielle et protègent la paroi des artères. \emph{Moins de 5 g de sel/j} limite la rétention d'eau et l'augmentation du volume sanguin, donc prévient l'hypertension, premier facteur de risque d'AVC. \emph{30 à 60 min d'activité physique/j} renforce le muscle cardiaque, améliore la circulation et aide à contrôler le poids, réduisant la surcharge de travail du cœur.

\textbf{5.} Exemple d'activité : organiser un stand de dépistage gratuit (mesure de la tension artérielle, du tour de taille et calcul de l'IMC par les élèves encadrés par un infirmier) suivi d'une séance collective de marche rapide ou de danse sportive et d'une dégustation de fruits locaux. Intérêt préventif : dépister tôt l'hypertension et le surpoids permet d'agir avant l'apparition des complications, et la démonstration pratique montre que la prévention est accessible à tous, sans matériel coûteux.

\emph{Barème indicatif (sur 10) :} Q1 : 1 pt ; Q2 : 2 pts (quatre facteurs correctement identifiés) ; Q3 : 3 pts (titre, slogan, trois recommandations chiffrées, message local) ; Q4 : 2 pts (justifications scientifiques) ; Q5 : 2 pts (activité concrète + intérêt préventif expliqué).

\emph{Points de vigilance :} distinguer les facteurs de risque « liés au mode de vie » (modifiables) des facteurs non modifiables (âge, hérédité) ; chaque recommandation chiffrée doit être accompagnée de sa justification scientifique ; l'activité proposée doit être réaliste dans le cadre d'un lycée.
\end{corrige}

\newpage

\coursec{Fiche bilan}

\begin{popfigure}
\centering
\begin{center}\tcbox[colback=popink,colframe=popink,coltext=white,arc=9pt,boxsep=4pt,left=6pt,right=6pt]{\bfseries\small Lutte : Mauvaise Alimentation \& Inactivité}\end{center}
\vspace{-2pt}
\begin{tcbraster}[raster columns=3,raster equal height=rows,raster column skip=6pt,raster row skip=6pt,enhanced,blankest]
\begin{tcolorbox}[enhanced,colback=popgreen,colframe=popgreen,coltext=white,boxrule=0.9pt,arc=3pt,left=4pt,right=4pt,top=3pt,bottom=3pt,boxsep=1pt,halign=left]\bfseries\scriptsize 1. Besoins \& Équilibre alimentaire\end{tcolorbox}
\begin{tcolorbox}[enhanced,colback=poporange,colframe=poporange,coltext=white,boxrule=0.9pt,arc=3pt,left=4pt,right=4pt,top=3pt,bottom=3pt,boxsep=1pt,halign=left]\bfseries\scriptsize 2. IMC : Calcul \& Interprétation\end{tcolorbox}
\begin{tcolorbox}[enhanced,colback=poppurple,colframe=poppurple,coltext=white,boxrule=0.9pt,arc=3pt,left=4pt,right=4pt,top=3pt,bottom=3pt,boxsep=1pt,halign=left]\bfseries\scriptsize 3. Conséquences Pathologiques\end{tcolorbox}
\begin{tcolorbox}[enhanced,colback=poppink,colframe=poppink,coltext=white,boxrule=0.9pt,arc=3pt,left=4pt,right=4pt,top=3pt,bottom=3pt,boxsep=1pt,halign=left]\bfseries\scriptsize 4. Recos Alimentation\end{tcolorbox}
\begin{tcolorbox}[enhanced,colback=popgold,colframe=popgold,coltext=white,boxrule=0.9pt,arc=3pt,left=4pt,right=4pt,top=3pt,bottom=3pt,boxsep=1pt,halign=left]\bfseries\scriptsize 5. Recos Activité Physique\end{tcolorbox}
\begin{tcolorbox}[enhanced,colback=popdark,colframe=popdark,coltext=white,boxrule=0.9pt,arc=3pt,left=4pt,right=4pt,top=3pt,bottom=3pt,boxsep=1pt,halign=left]\bfseries\scriptsize 6. Programme Prévention\end{tcolorbox}
\begin{tcolorbox}[enhanced,colback=popmuted,colframe=popmuted,coltext=white,boxrule=0.9pt,arc=3pt,left=4pt,right=4pt,top=3pt,bottom=3pt,boxsep=1pt,halign=left]\bfseries\scriptsize 7. Évaluation Ration Alimentaire\end{tcolorbox}
\end{tcbraster}

\legende{Carte mentale de la leçon ND10 : organisation des savoirs autour de la prévention des maladies liées au mode de vie.}
\end{popfigure}

\begin{bilan}
\courssub{Définitions clés}
\begin{itemize}
    \item \cle{Équilibre alimentaire} : Ration couvrant quantitativement et qualitativement les \cle{Besoins Nutritionnels Recommandés (BNR/ANC)} de l'individu (énergie, macronutriments, micronutriments, eau, fibres).
    \item \cle{Métabolisme de base (MB)} : Énergie minimale pour les fonctions vitales au repos complet (thermoneutralité, jeûne). Dépend du sexe, âge, masse, taille.
    \item \cle{Niveau d'Activité Physique (NAP)} : Facteur multiplicatif du MB (1,3 à 2,4) selon l'intensité de l'activité (sédentaire $\rightarrow$ très actif).
    \item \cle{Indice de Masse Corporelle (IMC)} : $IMC = \frac{m}{t^2}$ ($m$ en kg, $t$ en m). Indicateur de corpulence, non de composition corporelle.
    \item \cle{Sédentarité} : Temps d'éveil passé en position assise/allongée avec dépense énergétique $\le$ 1,5 MET (ex: écrans, bureau, transport). Distinct de l'inactivité physique (non-respect des seuils d'AP).
    \item \cle{Obésité} : Excès de masse grasse (IMC $\ge$ 30 kg/m$^2$), facteur de risque majeur de maladies chroniques non transmissibles (MCNT).
    \item \cle{Diabète de type 2 (DT2)} : Hyperglycémie chronique par résistance à l'insuline + défaut de sécrétion insulinique. Lié à surcharge pondérale et inactivité.
    \item \cle{Maladies Cardio-Vasculaires (MCV)} : Atteintes du cœur et vaisseaux (AVC, IDM, HTA). Favorsées par dyslipidémie, HTA, diabète, tabagisme, sédentarité.
\end{itemize}

\courssub{Formules \& Repères numériques à connaître}
\begin{center}
\begin{tabularx}{\linewidth}{|l|X|X|}\hline
\rowcolor{popblueL} \textbf{Notion} & \textbf{Formule / Valeur de référence} & \textbf{Commentaire / Unités} \\ \hline
IMC & $IMC = \dfrac{m}{t^2}$ & $m$ (kg), $t$ (m) \\ \hline
Interprétation IMC (Adulte OMS) & Maigreur $<$ 18,5 \textbar\ Normal 18,5--24,9 \textbar\ Surpoids 25--29,9 \textbar\ Obésité I 30--34,9 \textbar\ Obésité II 35--39,9 \textbar\ Obésité III $\ge$ 40 & Risque accru dès surpoids ; majoré si tour de taille $\uparrow$ (H $>$ 94 cm, F $>$ 80 cm) \\ \hline
Besoins Énergétiques (BE) & $BE = MB \times NAP$ & MB : équation Harris-Benedict ou Mifflin-St Jeor ; NAP voir tableau ci-dessous \\ \hline
NAP (Facteur) & Sédentaire 1,3 \textbar\ Léger 1,5 \textbar\ Modéré 1,75 \textbar\ Intense 2,0 \textbar\ Très intense 2,4 & Adapter selon profession \& loisirs \\ \hline
Répartition Énergétique (ANC) & Glucides 50--55\% \textbar\ Lipides 30--35\% \textbar\ Protéines 10--15\% & Glucides complexes $\uparrow$, Sucres ajoutés $<$ 10\% (idéal $<$ 5\%), AGS $<$ 12\%, AGT $<$ 1\% \\ \hline
Fibres & $\ge$ 25--30 g/j & Fruits, légumes, légumineuses, céréales complètes \\ \hline
Sel (NaCl) & $<$ 5 g/j (OMS) $\approx$ 2 g Na/j & Attention sel caché (pain, conserves, condiments type Maggi\footnote{Cube bouillon très utilisé au Cameroun, très riche en NaCl et glutamate.}) \\ \hline
Eau & 1,5--2 L/j (boissons) + eau aliments & Augmenter forte chaleur / activité intense \\ \hline
\end{tabularx}
\end{center}

\courssub{Méthodes express}
\begin{methode}[Calculer et interpréter l'IMC]
\begin{enumerate}
    \item Mesurer masse $m$ (kg, balance calibrée) et taille $t$ (m, stadiomètre, pieds nus).
    \item Calculer $IMC = m / t^2$. Arrondir au dixième.
    \item Situer la valeur dans la classification OMS (tableau ci-dessus).
    \item \textbf{Nuancer} : Préciser les limites (musculature, personnes âgées, grossesse, ethnies). Associer le tour de taille (indicateur obésité androïde/viscérale).
\end{enumerate}
\end{methode}

\begin{methode}[Évaluer une ration alimentaire (Journée type)]
\begin{enumerate}
    \item Lister aliments \& quantités (pesées ou portions usuelles : cuillère, bol, morceau).
    \item Utiliser table de composition (ex: Ciqual, table africaine) pour extraire : Énergie (kcal), Protéines (g), Lipides (g), Glucides (g), Fibres (g), Sel (g), Fer, Ca, Vit C...
    \item Totaliser par nutriment et pour la journée.
    \item Calculer \% énergie : $\%G = \frac{G \times 4}{E_{tot}} \times 100$, $\%L = \frac{L \times 9}{E_{tot}} \times 100$, $\%P = \frac{P \times 4}{E_{tot}} \times 100$.
    \item Comparer aux ANC (Besoins Énergétiques calculés + Répartition + Fibres + Sel + Micronutriments).
    \item Identifier écarts (excès/déficits) et proposer corrections concrètes (substitutions, portions, fréquence).
\end{enumerate}
\end{methode}

\begin{methode}[Élaborer un programme de prévention combiné (Alimentation + AP)]
\begin{enumerate}
    \item \textbf{Diagnostic} : IMC, Tour de taille, Habitudes alimentaires (fréquence, diversité), Niveau AP (questionnaire type IPAQ ou GPAQ), Antécédents familiaux, Facteurs socio-environnementaux (accès marché, sécurité, temps, budget).
    \item \textbf{Objectifs SMART} : Spécifiques, Mesurables, Atteignables, Réalistes, Temporels (ex: -5\% masse grasse en 3 mois ; marcher 30 min/j 5j/sem ; 5 fruits/légumes/j).
    \item \textbf{Volet Alimentation} : Rééquilibrage (pas de régime restrictif) : $\uparrow$ densité nutritionnelle, $\downarrow$ densité énergétique (produits ultra-transformés, boissons sucrées, fritures), fractionnement (3 repas + 1 collation si besoin), hydratation eau, cuissons douces.
    \item \textbf{Volet Activité Physique} : Recommandations OMS 18--64 ans : $\ge$ 150--300 min/sem AP modérée \textbf{OU} 75--150 min/sem AP soutenue \textbf{ET} renforcement musculaire 2j/sem. Réduire temps assis (lever toutes les 30--60 min). Adapter goûts, capacités, environnement (marche rapide, danse, vélo, sport collectif, travaux ménagers intenses).
    \item \textbf{Suivi \& Évaluation} : Carnet de bord / Appli, pesée mensuelle (pas quotidienne), réévaluation IMC/tour taille à 3 mois, ajustement programme. Soutien motivationnel (entourage, groupe, professionnel de santé).
\end{enumerate}
\end{methode}
\end{bilan}

\begin{aretenir}[Checklist avant le Probatoire]
\begin{itemize}
    \item $\square$ Je sais définir l'équilibre alimentaire, le métabolisme de base, le NAP, l'IMC, la sédentarité, l'obésité, le DT2, les MCV.
    \item $\square$ Je connais la formule de l'IMC ($m/t^2$) et je sais le calculer avec les bonnes unités (kg, m).
    \item $\square$ Je sais classer un IMC selon l'OMS (maigreur, normal, surpoids, obésité I/II/III) et interpréter le risque sanitaire associé.
    \item $\square$ Je connais les limites de l'IMC (ne distingue pas muscle/graisse, ni répartition) et l'intérêt du tour de taille.
    \item $\square$ Je sais calculer les Besoins Énergétiques (BE = MB $\times$ NAP) et citer les facteurs du MB et les valeurs du NAP.
    \item $\square$ Je connais les repères ANC de répartition énergétique (G 50-55\%, L 30-35\%, P 10-15\%), fibres ($\ge$ 25-30 g), sel ($<$ 5 g), eau (1,5-2 L).
    \item $\square$ Je sais expliquer les mécanismes liant sédentarité/mauvaise alimentation $\rightarrow$ obésité $\rightarrow$ résistance à l'insuline $\rightarrow$ DT2 / dyslipidémie / HTA $\rightarrow$ MCV.
    \item $\square$ Je sais lister les recommandations OMS d'AP pour l'adulte (volume, intensité, renforcement, lutte contre sédentarité).
    \item $\square$ Je sais évaluer une ration alimentaire (méthode en 6 étapes) et identifier les écarts par rapport aux ANC.
    \item $\square$ Je sais structurer un programme de prévention personnalisé (Diagnostic $\rightarrow$ Objectifs SMART $\rightarrow$ Volet Alim $\rightarrow$ Volet AP $\rightarrow$ Suivi).
    \item $\square$ Je sais donner des exemples concrets camerounais (aliments locaux riches en fibres/fer, pièges sel caché/bouillons, AP intégrée au quotidien).
\end{itemize}
\end{aretenir}

\findecours

\end{document}
